\documentclass[
  paper=b5,
  fontsize=13pt,
  twoside=true,
  open=any,
  bookmarkpackage=false,
  chapterprefix=true,
  headings=normal
]{scrbook}

\PassOptionsToPackage{dvipsnames}{xcolor}
\usepackage[T1]{fontenc}
\usepackage[utf8]{inputenc}
\usepackage[english]{babel}
\usepackage{amsmath,amsthm,amsfonts,mathtools}
\usepackage{xcolor}
\usepackage{newpxtext}
\usepackage{newpxmath}
\usepackage[
  paperwidth=176mm,
  paperheight=250mm,
  left=7mm,
  right=7mm,
  top=9.1mm,
  bottom=9.1mm,
  includehead,
  headheight=16pt,
  headsep=4pt
]{geometry}
\usepackage{microtype}
\usepackage{array,booktabs,longtable,tabularx}
\usepackage{enumitem}
\usepackage{fancyhdr}
\usepackage{xurl}

\definecolor{RieszLink}{HTML}{174A7E}
\definecolor{RieszCite}{HTML}{2E6B4F}
\definecolor{RieszURL}{HTML}{8A3B3B}

\usepackage[
  unicode=true,
  colorlinks=true,
  hyperfootnotes=true,
  hypertexnames=false,
  linkcolor=RieszLink,
  citecolor=RieszCite,
  urlcolor=RieszURL,
  bookmarksopen=true,
  bookmarksnumbered=true,
  pdfstartview=FitH
]{hyperref}

\hypersetup{
  pdftitle={Les systèmes d'équations linéaires à une infinité d'inconnues / Linear Systems of Equations with Infinitely Many Unknowns},
  pdfauthor={Frédéric Riesz},
  pdfsubject={Faithful English translation of the 1913 French edition}
}
\urlstyle{same}

\linespread{1.1}
\allowdisplaybreaks[4]
\setcounter{MaxMatrixCols}{30}
\setcounter{secnumdepth}{0}
\interfootnotelinepenalty=0
\setlength{\emergencystretch}{2em}
\setlength{\parindent}{1.25em}
\setlength{\parskip}{0pt}
\setlist{nosep}

\renewcommand*{\raggedchapter}{\centering}
\renewcommand*{\thechapter}{\Roman{chapter}}
\addto\captionsenglish{\renewcommand*{\contentsname}{Table of Contents}}
\RedeclareSectionCommand[tocnumwidth=3.2em]{chapter}
\setkomafont{disposition}{\normalfont\rmfamily\bfseries}
\setkomafont{chapter}{\normalfont\rmfamily\bfseries\Large}
\setkomafont{chapterprefix}{\normalfont\rmfamily\bfseries\Large}
\setkomafont{section}{\normalfont\rmfamily\bfseries\large}
\setkomafont{subsection}{\normalfont\rmfamily\bfseries\normalsize}

\newtheorem*{theorem}{Theorem}
\newtheorem*{vonkochtheorem}{von Koch's Theorem}
\newtheorem*{landautheorem}{Landau's Theorem}
\newtheorem*{selectionprinciple}{Selection Principle}
\newtheorem*{lemma}{Lemma}
\newtheorem*{proposition}{Proposition}
\newtheorem*{corollary}{Corollary}
\newtheorem*{criterion}{Criterion}
\theoremstyle{definition}
\newtheorem*{definition}{Definition}
\newtheorem*{remark}{Remark}

\newcommand{\numberedparagraph}[1]{%
  \par\medskip\phantomsection\label{sec:#1}%
  \noindent\textbf{#1.}\enspace
}
\newcommand{\literature}[2]{\href{#1}{#2}}
\newcounter{translatornote}
\makeatletter
\newcommand{\translatornotealpha}{%
  \ifnum\value{translatornote}>26
    a\@alph{\numexpr\value{translatornote}-26\relax}%
  \else
    \@alph{\value{translatornote}}%
  \fi
}
\makeatother
\newcommand{\translatornote}[1]{%
  \stepcounter{translatornote}%
  \begingroup
    \edef\savedfootnotenumber{\number\value{footnote}}%
    \renewcommand{\thefootnote}{\translatornotealpha}%
    \footnote{\emph{Translator's note:} #1}%
    \setcounter{footnote}{\savedfootnotenumber}%
  \endgroup
}

\fancyhf{}
\fancyhead[LE]{\thepage}
\fancyhead[RO]{\thepage}
\renewcommand{\headrulewidth}{0pt}
\renewcommand{\footrulewidth}{0pt}
\pagestyle{fancy}
\fancypagestyle{plain}{%
  \fancyhf{}
  \fancyhead[LE]{\thepage}
  \fancyhead[RO]{\thepage}
  \renewcommand{\headrulewidth}{0pt}
  \renewcommand{\footrulewidth}{0pt}
}

\begin{document}
\pagenumbering{arabic}
\setcounter{page}{1}

\begin{titlepage}
  \thispagestyle{empty}
  \centering
  {\small
    COLLECTION OF MONOGRAPHS ON THE THEORY OF FUNCTIONS\\
    PUBLISHED UNDER THE DIRECTION OF M.~ÉMILE BOREL
  \par}
  \vspace{2.2cm}
  {\bfseries\LARGE
    LES SYSTÈMES\\[0.35em]
    D'ÉQUATIONS LINÉAIRES\\[0.35em]
    À UNE INFINITÉ D'INCONNUES
  \par}
  \vspace{1.1cm}
  {\bfseries\Large
    Linear Systems of Equations\\[0.25em]
    with Infinitely Many Unknowns
  \par}
  \vspace{1.6cm}
  {\small BY\par}
  \vspace{0.2cm}
  {\large FRÉDÉRIC RIESZ\par}
  \vspace{0.35cm}
  {\small
    PROFESSOR AT THE ROYAL HUNGARIAN UNIVERSITY OF KOLOZSVÁR
  \par}
  \vfill
  {\small
    PARIS\\
    GAUTHIER--VILLARS, PRINTER--PUBLISHER\\
    TO THE BUREAU OF LONGITUDES AND THE ÉCOLE POLYTECHNIQUE\\
    Quai des Grands-Augustins, 55\\[0.45em]
    1913
  \par}
\end{titlepage}

\clearpage
\thispagestyle{fancy}
\vspace*{\fill}
\begin{center}
  All rights of translation, reproduction, and adaptation reserved in all countries.
\end{center}
\vspace*{\fill}

\chapter*{Preface}
\addcontentsline{toc}{chapter}{Preface}

In this volume I should like to give a rapid account of the fundamental ideas,
methods, and principal results of a theory which is due almost exclusively to
contemporary mathematicians. Personally, I have contributed very little to this
theory; and if I have undertaken the present work, it is because I was encouraged
to do so by my investigations into several neighboring subjects, among them
orthogonal systems of functions, integral equations, and functional operations.
Continued by several authors, these investigations have proved more or less
fruitful for applications of the present theory; and, on the other hand, they have
enabled me to present some parts of it from new points of view. Thus, for example,
I was able to connect the study of the spectrum of quadratic forms in infinitely
many variables with that of linear functional operations.

Our subject does not belong to the Theory of Functions in the strict sense. It
should rather be regarded as marking a first stage in the still nascent Theory of
Functions of Infinitely Many Variables, which may perhaps soon furnish the most
powerful methods in all of Analysis. In any event, I believe that the subject and
arrangement of the present volume accord well with the general plan of this
Collection, in which M.~Borel kindly offered to include it.

Finally, I wish to thank M.~Fréchet and M.~Marcel Riesz for the valuable advice
they gave me during the correction of the proofs.

\begin{flushright}
  Kolozsvár, June 12, 1913.\\[0.6em]
  \textsc{Frédéric Riesz}
\end{flushright}

\clearpage

% Begin inlined .parts/pages_001_060.tex
% [Source printed page 1]
\chapter{The Beginnings of the Theory}\label{chap:1}

\section{The Method of Undetermined Coefficients}

\numberedparagraph{1}
The theory of equations with infinitely many unknowns is not of long standing. Indeed, scarcely any such theory existed before 1886, the year in which Poincar\'e demonstrated the legitimacy of the method of determinants of infinite order. General attention, moreover, was not really drawn to it until two memoirs by Hilbert appeared in 1906 as part of a series of memoirs on Fredholm's equation; from that time onward the theory developed rapidly, and the importance and unexpected character of the results obtained caused almost everything already known before then to be forgotten.

Yet theories have their beginnings: vague allusions, unfinished attempts, particular problems; and even when these beginnings matter little in the present state of science, it would be wrong to pass them over in silence.

\numberedparagraph{2}
The study of systems of equations with infinitely many unknowns arose from the method of undetermined coefficients. This method, used since the seventeenth century for integrating differential equations by series and for other problems, consists in principle in replacing the function sought by a series with unknown coefficients. The data of the problem will then make it possible to establish among these unknown coefficients, infinite in number, infinitely many relations; this is the system of equations with infinitely many unknowns.

% [Source printed page 2]
For most problems treated in this way, however, the infinitude of the number of unknowns involved no difficulty at all and, it seems, no one even realized that a new idea was involved. This was because only recurrent systems arose, in which each equation by itself contained only finitely many unknowns, and one had to solve only finite systems, though infinitely many times.

\section{Fourier and the Principle of Sections}

\numberedparagraph{3}
While studying a special case of what is now called the Dirichlet problem, Fourier encountered, in his \emph{Analytical Theory of Heat}, a system that was no longer recurrent. He proposed to determine a solution $v(x,y)$ of the partial differential equation
\begin{equation}
  v_{xx}+v_{yy}=0,
  \tag{1}\label{eq:3:1}
\end{equation}
valid in the domain
\[
  x>0,\qquad -\frac{\pi}{2}<y<\frac{\pi}{2},
\]
equal to $1$ for $x=0$, vanishing for $y=\pm\pi/2$ and for $x=\infty$.\footnote{Art.~166 and following.} Here is how he proceeds. First he sets
\[
  v(x,y)=F(x)f(y)
\]
and seeks to satisfy equation~\eqref{eq:3:1} without regard to the boundary data. He thus finds the particular solution
\begin{equation}
  v(x,y)=e^{-mx}\cos my,
  \tag{2}\label{eq:3:2}
\end{equation}
where $m$ denotes an arbitrary real number. On choosing $m$ positive, integral, and odd, solution~\eqref{eq:3:2} satisfies the boundary conditions except the first. Fourier therefore sets
\[
  v(x,y)=\sum_{m=1}^{\infty}a_m e^{-(2m-1)x}\cos(2m-1)y;
\]
it remains only to determine the coefficients $a_m$ so that, for $x=0$, the sum $v(x,y)$ of the series reduces to $1$. This amounts, he says, to requiring
% [Source printed page 3]
\begin{equation}
  1=\sum_{m=1}^{\infty}a_m\cos(2m-1)y.
  \tag{3}\label{eq:3:3}
\end{equation}
To eliminate the variable $y$, he differentiates the equation term by term infinitely many times and puts $y=0$; this gives
\begin{equation}
\left\{
\begin{aligned}
  1&=\sum a_m,\\
  0&=\sum(2m-1)^2a_m,\\
  0&=\sum(2m-1)^4a_m,\\
  &\hspace{1em}\cdots
\end{aligned}
\right.
\tag{4}\label{eq:3:4}
\end{equation}
a system of infinitely many linear equations in the infinitely many unknowns $a_m$.

\numberedparagraph{4}
To solve it, Fourier takes the first $k$ equations, retains in them only the first $k$ unknowns, and neglects all terms depending on the remaining unknowns. Denoting the unknowns of this reduced system by
\[
  a_1^{(k)},a_2^{(k)},\ldots,a_k^{(k)},
\]
one obtains, by an easy calculation,
\[
  a_1^{(k)}=\frac{9\cdot25\cdots(2k-1)^2}{8\cdot24\cdots(4k^2-4k)},
\]
which, by a well-known formula of Wallis, gives
\[
  \lim_{k\to\infty}a_1^{(k)}=\frac4\pi.
\]
Moreover, for $m<k$, the reduced system gives
\[
  \frac{a_{m+1}^{(k)}}{a_m^{(k)}}
  =\frac{2m-1}{2m+1}\frac{m-k}{m+k};
\]
hence
\[
  \lim_{k\to\infty}\frac{a_{m+1}^{(k)}}{a_m^{(k)}}
  =-\frac{2m-1}{2m+1},
\]
and, by induction,
\[
  \lim_{k\to\infty}a_m^{(k)}=(-1)^{m-1}\frac4{\pi(2m-1)}.\footnote{We have just calculated this limiting value by a route different from Fourier's, whose calculation is somewhat more laborious. The essential principles of the argument are not thereby altered.}
\]

% [Source printed page 4]
Finally, putting
\[
  a_m=\lim_{k\to\infty}a_m^{(k)},
\]
we obtain
\begin{equation}
  v(x,y)=\frac4\pi\sum_{m=1}^{\infty}(-1)^{m-1}
  \frac{e^{-(2m-1)x}\cos(2m-1)y}{2m-1}.
  \tag{5}\label{eq:4:5}
\end{equation}

\numberedparagraph{5}
This result of Fourier is undoubtedly correct. Indeed, comparing series~\eqref{eq:4:5} and its derivatives of every order respectively with the series
\[
  \sum\frac{e^{-(2m-1)x_0}}{2m-1},\qquad
  \sum e^{-(2m-1)x_0},\qquad
  \sum(2m-1)e^{-(2m-1)x_0},\quad\ldots,
\]
one concludes that they converge uniformly for $x\geq x_0>0$, that they tend uniformly to zero as $x$ increases indefinitely, that $v(x,y)$ vanishes for $x>0$, $y=\pm\pi/2$, and finally, since termwise differentiation is legitimate, that $v(x,y)$ satisfies equation~\eqref{eq:3:1}. On the portion $x=0$, $-\pi/2<y<\pi/2$ of the boundary of the domain, series~\eqref{eq:4:5} converges to the limit $1$, uniformly on every segment $(-\pi/2+h,\pi/2-h)$; moreover $v(x,y)$ tends to this value when, from within the domain, one approaches such a segment indefinitely. At the points $(0,\pi/2)$ and $(0,-\pi/2)$, where the prescribed values undergo abrupt changes, $v(x,y)$ remains indeterminate, the limiting bounds of indeterminacy being $0$ and $1$.

All these facts are connected with the following observation. Introducing the variables
\[
  r=e^{-x},\qquad \theta=y,\qquad z=re^{i\theta},
\]
which makes the half-disk
\[
  r<1,\qquad -\frac\pi2<\theta<\frac\pi2
\]
correspond to the domain considered, the function $v(x,y)$ becomes the real part of one of the branches of the analytic function
\[
  f(z)=\frac1{2i}\log\frac{i-z}{i+z}.
\]

% [Source printed page 5]
This real part has a very simple geometric meaning: it is equal to half the angle whose sides join the point $z$ to the points $i$ and $-i$. All the stated consequences follow from this.

\numberedparagraph{6}
Fourier's result is therefore correct. But serious objections may be raised against his reasoning. Fourier was reasoning about an unknown series; he could not suspect the dangers that threatened him, and in that period confidence still prevailed. Since then, however, greater caution has been brought to every question involving infinity. Had Fourier shared our mistrust, would he have dared to take the path he followed?

Let us merely indicate the most perilous stage of his argument. He reduces his problem to that of determining a series of the form~\eqref{eq:3:3} which represents, throughout the interval $(-\pi/2,\pi/2)$, the function constantly equal to $1$. To calculate the $a_m$, he differentiates equation~\eqref{eq:3:3} infinitely many times and puts $y=0$, thereby obtaining system~\eqref{eq:3:4}. But from the second equation onward, the series occurring there diverge after the $a_m$ have been replaced by the values $(-1)^{m-1}4/[\pi(2m-1)]$ that were found. Thus at first sight the values obtained for the $a_m$ do not appear to satisfy system~\eqref{eq:3:4}.

Fortunately, we now know how to interpret broad classes of divergent series by assigning them sums that retain certain principal properties of sums of convergent series. From this point of view, to justify Fourier's procedure one need only replace system~\eqref{eq:3:4} by
\begin{equation}
\left\{
\begin{aligned}
  1&=\lim_{r\to1-0}\sum a_m r^{2m-1},\\
  0&=\lim_{r\to1-0}\sum(2m-1)^2a_m r^{2m-1},\\
  0&=\lim_{r\to1-0}\sum(2m-1)^4a_m r^{2m-1},\\
  &\hspace{1em}\cdots
\end{aligned}
\right.
\tag{6}\label{eq:6:6}
\end{equation}
This interpretation of system~\eqref{eq:3:4} is entirely in accord with the problem being treated.\footnote{Other summation methods might also have been used; for example, the series in question
\[
  \sum\frac{(-1)^{m-1}}{(2m-1)^{2k+1}}
\]
is summable by arithmetic means of order $2k$, and this procedure likewise gives the sum $0$ for $k>0$.} Indeed, the values found for the $a_m$
% [Source printed page 6]
satisfy system~\eqref{eq:6:6}, and this fact is connected with the holomorphy at $z=1$ of the function $f(z)$ defined above.

\numberedparagraph{7}
Fourier also applies his method to a more general problem: to expand in a trigonometric series an odd function given by its Taylor series
\[
  f(x)=f'(0)x+\frac{f'''(0)}{3!}x^3+\cdots.\footnote{Art.~207 and following.}
\]
Since the function is assumed odd, one must evidently find an expansion of the form
\begin{equation}
  f(x)=\sum_{m=1}^{\infty}b_m\sin mx.
  \tag{7}\label{eq:7:7}
\end{equation}
To determine the $b_m$, Fourier uses a procedure like the one just described: he differentiates both sides of~\eqref{eq:7:7} infinitely many times and sets $x=0$; or rather, what amounts to the same thing, he expands both sides in power series and compares coefficients of like powers. This procedure gives him an infinite system of equations in the unknowns $b_m$; he takes the first $k$, retains only the first $k$ unknowns, and solves the reduced system; he then passes to the limit, if one may call the extremely bold calculation he performs a passage to the limit. That calculation gives
\[
  b_m=(-1)^{m+1}\frac2{m\pi}
  \left[f(\pi)-\frac1{m^2}f''(\pi)+\frac1{m^4}f^{\mathrm{IV}}(\pi)-\cdots\right].
\]
To deduce the final result, Fourier observes that the series
\[
  s(x)=f(x)-\frac1{m^2}f''(x)+\frac1{m^4}f^{\mathrm{IV}}(x)-\cdots
\]
% [Source printed page 7]
satisfies the differential equation
\[
  \frac1{m^2}s''(x)+s(x)=f(x).
\]
The general integral of this equation is
\begin{align*}
  s(x)={}&C_1\cos mx+C_2\sin mx
  +m\sin mx\int_0^x f(t)\cos mt\,dt\\
  &{}-m\cos mx\int_0^x f(t)\sin mt\,dt.
\end{align*}
Since $s(x)$, like $f(x)$, is odd, one has
\[
  C_1=s(0)=0;
\]
and, on putting $x=\pi$, one obtains the well-known formula
\begin{equation}
  b_m=\frac2\pi\int_0^\pi f(x)\sin mx\,dx.
  \tag{8}\label{eq:7:8}
\end{equation}

\numberedparagraph{8}
Fourier's reasoning is not correct, and could be made so only under very restrictive conditions on the given function $f(x)$. On the other hand, since the theory of trigonometric series has subsequently been based on firmer foundations, we now know that the validity of formula~\eqref{eq:7:8} requires only very broad conditions. For the theory of trigonometric series, Fourier's argument is merely an interesting curiosity of purely historical value. For us, however, it contains something very valuable: it implies an exceedingly fruitful principle from the standpoint of our equations. Here is that principle, which must of course still be made much more precise: to solve an infinite system of equations with infinitely many unknowns, restrict the system to its first $k$ equations and neglect in them all unknowns except the first $k$. As $k$ tends to infinity, the solutions of these systems tend to the solution of the proposed system.

This principle, whose legitimacy can be justified under broad hypotheses, will prove very fruitful for the theory before us. We shall call it the \emph{principle of sections}.

\numberedparagraph{9}
We have just examined the principle of sections from the point of
% [Source printed page 8]
view of Fourier's use of it. Let us now consider the same principle starting from a classical theorem in the Theory of Functions. According to Weierstrass, given an infinite sequence of quantities $a_i$ such that $|a_i|\to\infty$, there exists an entire function
\[
  \mathfrak F(z)=C_0+C_1z+C_2z^2+\cdots
\]
having the $a_i$ as zeros and no others. The coefficients $C_0,C_1,C_2,\ldots$ must satisfy the equations
\[
  C_0+C_1a_i+C_2a_i^2+\cdots=0\qquad(i=1,2,\ldots).
\]
For definiteness, suppose the $a_i$ are distinct and nonzero. The principle of sections, applied to these equations, then leads to the successive products
\[
  \prod_{i=1}^{n}\left(1-\frac z{a_i}\right)
  =C_0^{(n)}+C_1^{(n)}z+\cdots+C_n^{(n)}z^n,
  \qquad C_0^{(n)}=1.
\]
Indeed, each of these products corresponds to one of the reduced systems, so that this system is satisfied by $C_0^{(n)},\ldots,C_n^{(n)}$. But this procedure is known to converge only if the $|a_k|$ increase sufficiently rapidly. Thus our equations can always be solved, but the principle of sections applies only under very restrictive conditions.

Weierstrass's method consists in adjoining exponential factors to the products above so as to accelerate convergence. One may hope that this method will some day be generalized so as to perfect the principle of sections. Up to the present, this has not been attempted.

\section{F\"urstenau and K\"otteritzsch}

\numberedparagraph{10}
Although Fourier's \emph{Analytical Theory} has always remained the source of numerous investigations, the portions just recalled have fallen almost entirely into oblivion. For more than half a century, the authors who dealt with our subject were led to it independently of Fourier and without referring to one another. Only in a memoir by G.~Piola
% [Source printed page 9]
do we find Fourier's investigations cited; that author applied the same principle of sections to a particular problem.\footnote{Piola, \emph{Sulla teoria delle funzioni discontinue}, \emph{Memorie della Societ\`a Italiana delle Scienze}, vol.~XX, pp.~573--639. We cite this memoir from the reference supplied by G.~Loria, \emph{Supplemento ai Rendiconti}, Palermo, 1907, p.~34.}

The next two authors in historical order, E.~F\"urstenau and Th.~K\"otteritzsch, are of no importance for the development of the theory. Their works nevertheless contain methods in embryonic form which, rediscovered and made rigorous in recent years, have proved very fruitful for the theory under consideration.

\numberedparagraph{11}
F\"urstenau considers the following problem:

\emph{Given the equation}
\begin{equation}
  c_0+c_1x+\cdots+c_mx^m=0\qquad(c_0\ne0),
  \tag{9}\label{eq:11:9}
\end{equation}
\emph{calculate the root having the least absolute value}.\footnote{F\"urstenau, \emph{Darstellung der reellen Wurzeln algebraischer Gleichungen durch Determinanten der Coeffizienten}, program dissertation, Marburg, 1860; \emph{Neue Methode zur Darstellung und Berechnung der imagin\"aren Wurzeln algebraischer Gleichungen durch Determinanten der Coeffizienten}, program dissertation, Marburg, 1867.} For this purpose he successively multiplies the equation by higher and higher powers of $x$, and replaces $x^n$ by $x_n$; this procedure leads him to the infinite system
\[
\left\{
\begin{aligned}
  -c_0&=c_1x_1+c_2x_2+\cdots+c_mx_m,\\
  0&=c_0x_1+c_1x_2+\cdots+c_{m-1}x_m+c_mx_{m+1},\\
  0&=c_0x_2+\cdots+c_{m-2}x_m+c_{m-1}x_{m+1}+c_mx_{m+2},\\
  &\hspace{1em}\cdots
\end{aligned}
\right.
\]
To derive an expression for $x_1$, he applies the principle of sections. He obtains
\begin{equation}
  x_1=-c_0\lim_{n\to\infty}\frac{\Delta_{n-1}}{\Delta_n},
  \tag{10}\label{eq:11:10}
\end{equation}
where $\Delta_n$ is the determinant formed from the coefficients of $x_1,\ldots,x_n$ in the first $n$ equations.

When there is only one root of least absolute value and this root is simple, it is given by expression~\eqref{eq:11:10}. F\"urstenau's memoir proves this result with complete rigor, but by a rather laborious calculation. According to an observation of H.~Naegelsbach, moreover, everything reduces to the fact that the function
% [Source printed page 10]
\[
  R(x)=\frac1{c_0+c_1x+\cdots+c_mx^m}
\]
has the power-series expansion
\[
  R(x)=\frac1{c_0}-\frac{\Delta_1}{c_0^2}x
       +\frac{\Delta_2}{c_0^3}x^2
       -\frac{\Delta_3}{c_0^4}x^3+\cdots.\footnote{Naegelsbach, \emph{Studien zu F\"urstenaus neuer Methode der Darstellung und Berechnung der Wurzeln algebraischer Gleichungen durch Determinanten der Coeffizienten}, \emph{Archiv f\"ur Mathematik und Physik}, vol.~LIX, 1876, pp.~147--192.}
\]
Indeed, let $\alpha$ be the root in question and let $A$ be the residue of $R(x)$ at $\alpha$. Since the poles of the rational fraction
\[
  r(x)=R(x)-\frac{A}{1-x/\alpha}
\]
are farther from the origin than the point $\alpha$, the series
\[
  r(x)=\frac1{c_0}-A+
  \left(-\frac{\Delta_1}{c_0^2}-\frac A\alpha\right)x+
  \left(\frac{\Delta_2}{c_0^3}-\frac A{\alpha^2}\right)x^2+\cdots
\]
converges for $x=\alpha$; in particular its terms tend to zero, and therefore
\[
  \frac{\Delta_n(-\alpha)^n}{c_0^{n+1}}\longrightarrow A.
\]
It follows that
\[
  -c_0\frac{\Delta_{n-1}}{\Delta_n}
  =\alpha\,
   \frac{\Delta_{n-1}(-\alpha)^{n-1}/c_0^n}
        {\Delta_n(-\alpha)^n/c_0^{n+1}}
  \longrightarrow\alpha.
\]

F\"urstenau's method extends immediately to transcendental equations. Observe also that, if for brevity we write
\[
  R(x)=a_0+a_1x+a_2x^2+\cdots,
\]
% [Source printed page 11]
our result becomes
\[
  \alpha=\lim_{n\to\infty}\frac{a_{n-1}}{a_n}.
\]
This elementary and well-known result concerns the power-series expansion of any function having a simple pole $\alpha$, all its other singular points being farther from the origin.\footnote{K\"onig, ``Ueber eine Eigenschaft der Potenzreihen,'' \emph{Mathematische Annalen}, vol.~XXIII (1884), pp.~447--449.} It was generalized by Hadamard, who sought to determine several poles at once.\footnote{Hadamard, ``Sur la recherche des discontinuit\'es polaires,'' \emph{Comptes rendus}, April~8, 1889. See also Hadamard, \emph{La s\'erie de Taylor et son prolongement analytique}, pp.~38--43. Compare also the little-known memoir Worpitzky, \emph{Beitr\"age zur Functionentheorie}, program dissertation, Berlin, 1870, whose existence I learned only during the final proof corrections.} F\"urstenau himself extended his method to multiple roots and to the simultaneous calculation of several roots. It would be of some interest to compare his calculations with Hadamard's.

\numberedparagraph{12}
Up to this point only particular systems had been considered. K\"otteritzsch begins at once with the general system
\begin{equation}
  \sum_{k=1}^{\infty}a_{ik}x_k=\alpha_i
  \qquad(i=1,2,\ldots),
  \tag{11}\label{eq:12:11}
\end{equation}
subject only to rather weak restrictions.\footnote{K\"otteritzsch, ``Ueber die Aufl\"osung eines Systems von unendlich vielen linearen Gleichungen,'' \emph{Zeitschrift f\"ur Mathematik und Physik}, vol.~XV (1870), pp.~1--15, 229--268.}

To solve it, he first reduces it to the more special form
\begin{equation}
\left\{
\begin{aligned}
 b_{11}x_1+b_{12}x_2+b_{13}x_3+\cdots&=\beta_1,\\
 b_{22}x_2+b_{23}x_3+\cdots&=\beta_2,\\
 b_{33}x_3+\cdots&=\beta_3,\\
 &\hspace{1em}\cdots
\end{aligned}
\right.
\tag{12}\label{eq:12:12}
\end{equation}
This reduction is effected by an entirely elementary calculation. Suppose, in fact, that the diagonal minors
\[
  |a_{ik}|_1=a_{11},\qquad
  |a_{ik}|_2=a_{11}a_{22}-a_{12}a_{21},\qquad\ldots
\]
do not vanish; then, to obtain the corresponding system~\eqref{eq:12:12},
% [Source printed page 12]
one need only put $b_{1k}=a_{1k}$, $\beta_1=\alpha_1$, and, for $n>1$,
\[
 b_{nk}=
 \begin{vmatrix}
 a_{1,1}&a_{1,2}&\cdots&a_{1,n-1}&a_{1,k}\\
 \vdots&\vdots&&\vdots&\vdots\\
 a_{n,1}&a_{n,2}&\cdots&a_{n,n-1}&a_{n,k}
 \end{vmatrix},
 \qquad
 \beta_n=
 \begin{vmatrix}
 a_{1,1}&a_{1,2}&\cdots&a_{1,n-1}&\alpha_1\\
 \vdots&\vdots&&\vdots&\vdots\\
 a_{n,1}&a_{n,2}&\cdots&a_{n,n-1}&\alpha_n
 \end{vmatrix}.
\]
It remains to solve system~\eqref{eq:12:12}. Suppose that the unknown $x_n$ is to be evaluated. The first $n-1$ equations do not enter the calculation; from the remaining ones, $x_{n+1},x_{n+2},\ldots$ are successively eliminated. This procedure yields an expansion of the form
\[
  x_n=B_{n,n}\beta_n+B_{n,n+1}\beta_{n+1}+\cdots,
\]
where the $B$ are expressions formed from the coefficients $b_{ik}$ and readily calculated by determinants.

It is easy to see that K\"otteritzsch's method is at bottom the same as Fourier's; it differs only in containing a formal computational procedure that may be convenient in certain cases. From the standpoint of infinity, however, the two methods do not differ. On this point K\"otteritzsch teaches us nothing new; he merely says that for large indices the terms must be very small, so that they may be neglected.

Finally, K\"otteritzsch applies his method to several special cases and extends to them formulas known for finite systems. Here too, however, almost everything rests on analogy.\footnote{Mention should be made here of a paper by von Koch which was to appear in the \emph{Proceedings of the International Congress at Cambridge}, 1912. Von Koch studies systems of type~\eqref{eq:12:12} under very broad conditions. I confine myself to stating one special result, connected among other things with F\"urstenau's problem. Let $f(z)=1+c_1z+c_2z^2+\cdots$ be an entire function, and let $\alpha$ be the zero nearest the origin. Further, let $x_1,x_2,\ldots$ be a solution of
\[
\begin{aligned}
0&=x_0+c_1x_1+c_2x_2+c_3x_3+\cdots,\\
0&=x_1+c_1x_2+c_2x_3+c_3x_4+\cdots,\\
&\hspace{1em}\cdots.
\end{aligned}
\]
Then $\limsup_{n\to\infty}|x_n|^{1/n}\geq|\alpha|$. See also Borel's thesis, discussed below, and St\"ackel, ``Periodische Funktionen und Systeme von unendlich vielen linearen Gleichungen,'' \emph{Festschrift H.~Weber}, 1912, pp.~396--409.}

\section{Appell's Note and Poincar\'e's Criticism}

\numberedparagraph{13}
Poincar\'e deserves the credit for making the first critical investigations of our subject. His attention was drawn to it by the works of Hill (1877) and Appell (1885). We shall return later to Hill's method.

% [Source printed page 13]
Appell proposes to expand the elliptic function $\theta_1/\theta$ in a trigonometric series by the method of undetermined coefficients.\footnote{Appell, ``Sur une m\'ethode \'el\'ementaire pour obtenir les d\'eveloppements en s\'eries trigonom\'etriques des fonctions elliptiques,'' \emph{Bulletin de la Soci\'et\'e math\'ematique de France}, vol.~XIII (1885), pp.~13--18.} In Jacobi's notation,
\[
 \theta\!\left(\frac{Kx}{\pi i}\right)
   =\sum_{n=-\infty}^{\infty}(-1)^n q^{n^2}e^{nx},\qquad
 \theta_1\!\left(\frac{Kx}{\pi i}\right)
   =\sum_{n=-\infty}^{\infty}q^{n^2}e^{nx},\qquad
 q=e^{-\pi K'/K};
\]
the periods $2K$ and $2iK'$ are arranged so that the real part of $K'/K$ is positive and hence $|q|<1$. One seeks the coefficients $\Lambda_n$ in the expansion
\[
  \frac{\theta_1(Kx/\pi i)}{\theta(Kx/\pi i)}
  =\sum_{n=-\infty}^{\infty}\Lambda_ne^{nx}.
\]
After clearing the denominator, multiplying out the right-hand side, and equating the coefficients of $e^{nx}$ for every $n$, one obtains
\[
 q^{n^2}=\sum_{\mu=-\infty}^{\infty}
 (-1)^{n-\mu}q^{(n-\mu)^2}\Lambda_\mu
 \qquad(n=0,1,2,\ldots),
\]
or, after simplification,
\[
 (-1)^n=\sum_{\mu=-\infty}^{\infty}
 (-1)^\mu q^{\mu^2-2\mu n}\Lambda_\mu
 \qquad(n=0,1,2,\ldots).
\]

% [Source printed page 14]
Since the quotient $\theta_1/\theta$ is even, $\Lambda_{-\mu}=\Lambda_\mu$; consequently the system reduces to\translatornote{The French text calls both $\theta_1$ and $\theta$ odd.  The displayed bilateral series show that each is even; only the evenness of their quotient is used here.}
\[
 (-1)^n=\Lambda_0+
 \sum_{\mu=1}^{\infty}(-1)^\mu q^{\mu^2}
       \bigl(q^{2n\mu}+q^{-2n\mu}\bigr)\Lambda_\mu
 \qquad(n=0,1,2,\ldots).
\]
Putting
\[
  \omega=\frac{2\pi K'i}{K},
\]
the system becomes
\[
 (-1)^n=\Lambda_0+2\sum_{\mu=1}^{\infty}
 (-1)^\mu q^{\mu^2}\Lambda_\mu\cos(n\mu\omega)
 \qquad(n=0,1,2,\ldots).
\]
To solve this system, Appell uses the principle of sections: he takes the first $m+1$ equations, retains in them only the first $m+1$ unknowns, suppresses every term depending on the others, and finally lets $m$ increase indefinitely. If $\Lambda_0^{(m)},\Lambda_1^{(m)},\ldots,\Lambda_m^{(m)}$ is the solution of the reduced system, then
\[
 \Lambda_0^{(m)}=\frac{\Delta(\pi,\omega,2\omega,\ldots,m\omega)}
                       {\Delta(0,\omega,2\omega,\ldots,m\omega)},
\]
and, for $\mu>0$,
\[
 (-1)^\mu q^{\mu^2}\Lambda_\mu^{(m)}
 =\frac{\Delta[0,\omega,\ldots,(\mu-1)\omega,\pi,
        (\mu+1)\omega,\ldots,m\omega]}
       {\Delta(0,\omega,2\omega,\ldots,m\omega)},
\]
where
\[
 \Delta(a,b,\ldots,l)=
 \begin{vmatrix}
 1&1&\cdots&1\\
 \cos a&\cos b&\cdots&\cos l\\
 \vdots&\vdots&&\vdots\\
 \cos ma&\cos mb&\cdots&\cos ml
 \end{vmatrix}
 =C\prod(\cos a-\cos b).
\]
Here $C$ is a numerical constant, and the product extends over all pairwise differences among $\cos a,\cos b,\ldots,\cos l$. Returning from $\omega$ to $q$, one obtains
\[
 \Lambda_0^{(m)}=\prod_{\nu=1}^{m}
       \frac{(1+q^{2\nu})^2}{(1-q^{2\nu})^2},
\]
\[
 \Lambda_\mu^{(m)}=\frac{2q^\mu}{1+q^{2\mu}}
 \frac{\displaystyle\prod_{\nu=1}^{m}(1+q^{2\nu})^2}
      {\displaystyle
       \prod_{\nu=1}^{m-\mu}(1-q^{2\nu})
       \prod_{\nu=1}^{m+\mu}(1-q^{2\nu})}.
\]

% [Source printed page 15]
Finally, on letting the index $m$ increase indefinitely, one obtains
\[
 \Lambda_0=\prod_{\nu=1}^{\infty}
 \frac{(1+q^{2\nu})^2}{(1-q^{2\nu})^2}
 =\frac{\pi}{2qK}\sqrt{\frac1{k'}},
 \qquad
 \Lambda_\mu=\frac{2q^\mu}{1+q^{2\mu}}\Lambda_0.
\]
Thus one recovers the known expansion
\[
 \frac{\theta_1(z)}{\theta(z)}
 =\frac{\pi}{2qK}\sqrt{\frac1{k'}}
 \left(1+4\sum_{\mu=1}^{\infty}
 \frac{q^\mu}{1+q^{2\mu}}\cos\frac{\mu\pi z}{K}\right).
\]

\numberedparagraph{14}
In the same \emph{Bulletin} that contained Appell's note just discussed, immediately following it, appeared Poincar\'e's first critical remarks.\footnote{Poincar\'e, ``Remarques sur l'emploi de la m\'ethode pr\'ec\'edente,'' \emph{Bulletin de la Soci\'et\'e math\'ematique de France}, vol.~XIII (1885), pp.~19--27.} He asks in what cases the method that succeeded for Appell may legitimately be used; that is, he examines, in certain cases, the legitimacy of the principle of sections.

He first considers the system
\begin{equation}
  \sum_{k=1}^{\infty}a_k^i x_k=0
  \qquad(i=0,1,2,\ldots),
  \tag{13}\label{eq:14:13}
\end{equation}
assuming $|a_{k+1}|>|a_k|$ and $|a_k|\to\infty$ as $k\to\infty$, and seeks to satisfy it in such a way that the series on the left converge.

His method rests entirely on the elements of the Theory of Functions. Form the entire function $F(z)$ having the numbers $a_k$ as simple zeros and having no other zeros. Since Weierstrass we have known that, under the hypothesis $|a_k|\to\infty$, such functions exist and may be calculated by infinite products. For definiteness, suppose that $\sum1/|a_k|$ converges; then put
\[
 F(z)=\prod_{k=1}^{\infty}\left(1-\frac z{a_k}\right).
\]
Let $C_1,C_2,\ldots,C_k,\ldots$ be infinitely many circles centered at the origin, such that $C_k$ separates the points $z=a_k$
% [Source printed page 16]
and $z=a_{k+1}$. Suppose that
\begin{equation}
  \int_{C_k}\frac{z^i\,dz}{F(z)}\longrightarrow0
  \qquad(k\to\infty),
  \tag{14}\label{eq:14:14}
\end{equation}
for every $i$. By Cauchy's theorem the preceding relation may be written
\[
  \sum_{k=1}^{\infty}a_k^i A_k=0,
\]
where $A_k$ denotes the residue of $1/F(z)$ at $z=a_k$. Thus, under the hypothesis made, the residues of $1/F(z)$ furnish a solution of system~\eqref{eq:14:13}. Now
\[
 A_k=\frac{-a_k}
 {(1-a_k/a_1)\cdots(1-a_k/a_{k-1})(1-a_k/a_{k+1})\cdots},
\]
and this is the same solution to which the principle of sections would lead.

\numberedparagraph{15}
Fourier's system~\eqref{eq:3:4} falls in principle within the type just considered. Indeed, omitting its first equation and putting
\[
  (2m-1)^2=b_m,\qquad (2m-1)^2a_m=x_m,
\]
the system becomes
\[
  \sum_{m=1}^{\infty}b_m^i x_m=0
  \qquad(i=0,1,2,\ldots),
\]
and the $b_m$ increase indefinitely. Let us calculate $F(z)$. We have
\[
  F(z)=\prod_{m=1}^{\infty}\left[1-\frac z{(2m-1)^2}\right]
      =\cos\frac\pi2\sqrt z.
\]
The residues of $1/F(z)$ are\translatornote{The French original prints the factor $\pi/4$ in both places; it is corrected here to $4/\pi$, as follows by differentiating $F(z)=\cos((\pi/2)\sqrt z)$.}
\[
 B_m=\frac1{F'(b_m)}
 =-\frac4\pi\frac{\sqrt{b_m}}{\sin(\frac\pi2\sqrt{b_m})}
 =\pm\frac4\pi(2m-1).
\]
Thus, apart from a factor, they represent Fourier's solution.

% [Source printed page 17]
It is not, however, a solution of the kind Poincar\'e required: the series $\sum b_m^iB_m$ do not converge. This example shows very clearly that hypothesis~\eqref{eq:14:14} cannot be omitted.

\numberedparagraph{16}
Here is a very interesting observation of Poincar\'e, which at first sight will appear paradoxical. Suppose that a solution $x_1=\alpha_1,x_2=\alpha_2,\ldots$ of system~\eqref{eq:14:13} has been found, and suppose also that the $\alpha_k$ do not vanish. Let $p$ be any positive integer. It is clear that $x_1=a_1^p\alpha_1,x_2=a_2^p\alpha_2,\ldots$ also satisfy system~\eqref{eq:14:13}. The same is true of $x_1=f(a_1)\alpha_1,x_2=f(a_2)\alpha_2,\ldots$, where $f(z)$ is any polynomial. So far this is correct. Now consider, in place of polynomials, the entire function
\[
  f(z)=c_0+c_1z+c_2z^2+\cdots.
\]
Multiply the equations~\eqref{eq:14:13}, beginning with the $(i+1)$st, successively by $c_0,c_1,c_2,\ldots$, and add term by term; this would give
\begin{equation}
  \sum_{k=1}^{\infty}a_k^i f(a_k)\alpha_k=0
  \qquad(i=0,1,2,\ldots).
  \tag{15}\label{eq:16:15}
\end{equation}
The following paradoxical consequence would then follow at once: if the reasoning were correct, every arbitrary sequence $\beta_1,\beta_2,\ldots$ would satisfy system~\eqref{eq:14:13}. Indeed, one can choose the entire function $f(z)$ so that at $z=a_1,a_2,a_3,\ldots$ it takes respectively the values $\beta_1/\alpha_1,\beta_2/\alpha_2,\beta_3/\alpha_3,\ldots$. To obtain such a function, let $F(z)$ be the entire function already considered, vanishing at $z=a_1,a_2,\ldots$ and having no other zeros. By Mittag-Leffler's theorem, one can construct a meromorphic function $R(z)$ having the $a_k$ as simple poles with residues $\beta_k/[\alpha_kF'(a_k)]$, and no other poles. Thus $f(z)=F(z)R(z)$ is the required entire function. Applied to this $f(z)$, formulas~\eqref{eq:16:15} give
\[
  \sum_{k=1}^{\infty}a_k^i\beta_k=0
  \qquad(i=0,1,2,\ldots).
\]

% [Source printed page 18]
If the reasoning were legitimate, these equations would have to hold for every choice of the quantities $\beta_k$. Obviously it is not legitimate, for in general it is not permissible to add infinitely many series term by term.

To make the reasoning valid, assume that for the solution $x_1=\alpha_1,x_2=\alpha_2,\ldots$ the left-hand sides of equations~\eqref{eq:14:13} are absolutely convergent series. Put
\[
  \sum_{k=1}^{\infty}|a_k^i\alpha_k|=S_i
  \qquad(i=0,1,2,\ldots).
\]
Then whenever the series
\[
  |c_0|S_0+|c_1|S_1+|c_2|S_2+\cdots
\]
converges, our argument is valid, since termwise addition is permitted. Since $f(z)$ depends only on the quantities $a_k,\alpha_k,\beta_k$, its coefficients $c_k$ will be expressible in some definite way through those same quantities. Everything therefore reduces to the convergence of a certain series whose terms are well-determined functions of the $a_k,\alpha_k$, and $\beta_k$. When that series converges, the $\beta_k$ give a solution of system~\eqref{eq:14:13}. Thus we have connected the problem of solving system~\eqref{eq:14:13}, to a certain extent, with a simple question of convergence.

\numberedparagraph{17}
\begingroup
\let\OriginalFootnote\footnote
\newcount\FootnoteSelector
\FootnoteSelector=0
\long\def\footnote#1{%
  \advance\FootnoteSelector by 1\relax
  \ifnum\FootnoteSelector=1\relax
    \OriginalFootnote{#1}%
  \else
    \footnotemark{}%
    \gdef\BorelLongNote{#1}%
  \fi
}
The equations treated by Appell do not belong to the type just considered, but to the type
\[
  \sum_{k=-\infty}^{\infty}a_k^i x_k=0
  \qquad(i=0,1,2,\ldots),
\]
where $|a_k|$ increases toward $\infty$ as $k\to\infty$ and decreases toward $0$ as $k\to-\infty$. To treat this type, Poincar\'e uses a method analogous to that of \hyperref[sec:14]{no.~14}; the principal difference is that the function $F(z)$ having the $a_k$ as zeros is no longer an entire function, since the origin, being a limit point of the $a_k$, is necessarily an essential singular point.

The detailed study of Poincar\'e's note and of the supplementary remarks that followed it in a second note\footnote{Poincar\'e, ``Sur les d\'eterminants d'ordre infini,'' \emph{Bulletin de la Soci\'et\'e math\'ematique de France}, vol.~XIV (1886), pp.~77--90.} would require us to penetrate still further into the Theory of Functions. Besides, the matter there is not one of methods in the strict sense; rather, these are ideas still awaiting development.\footnote{The only progress in this direction is due to Borel. He was led to it by the following problem treated in his thesis, ``Sur quelques points de la Th\'eorie des fonctions,'' \emph{Annales de l'\'Ecole Normale sup\'erieure}, 3rd series, vol.~XII (1895), pp.~9--55: determine a power series $f(z)=\sum b_nz^n$ convergent for $z=1$ together with its derivatives, prescribed values $f^{(n)}(1)=F_n$ being given for the series and its derivatives at $z=1$. The essential idea of Borel's reasoning is this: to solve the equations $f^{(n)}(1)=F_n$, it is enough to know how to solve them with a certain approximation. Suppose a power series $g(z)=\sum b_nz^n$ has been found such that $|g^{(n)}(1)-F_n|<C$ for all $n$. Put $g^{(n)}(1)-F_n=c_n$; then $\sum(c_n/n!)(z-1)^n$ defines an entire function $h(z)$ and $f(z)=g(z)-h(z)$. Thus everything reduces to solving $f^{(n)}(1)=F_n$ with the stated approximation. Put $b_0=0$, $b_k=\pm1/k$ for $k=1,2,\ldots,n$, choosing the signs and $n$ so that $|b_1+\cdots+b_n-F_0|<1$. Next put $b_k=\pm1/k^2$ for $k=n+1,\ldots,n'$, choosing the signs and $n'$ so that $|b_1+2b_2+\cdots+n'b_{n'}-F_1|<1$. Then put $b_k=\pm1/k^3$ for $k=n'+1,\ldots,n''$, so that
\[
 |2b_2+2\cdot3b_3+\cdots+(n''-1)n''b_{n''}-F_2|<1.
\]
Continued in this fashion, the procedure leads to a function $g(z)$ such that $|g^{(n)}(1)-F_n|<3$. This follows readily from the divergence of $\sum1/k$ and the inequality $\sum1/k^2<2$. The reasoning assumes the data $F_n$ real, but this hypothesis is plainly inessential.

Borel applies an analogous procedure to the systems
\[
 \sum_{k=1}^{\infty}a_k^i x_k=A_i\qquad(i=1,2,\ldots),
\]
where $a_k$ increases indefinitely with $k$. The procedure consists in putting, according to the indices, $x_k=\pm1/k,\ \pm1/(ka_k),\ \pm1/(ka_k^2),\ldots$, so that the quantities $B_i=A_i-\sum a_k^i x_k$ remain bounded as a whole. Thus the general case may be reduced to the special case in which the quantities $A_i$ are bounded. In this special case Borel solves the system by slightly modifying Poincar\'e's idea from \hyperref[sec:14]{no.~14}. Let $F(z)$ have the $a_k$ as simple zeros, and choose $\theta_k$ so that the series $\sum|a_k^i/[F'(a_k)\theta_k]|$ converge. Construct an entire function $\theta(z)$ such that $\theta(a_k)=\theta_k$. On writing
\[
 \frac{F(z)\theta(z)}{z-a_k}=\sum c_n^{(k)}z^n,
\]
a solution is given by
\[
 x_k=\frac{\sum_i c_i^{(k)}A_i}{F'(a_k)\theta_k}.
\]}
\endgroup
\footnotetext{\fontsize{9pt}{9.8pt}\selectfont\BorelLongNote}

% [Source printed page 19]
% The preceding source footnote occupies the whole of printed page 19 and
% continues onto printed page 20.

% [Source printed page 20]
What I have said, however, applies only to the first note and to the first part of the second. In the second part, abruptly leaving the line of ideas followed thus far, Poincar\'e lays the first foundations of a theory of infinite determinants, which will form the subject of the \hyperref[chap:2]{next chapter}.

% [Source printed page 21]
\chapter{Infinite Determinants}\label{chap:2}

\section{Historical Remarks and Generalities}

\numberedparagraph{18}
The extension of algebraic methods to systems of linear equations with infinitely many unknowns naturally led to the consideration of infinite determinants. Allusions to these algorithms already occur in the work of K\"otteritzsch cited above. Their introduction is generally attributed to G.~W.~Hill. That eminent astronomer used them in a very important memoir on the motion of the lunar perigee,\footnote{Hill, \emph{On the Part of the Motion of the Lunar Perigee Which Is a Function of the Mean Motions of the Sun and Moon}, Cambridge, Massachusetts, U.S.A., 1877 (reprinted in \emph{Acta Mathematica}, vol.~VIII, 1886, p.~1). Compare also Adams's communication in the \emph{London Astronomical Society Monthly Notices}, vol.~XXXVIII (1877), p.~13; there he announces that his earlier study of the motion of the Moon's node had already led him to an infinite determinant, but he did not publish his results.} and it was through this memoir that Poincar\'e's attention was drawn to the matter. Hill considers the equation
\begin{equation}
  \frac{d^2w}{dt^2}+\theta w=0,
  \tag{1}\label{eq:18:1}
\end{equation}
where $\theta$ is a series of the form
\[
  \theta=\theta_0+2\theta_1\cos t+2\theta_2\cos2t+\cdots.
\]
Putting $\theta_{-n}=\theta_n$, one may also write
\begin{equation}
  \theta=\sum_{n=-\infty}^{\infty}\theta_ne^{int}.
  \tag{2}\label{eq:18:2}
\end{equation}

% [Source printed page 22]
The general integral of equation~\eqref{eq:18:1} may be put in the form
\[
  w=\sum_{n=-\infty}^{\infty}b_ne^{i(n+c)t},
\]
where the $b_n$ and $c$ are suitably chosen constants. Substituting this series for $w$ and series~\eqref{eq:18:2} for $\theta$ in equation~\eqref{eq:18:1}, and comparing coefficients, gives infinitely many equations
\begin{equation}
  \sum_{k=-\infty}^{\infty}\theta_{n-k}b_k-(n+c)^2b_n=0
  \qquad(n=-\infty,\ldots,+\infty),
  \tag{3}\label{eq:18:3}
\end{equation}
and one must determine the $b_n$ and $c$ so that these equations are satisfied.

Hill treats system~\eqref{eq:18:3} by the procedure most often used for finite systems, namely by determinants. He introduces \emph{infinite determinants}, applies to them the ordinary rules of determinant calculus, and his boldness is justified by success, since the results agree with observation. He did not, however, prove the legitimacy of his method. This gap was soon filled by Poincar\'e, who on this occasion developed the first principles of a theory of infinite determinants.\footnote{See the note cited at \hyperref[sec:17]{no.~17}, or the same author's \emph{M\'ethodes nouvelles de la M\'ecanique c\'eleste}, vol.~II, p.~260.} Poincar\'e's investigations were continued until recent times and greatly deepened by von Koch, to whom, one may say, almost all the essential results of the theory are due.\footnote{For a list of von Koch's works on infinite determinants, compare his report ``Sur les syst\`emes d'une infinit\'e d'\'equations lin\'eaires \`a une infinit\'e d'inconnues,'' \emph{Proceedings of the Congress of Mathematicians held at Stockholm}, 1909. Add to that list the two notes ``Sur un nouveau crit\`ere de convergence pour les d\'eterminants infinis,'' \emph{Arkiv f\"or matematik, astronomi och fysik}, vol.~VII (1911), no.~4, and ``Sur certains syst\`emes d'\'equations lin\'eaires \`a une infinit\'e d'inconnus,'' ibid., vol.~VIII (1912), no.~9, together with the communication mentioned at \hyperref[sec:12]{no.~12}.}

\numberedparagraph{19}
Consider the system
\begin{equation}
  \sum_{k=1}^{\infty}a_{ik}x_k=c_i
  \qquad(i=1,2,\ldots),
  \tag{4}\label{eq:19:4}
\end{equation}
and suppose that the principle of sections may legitimately be applied to it. Suppose further that, for all sufficiently large $n$,
% [Source printed page 23]
\[
  \Delta_n=
  \begin{vmatrix}
    a_{11}&\cdots&a_{1n}\\
    \vdots&&\vdots\\
    a_{n1}&\cdots&a_{nn}
  \end{vmatrix}\ne0.
\]
Then the unknown $x_k$ is given by
\[
  \lim_{n\to\infty}\frac{\Delta_n^{(k)}}{\Delta_n},
\]
where $\Delta_n^{(k)}$ denotes the result of replacing the terms of column $k$ of $\Delta_n$ by $c_1,\ldots,c_n$. It may happen that the denominator $\Delta_n$ itself tends, as $n\to\infty$, to a limit $\Delta$; if this limit is nonzero, one can assert that the numerator also tends to a limit $\Delta^{(k)}$ and that
\[
  x_k=\frac{\Delta^{(k)}}{\Delta}.
\]

This, however, is only the simplest case; the homogeneous system~\eqref{eq:18:3} does not belong to it and requires a more delicate discussion. To approach that system by infinite determinants, these determinants must be studied somewhat more closely. It should also be observed that even in the simpler case just considered we assumed \emph{a priori} that the principle of sections applies. This is not self-evident and may in fact fail; the limits $\Delta\ne0$ and $\Delta^{(k)}$ $(k=1,2,\ldots)$ may exist without the values $x_k=\Delta^{(k)}/\Delta$ satisfying system~\eqref{eq:19:4}. For example, if $a_{ik}=0$ for $k<i$ and $a_{ik}=1$ for $k\geq i$, while $c_i=(-1)^i$, then $\Delta=1$ and $\Delta^{(k)}=2(-1)^k$, which gives $x_k=2(-1)^k$; when these values are substituted in equations~\eqref{eq:19:4}, the left-hand sides diverge.\footnote{Cazzaniga gave a series of examples of this kind, showing the care required with infinite determinants: ``Sui determinanti d'ordine infinito,'' \emph{Annali di Matematica}, 2nd series, vol.~XXVI (1897), pp.~143--218; ``Intorno ad un tipo di determinanti nulli d'ordine infinito,'' ibid., 3rd series, vol.~I (1898), pp.~83--94.}

These remarks suffice to show that, if one wishes to base
% [Source printed page 24]
a theory of infinite systems of equations on the study of infinite determinants, the conditions under which one is working must first be clearly specified. Well-chosen hypotheses will make it possible to extend to infinite determinants the ordinary rules for determinants of finite order and consequently to apply them to the theory with which we are concerned.

Some of the simplest rules undoubtedly extend at once to all infinite determinants. For example, rows and columns may be interchanged; interchanging two rows amounts to multiplying the determinant by $-1$; multiplying all elements of one row by the same number amounts to multiplying the determinant itself by that number; the determinant vanishes if two of its rows coincide; it is unchanged when the elements of one row are increased by the corresponding elements of another, and so forth. All this follows immediately from the analogous properties of finite-order determinants.

\section{Normal Determinants}

\numberedparagraph{20}
Consider the doubly indexed array
\[
\begin{matrix}
  1+a_{11},&a_{12},&a_{13},&\ldots,\\
  a_{21},&1+a_{22},&a_{23},&\ldots,\\
  a_{31},&a_{32},&1+a_{33},&\ldots,\\
  \vdots&\vdots&\vdots&
\end{matrix}
\]
and suppose that the doubly infinite series $\sum_{i,k=1}^{\infty}|a_{ik}|$ converges. Put
\[
  \Lambda_k=\sum_{i=1}^{\infty}|a_{ik}|;
\]
by hypothesis $\sum_{k=1}^{\infty}\Lambda_k$ converges. A well-known convergence rule then implies that the products
\[
  \Pi_n=\prod_{k=1}^{n}(1+\Lambda_k)
\]
tend, as $n\to\infty$, to a limit $\Pi$.

% [Source printed page 25]
Now compare the two determinants
\[
 \Delta_n=
 \begin{vmatrix}
   1+a_{11}&a_{12}&\cdots&a_{1n}\\
   \vdots&\vdots&&\vdots\\
   a_{n1}&a_{n2}&\cdots&1+a_{nn}
 \end{vmatrix},
\]
\[
 \Delta_{n+p}=
 \begin{vmatrix}
   1+a_{11}&a_{12}&\cdots&a_{1,n+p}\\
   \vdots&\vdots&&\vdots\\
   a_{n+p,1}&a_{n+p,2}&\cdots&1+a_{n+p,n+p}
 \end{vmatrix}.
\]
They are sums of certain products of the quantities $a_{ik}$ with appropriate signs. The same products occur in absolute value in the expansions of $\Pi_n$ and $\Pi_{n+p}$ in powers of the $|a_{ik}|$. Moreover, $\Delta_{n+p}$ contains all terms of $\Delta_n$, and its other terms occur in absolute value in the expansion of $\Pi_{n+p}-\Pi_n$. Hence
\[
  |\Delta_n|\leq\Pi_n,
  \qquad
  |\Delta_{n+p}-\Delta_n|\leq\Pi_{n+p}-\Pi_n;
\]
consequently the series
\[
  \Delta_1+(\Delta_2-\Delta_1)+(\Delta_3-\Delta_2)+\cdots
\]
converges absolutely to a quantity $\Delta=\lim\Delta_n$ such that $|\Delta|\leq\Pi$.

\numberedparagraph{21}
Replace in our array the elements of column $k$ by quantities $c_i$, assumed bounded as a whole, that is, bounded in absolute value by some quantity $C$. Let $n>k$, and denote by $\Delta_n^{(k)}$ the result of applying the indicated modification to $\Delta_n$, and by $\Pi_n^{(k)}$ the result of suppressing the factor $1+\Lambda_k$ from $\Pi_n$. The terms of $\Delta_n^{(k)}$ and of $\Delta_{n+p}^{(k)}-\Delta_n^{(k)}$ occur, in absolute value and apart from a factor $|c_i|$, in the expansions of $\Pi_n^{(k)}$ and $\Pi_{n+p}^{(k)}-\Pi_n^{(k)}$. Hence
\[
 |\Delta_n^{(k)}|\leq C\Pi_n^{(k)},\qquad
 |\Delta_{n+p}^{(k)}-\Delta_n^{(k)}|
 \leq C(\Pi_{n+p}^{(k)}-\Pi_n^{(k)}).
\]
Thus the determinant $\Delta^{(k)}=\lim\Delta_n^{(k)}$ converges and $|\Delta^{(k)}|\leq C\Pi$.

In particular, put $c_i=1$ and $c_j=0$ for $j\ne i$; $\Delta^{(k)}$ then becomes what is called the minor
% [Source printed page 26]
$\left(\begin{smallmatrix}i\\k\end{smallmatrix}\right)$ of the determinant $\Delta_n$. Similarly, we shall call $\Delta^{(k)}$ the minor $\left(\begin{smallmatrix}i\\k\end{smallmatrix}\right)$ of $\Delta$. It follows at once from the analogous fact for finite determinants that the minor $\left(\begin{smallmatrix}i\\k\end{smallmatrix}\right)$ is unchanged if the elements of row $i$, except the $k$th, are replaced by zeros or by arbitrary quantities. This observation shows that rows and columns play the same role in forming minors.

We shall prove the following theorem, whose first part is the familiar Laplace rule.

\begin{theorem}
Whatever the quantities $c_i$, provided they are bounded as a whole,
\begin{equation}
  \Delta^{(k)}=\sum_{i=1}^{\infty}
  \left(\begin{matrix}i\\k\end{matrix}\right)c_i,
  \tag{5}\label{eq:21:5}
\end{equation}
and the series on the right converges absolutely. In particular, the series
\begin{equation}
  \sum_{i=1}^{\infty}
  \left|\left(\begin{matrix}i\\k\end{matrix}\right)\right|
  \tag{6}\label{eq:21:6}
\end{equation}
converges, and its sum does not exceed $\Pi$.
\end{theorem}

\begin{proof}
Absolute convergence of the series on the right of~\eqref{eq:21:5} follows immediately from convergence of~\eqref{eq:21:6}. Since the latter series has positive terms, it is enough to prove that its partial sums do not exceed a finite bound. I claim that they do not exceed $\Pi$. In $\Delta^{(k)}$ put
\[
 c_i=\operatorname{sgn}^{*}\!\left(\begin{matrix}i\\k\end{matrix}\right)
 \quad(i\leq n),\qquad c_i=0\quad(i>n),
\]
where $\operatorname{sgn}^{*}z$ denotes $0$ when $z=0$ and $|z|/z$ otherwise.\footnote{Thus, if $z=re^{i\varphi}$, then $\operatorname{sgn}^{*}z=e^{-i\varphi}$. For real positive or negative $z$, $\operatorname{sgn}^{*}z=\pm1$, according to the sign of $z$.} Then $C=1$, and $\Delta^{(k)}$ is the sum of the first $n$ terms of~\eqref{eq:21:6}; the inequality $|\Delta^{(k)}|\leq C\Pi$ therefore shows that this sum cannot exceed $\Pi$.

It remains to verify identity~\eqref{eq:21:5}. The sum $s_n$ of the
% [Source printed page 27]
first $n$ terms of the series on the right is the result of replacing by zeros in $\Delta^{(k)}$ the $c_i$ beginning with $c_{n+1}$. Hence the difference $\Delta^{(k)}-s_n$ is what $\Delta^{(k)}$ becomes after its first $n$ quantities $c_i$ have been set equal to zero. Denote this last determinant by $\Delta^{(k,n)}$, and let $\Delta_n^{(k,n)}$ be the finite determinant formed from its first $n^2$ elements. Then
\[
  |\Delta^{(k,n)}-\Delta_n^{(k,n)}|\leq C(\Pi-\Pi_n).
\]
For $n\geq k$, however, $\Delta_n^{(k,n)}$ vanishes because it contains an entire column of zeros. Hence
\[
 |\Delta^{(k)}-s_n|=|\Delta^{(k,n)}|\leq C(\Pi-\Pi_n).
\]
As $n\to\infty$, $\Pi-\Pi_n\to0$, and therefore $s_n\to\Delta^{(k)}$, as required.
\end{proof}

We record the following special cases of~\eqref{eq:21:5}:
\begin{equation}
 \left(\begin{matrix}j\\k\end{matrix}\right)
 +\sum_{i=1}^{\infty}a_{ij}
 \left(\begin{matrix}i\\k\end{matrix}\right)
 =\begin{cases}\Delta,&k=j,\\0,&k\ne j.
 \end{cases}
 \tag{7}\label{eq:21:7}
\end{equation}
Convergence of~\eqref{eq:21:6} also immediately implies absolute convergence of the double series
\begin{equation}
 \sum_{i,k=1}^{\infty}a_{jk}c_i
 \left(\begin{matrix}i\\k\end{matrix}\right).
 \tag{8}\label{eq:21:8}
\end{equation}
Throughout the preceding discussion, rows and columns may exchange their roles. In particular,
\begin{equation}
 \left(\begin{matrix}i\\j\end{matrix}\right)
 +\sum_{k=1}^{\infty}a_{jk}
 \left(\begin{matrix}i\\k\end{matrix}\right)
 =\begin{cases}\Delta,&i=j,\\0,&i\ne j.
 \end{cases}
 \tag{9}\label{eq:21:9}
\end{equation}

\numberedparagraph{22}
Infinite determinants $\Delta$ of the type just considered are called \emph{normal determinants}. Historically, Poincar\'e studied the determinants $\Delta$ and $\Delta^{(k)}$ under the additional restrictive condition $a_{ii}=0$ for all $i$. It is nearly evident that the type considered here is easily reduced to that more special type. For example, when $a_{ii}\ne-1$, one need only divide column $i$ by $1+a_{ii}$. Such a reduction is unnecessary, however; von Koch observed that Poincar\'e's method,
% [Source printed page 28]
with only slight modifications, extends to all normal determinants, and this is precisely the method we have just presented.

The preceding considerations, and those that follow, extend immediately to \emph{infinite determinants in four directions}, that is, determinants corresponding to arrays whose indices $i,k$ range from $-\infty$ to $+\infty$. In this case, the infinite determinant means the limit, as $m\to-\infty$ and $n\to+\infty$, of the finite determinant whose indices range from $m$ to $n$. Hill's system~\eqref{eq:18:3} leads to determinants of this kind.

\section{Application of Normal Determinants to Systems of Equations. Minors of Higher Order and von Koch's Theorem}

\numberedparagraph{23}
After these preliminaries, consider the system of equations
\begin{equation}
\left\{
\begin{aligned}
 (1+a_{11})x_1+a_{12}x_2+\cdots&=c_1,\\
 a_{21}x_1+(1+a_{22})x_2+\cdots&=c_2,\\
 &\hspace{1em}\cdots
\end{aligned}
\right.
\tag{10}\label{eq:23:10}
\end{equation}
where the $x$ are the unknowns and the $a$ and $c$ are given quantities.

We suppose that $\sum_{i,k}|a_{ik}|=\sum_k\Lambda_k=\Lambda$ converges, and ask whether system~\eqref{eq:23:10} can be satisfied by quantities $x_k$ that in addition remain bounded in absolute value. Under these conditions the left-hand sides of equations~\eqref{eq:23:10} are absolutely convergent series.

A necessary condition is supplied by
\[
 |c_i|\leq|x_i|+\sum_{k=1}^{\infty}|a_{ik}x_k|
 \leq(1+\Lambda)\sup_k|x_k|.
\]
It says that the quantities $|c_i|$ must likewise remain below some finite bound $C$. Assume this condition satisfied.

We must distinguish two cases according as the determinant of the system vanishes or does not vanish.

\emph{First case: $\Delta\ne0$.} Under the conditions stated,
% [Source printed page 29]
the system has the solution
\begin{equation}
  x_1=\frac{\Delta^{(1)}}{\Delta},\qquad
  x_2=\frac{\Delta^{(2)}}{\Delta},\qquad\ldots.
  \tag{11}\label{eq:23:11}
\end{equation}
Indeed, using~\eqref{eq:21:7} and~\eqref{eq:21:9}, and observing that the double series~\eqref{eq:21:8} converges absolutely, we have
\begin{align*}
 \Delta^{(j)}+\sum_{k=1}^{\infty}a_{jk}\Delta^{(k)}
 &=\sum_{i=1}^{\infty}c_i
   \left(\begin{matrix}i\\j\end{matrix}\right)
 +\sum_{k=1}^{\infty}a_{jk}\sum_{i=1}^{\infty}c_i
   \left(\begin{matrix}i\\k\end{matrix}\right)\\
 &=\sum_{i=1}^{\infty}c_i
 \left[
   \left(\begin{matrix}i\\j\end{matrix}\right)
   +\sum_{k=1}^{\infty}a_{jk}
   \left(\begin{matrix}i\\k\end{matrix}\right)
 \right]
 =c_j\Delta;
\end{align*}
that is, the quantities~\eqref{eq:23:11} satisfy system~\eqref{eq:23:10}. Moreover, the inequalities
\[
 |\Delta^{(k)}|\leq\sum_{i=1}^{\infty}|c_i|
 \left|\left(\begin{matrix}i\\k\end{matrix}\right)\right|
 \leq C\Pi
\]
show that these quantities remain bounded in absolute value.

Solution~\eqref{eq:23:11} is unique: more precisely, system~\eqref{eq:23:10} has no other solution with the required boundedness property. Otherwise the homogeneous system
\begin{equation}
\left\{
\begin{aligned}
 (1+a_{11})x_1+a_{12}x_2+\cdots&=0,\\
 a_{21}x_1+(1+a_{22})x_2+\cdots&=0,\\
 &\hspace{1em}\cdots
\end{aligned}
\right.
\tag{12}\label{eq:23:12}
\end{equation}
would also have a bounded solution $x_1,x_2,\ldots$ in which at least one $x_i$ did not vanish. But the following double series converge absolutely, and hence
\[
 0=\sum_{i=1}^{\infty}
 \left(\begin{matrix}i\\j\end{matrix}\right)
 \left[x_i+\sum_{k=1}^{\infty}a_{ik}x_k\right]
 =\sum_{k=1}^{\infty}x_k
 \left[
 \left(\begin{matrix}k\\j\end{matrix}\right)
 +\sum_{i=1}^{\infty}a_{ik}
 \left(\begin{matrix}i\\j\end{matrix}\right)
 \right]
 =x_j\Delta,
\]
so $x_j=0$ for every $j$.

\emph{Second case: $\Delta=0$.} The homogeneous system~\eqref{eq:23:12} has the solution
\begin{equation}
 x_1=\left(\begin{matrix}i\\1\end{matrix}\right),\qquad
 x_2=\left(\begin{matrix}i\\2\end{matrix}\right),\qquad\ldots,
 \tag{13}\label{eq:23:13}
\end{equation}
% [Source printed page 30]
where $i$ is any index. This is an immediate consequence of formulas~\eqref{eq:21:9}.

Suppose $i$ and $k$ are chosen so that the minor $\left(\begin{smallmatrix}i\\k\end{smallmatrix}\right)\ne0$. Then solution~\eqref{eq:23:13} differs from the evident solution $x_1=x_2=\cdots=0$ and is, moreover, unique up to a factor. Indeed, suppose the system has the solution $x_1=\xi_1,x_2=\xi_2,\ldots$, and replace its $i$th equation by
\[
  x_k=\xi_k.
\]
The determinant of the new system is $\left(\begin{smallmatrix}i\\k\end{smallmatrix}\right)\ne0$; hence the system belongs to the first case and has the unique solution
\[
 x_1=\frac{\left(\begin{smallmatrix}i\\1\end{smallmatrix}\right)}
           {\left(\begin{smallmatrix}i\\k\end{smallmatrix}\right)}\xi_k,
 \qquad
 x_2=\frac{\left(\begin{smallmatrix}i\\2\end{smallmatrix}\right)}
           {\left(\begin{smallmatrix}i\\k\end{smallmatrix}\right)}\xi_k,
 \qquad\ldots.
\]

\numberedparagraph{24}
It may happen, however, that all the minors $\left(\begin{smallmatrix}i\\k\end{smallmatrix}\right)$ vanish. To discuss this case, one must first extend the notion of a minor by introducing, with von Koch, minors
\[
 \left(\begin{matrix}i_1,\ldots,i_r\\k_1,\ldots,k_r\end{matrix}\right)
\]
of order $r$. By this we mean the result of replacing in $\Delta$ the $r$ elements lying simultaneously in row $i_p$ and column $k_p$ $(p=1,\ldots,r)$ by $1$, and all other elements of these rows (or, equivalently, of these columns) by zero. Our general convergence rule applies also to these new determinants.

\begin{vonkochtheorem}
There is always a minor of $\Delta$ of some finite order that does not vanish. In particular, if $m$ is sufficiently large, the minor
\[
 \left(\begin{matrix}1,\ldots,m\\1,\ldots,m\end{matrix}\right)
\]
is certainly nonzero.
\end{vonkochtheorem}

\begin{proof}
Choose $m$ such that
\[
  R_m=\sum_{i=m+1}^{\infty}\Lambda_i<\frac12.
\]
Then
\[
 \left|
 \left(\begin{matrix}1,\ldots,m\\1,\ldots,m\end{matrix}\right)
 \right|
  \geq1-\sum_{i=m+1}^{\infty}\Lambda_i
       -\sum_{i,k=m+1}^{\infty}\Lambda_i\Lambda_k-\cdots
 \geq1-R_m-R_m^2-\cdots
 =\frac{1-2R_m}{1-R_m}>0.
\]
\end{proof}

% [Source printed page 31]
\numberedparagraph{25}
By the theorem just proved, in the case now under consideration there is a number $r$ such that $\Delta$ and all its minors of orders $1,\ldots,r-1$ vanish, while at least one minor of order $r$ is nonzero. Let
\[
 D=\left(\begin{matrix}i_1,\ldots,i_r\\k_1,\ldots,k_r\end{matrix}\right)\ne0
\]
be such a minor. In~\eqref{eq:23:12}, replace the $r$ equations corresponding to $i_1,\ldots,i_r$ by
\[
  x_{k_1}=\xi_{k_1},\qquad x_{k_2}=\xi_{k_2},\qquad\ldots,
  \qquad x_{k_r}=\xi_{k_r}.
\]
The determinant of the modified system is $D\ne0$, so the system belongs to the first case and has the unique solution
\begin{equation}
 x_k=\frac1D\sum_{p=1}^{r}
 \left(\begin{matrix}
 i_1,\ldots,i_r\\
 k_1,\ldots,k_{p-1},k,k_{p+1},\ldots,k_r
 \end{matrix}\right)\xi_{k_p}.
 \tag{14}\label{eq:25:14}
\end{equation}

Give arbitrary values to the indeterminates $\xi_{k_1},\ldots,\xi_{k_r}$. The resulting $x_k$ satisfy the homogeneous equations~\eqref{eq:23:12}, except perhaps those corresponding to $i_1,\ldots,i_r$, which were temporarily excluded. But the minors
\[
 \left(\begin{matrix}
 i_1,\ldots,i_{p-1},i,i_{p+1},\ldots,i_r\\
 k_1,\ldots,k_{p-1},k,k_{p+1},\ldots,k_r
 \end{matrix}\right),
\]
of order $r$ with respect to $\Delta$, are at the same time first-order minors of the determinant obtained by suppressing the other $r-1$ selected rows and columns. Hence, by~\eqref{eq:21:9},
\begin{align*}
 &\left(\begin{matrix}i_1,\ldots,i,\ldots,i_r\\
 k_1,\ldots,i,\ldots,k_r\end{matrix}\right)
 +\sum_k a_{ik}
 \left(\begin{matrix}i_1,\ldots,i,\ldots,i_r\\
 k_1,\ldots,k,\ldots,k_r\end{matrix}\right)\\
 &\hspace{5em}=
 \left(\begin{matrix}i_1,\ldots,i_{p-1},i_{p+1},\ldots,i_r\\
 k_1,\ldots,k_{p-1},k_{p+1},\ldots,k_r\end{matrix}\right)=0,
\end{align*}
where $i$ and the summation indices are first assumed distinct from the fixed indices. The summation may be extended to all $k$, and the formula is valid for all $i$, if by definition one sets
% [Source printed page 32]
\[
 \left(\begin{matrix}i_1,\ldots,i_r\\k_1,\ldots,k_r\end{matrix}\right)=0
\]
whenever two $i$-indices or two $k$-indices coincide. Substituting successively $k=k_1,\ldots,k_r$ in this formula, multiplying respectively by $\xi_{k_1},\ldots,\xi_{k_r}$, and adding term by term, equation~\eqref{eq:25:14} gives
\[
  x_i+\sum_{k=1}^{\infty}a_{ik}x_k=0
\]
for every $i$.

Thus homogeneous system~\eqref{eq:23:12} has the solutions~\eqref{eq:25:14}, depending on $r$ parameters, and no others.

Now consider the nonhomogeneous system~\eqref{eq:23:10} and suppose that it has a solution $x'_1,x'_2,\ldots$. Any solution of homogeneous system~\eqref{eq:23:12} may plainly be added to it. Add in particular the solution corresponding to $\xi_{k_1}=-x'_{k_1},\ldots,\xi_{k_r}=-x'_{k_r}$; the resulting solution $x''_1,x''_2,\ldots$ of~\eqref{eq:23:10} has $x''_{k_1}=\cdots=x''_{k_r}=0$. Thus, whenever~\eqref{eq:23:10} has solutions, one of them has $x_{k_1}=\cdots=x_{k_r}=0$, and all the others are obtained from it by adding the solutions~\eqref{eq:25:14} of the homogeneous system.

It remains, then, to determine a solution $(x_k)$ of~\eqref{eq:23:10} with $x_{k_1}=\cdots=x_{k_r}=0$. The other unknowns then satisfy the system obtained from~\eqref{eq:23:10} by suppressing the equations numbered $i_1,\ldots,i_r$ and the unknowns $x_{k_1},\ldots,x_{k_r}$. The determinant of this modified system is $D\ne0$, so it has exactly one solution. Hence~\eqref{eq:23:10} is solvable if and only if the values $x_{k_1}=\cdots=x_{k_r}=0$, together with the values of the remaining $x_k$ furnished by the modified system, satisfy the equations of ranks $i_1,\ldots,i_r$. This gives $r$ relations among the right-hand sides of~\eqref{eq:23:10}; these relations express the necessary and sufficient condition for solvability.

This condition can be put into several forms analogous to
% [Source printed page 33]
those known for ordinary finite systems. We shall not dwell on them. The essential fact resulting from the preceding considerations is that the classical theory of determinants extends to normal infinite determinants and to the corresponding systems of equations.

\section{Absolutely Convergent Determinants}

\numberedparagraph{26}
In seeking to extend the preceding results to more general cases, von Koch soon observed that a very slight modification of the reasoning yields a considerable extension.\footnote{Von Koch, ``Sur la convergence des d\'eterminants d'ordre infini,'' \emph{Bihang till K.~Svenska Vetenskaps-Akademiens Handlingar}, vol.~XXII (1896), no.~4, pp.~1--31.} The majorant expression $\Pi$ is in fact much richer in terms than the determinant $\Delta$: it contains all products of the $|a_{ik}|$ whose factors belong to distinct columns. Von Koch's procedure consists in constructing a majorant expression better adapted to the structure of determinants. He starts from the following definition.

Let a doubly indexed array $(A_{ik})$ be given, and suppose the product $P$ of all diagonal elements $A_{ii}$ converges absolutely (equivalently, suppose $\sum|A_{ii}-1|$ converges). Form a family of new products $P'$ by permuting in $P$ the second indices in every possible way, assigning to each resulting product the sign $+$ or $-$ according as the number of transpositions required to pass from the initial product $P$ to the product in question is even or odd. If the series $\Delta$ formed from all these products converges absolutely even when every $A_{ii}$ is replaced by $1+|A_{ii}-1|$, we agree to regard $\Delta$ as the determinant of the $A_{ik}$ and to say that the determinant \emph{converges absolutely}.

\numberedparagraph{27}
Put $a_{ii}=A_{ii}-1$ and $a_{ik}=A_{ik}$ for $i\ne k$. Expanding $P,P'$ in products of the quantities $a$, von Koch's hypothesis plainly amounts to absolute convergence of the sum of all these products. Since the terms
% [Source printed page 34]
of an absolutely convergent sum may be arranged arbitrarily, $\Delta$ admits the expansions
\[
 \Delta_1+(\Delta_2-\Delta_1)+(\Delta_3-\Delta_2)+\cdots,
\]
and
\[
 1+\sum_{i=1}^{\infty}a_{ii}
 +\sum_{i,j=1}^{\infty}
 \begin{vmatrix}a_{ii}&a_{ij}\\a_{ji}&a_{jj}\end{vmatrix}
 +\sum_{i,j,k=1}^{\infty}
 \begin{vmatrix}
 a_{ii}&a_{ij}&a_{ik}\\
 a_{ji}&a_{jj}&a_{jk}\\
 a_{ki}&a_{kj}&a_{kk}
 \end{vmatrix}+\cdots.
\]
Furthermore, if $\Delta$ converges absolutely, so do its principal minors
\[
 \left(\begin{matrix}i_1,\ldots,i_r\\i_1,\ldots,i_r\end{matrix}\right),
\]
for all products of the quantities $a$ composing these minors also enter the expansion of $\Delta$.

These principal minors cannot all vanish. Indeed, as $m\to\infty$, the sum of the products of $a$ composing
$\left(\begin{smallmatrix}1,\ldots,m\\1,\ldots,m\end{smallmatrix}\right)$ tends to zero; hence that minor tends to $1$ and is nonzero for all sufficiently large $m$.

For nonprincipal minors, it should be observed that absolute convergence of $\Delta$ does not always imply absolute convergence of its minors. But if none of the quantities $a_{ik}$ $(i\ne k)$ vanishes, absolute convergence of every minor follows at once from that of $\Delta$ itself.

Suppose in particular that $\Delta$ and, for a fixed index $k$, the minors $\left(\begin{smallmatrix}i\\k\end{smallmatrix}\right)$ converge absolutely. By suitably arranging the terms one then obtains Laplace's rule
\[
 \Delta=\left(\begin{matrix}k\\k\end{matrix}\right)
 +\sum_{i=1}^{\infty}a_{ik}
  \left(\begin{matrix}i\\k\end{matrix}\right).
\]
These results make it easy to see that the determinant method applies to infinite systems of linear equations whenever the determinants and minors entering the calculation converge absolutely.

\numberedparagraph{28}
Von Koch gave a necessary and sufficient condition for absolute convergence of $\Delta$. Consider the
% [Source printed page 35]
cyclic products of the quantities $a_{ik}$; by a cyclic product of order $r$ we mean any product of the form
\[
 a_{i_1i_2}a_{i_2i_3}\cdots a_{i_{r-1}i_r}a_{i_ri_1},
\]
where $i_1,\ldots,i_r$ are distinct. By a cyclic product of first order we mean a quantity $a_{ii}$. Since every cyclic product occurs in the expansion of $\Delta$, absolute convergence of the series formed from all cyclic products is necessary for absolute convergence of $\Delta$. This necessary condition is also sufficient. Indeed, form the infinite product $\prod(1+|p|)$ over all cyclic products $p$. Its expansion contains in absolute value every term in the expansion of $\Delta$. If $\sum|p|$ converges, then $\prod(1+|p|)$ converges also, and consequently $\Delta$ converges absolutely.

\begin{theorem}
The infinite determinant $\Delta$ converges absolutely if and only if the series formed from the cyclic products of the quantities $a_{ik}$ converges absolutely.
\end{theorem}

\numberedparagraph{29}
It is almost evident that normal determinants converge absolutely, as do the determinants obtained from them by replacing the elements of one column by quantities bounded as a whole. For normal determinants the infinite product $\prod_{i,k}(1+|a_{ik}|)$ converges, and its expansion contains in absolute value all cyclic products. For determinants of the second type, a corresponding majorant is obtained by adjoining the factor $1+C$ and suppressing the factors corresponding to the index $k$. Their minors belong to the same type and therefore also converge absolutely.

In seeking more general rules, von Koch constructed a succession of increasingly precise criteria for absolute convergence of $\Delta$ and all its minors. We shall give only one of the simplest, which is important for applications to integral equations.\footnote{Von Koch, ``Sur la convergence des d\'eterminants infinis,'' \emph{Rendiconti del Circolo matematico di Palermo}, vol.~XXVIII (1909), pp.~255--266.}

\begin{theorem}
For $\Delta$ and all its minors to converge absolutely, it is sufficient that the two series
\[
  \sum_{i=1}^{\infty}|a_{ii}|,
  \qquad
  \sum_{i,k=1}^{\infty}|a_{ik}|^2
\]
  converge
% [Source printed page 36]
  absolutely.
\end{theorem}
\begin{proof}
The minors belong to the same type as $\Delta$, so it suffices to prove absolute convergence of $\Delta$. By \hyperref[sec:28]{the theorem just established}, it is enough to prove convergence of the sum of the moduli of the cyclic products. Those of first order have been covered explicitly by the hypothesis; below we consider only orders greater than $1$.

Following von Koch's proof, define \(\sigma\) by the following formula.\translatornote{The French original prints \(\sigma^2\) on the left.  The estimates that follow require \(\sigma=(\sum_{i\ne k}|a_{ik}|^2)^{1/2}\), equivalently \(\sigma^2=\sum_{i\ne k}|a_{ik}|^2\).}
\[
  \sigma=\left(\sum_{i\ne k}|a_{ik}|^2\right)^{1/2}.
\]
Denote by $|i_1,i_2,\ldots,i_r|$ the absolute value of the cyclic product $a_{i_1i_2}a_{i_2i_3}\cdots a_{i_ri_1}$, and by $|i_1,i_2,\ldots,i_r;i_{r+1}|$ that of $a_{i_1i_2}a_{i_2i_3}\cdots a_{i_ri_{r+1}}$. The indices are required to be distinct; whenever this condition fails, both symbols are defined to be zero.

We shall use the Lagrange--Cauchy inequality\footnote{Cauchy, \emph{Cours d'Analyse de l'\'Ecole royale Polytechnique}, Paris, 1824, Note~II; \emph{Oeuvres}, 2nd series, vol.~III.}
\begin{equation}
 \left|\sum_i a_ib_i\right|^2
 \leq\left(\sum_i|a_ib_i|\right)^2
 \leq\sum_i|a_i|^2\sum_i|b_i|^2,
 \tag{15}\label{eq:29:15}
\end{equation}
valid for arbitrary numbers. It follows from Lagrange's identity
\[
 \sum_i|a_i|^2\sum_i|b_i|^2
 =\left(\sum_i|a_ib_i|\right)^2
 +\sum_{i<k}(|a_ib_k|-|a_kb_i|)^2.
\]

% [Source printed page 37]
The sums may contain infinitely many terms, provided the two series on the right of~\eqref{eq:29:15} converge; absolute convergence of the other series then follows immediately.

First suppose $\sigma<1$. From~\eqref{eq:29:15},
\[
 \left(\sum_j|i,j;k|\right)^2
 \leq\sum_j|i,j|^2\sum_j|j,k|^2.
\]
For $i=k$ this gives
\[
 \sum_{i,j}|i,j|\leq\sigma^2,
\]
while for arbitrary $i$,
\[
 \sum_k\left(\sum_j|i,j;k|\right)^2
 \leq\sigma^2\sum_j|i,j|^2.
\]
Consequently the inequalities
\begin{equation}
 \sum_{i_1,\ldots,i_n}|i_1,i_2,\ldots,i_n|\leq\sigma^n,
 \tag{16}\label{eq:29:16}
\end{equation}
and
\begin{equation}
 \sum_j\left(\sum_{i_2,\ldots,i_n}
 |i_1,i_2,\ldots,i_n;j|\right)^2
 \leq\sigma^{2n-2}\sum_j|i_1,j|^2
 \tag{17}\label{eq:29:17}
\end{equation}
hold for $n=2$. Suppose they hold for a given $n$. Then
\begin{align*}
 \sum_{i_1,\ldots,i_{n+1}}|i_1,\ldots,i_{n+1}|
 &=\sum_{i_1,i_{n+1}}
 \left(
   |i_{n+1},i_1|
   \sum_{i_2,\ldots,i_n}|i_1,\ldots,i_n;i_{n+1}|
 \right)\\
 &\leq
 \sum_{i_1}
 \left[
   \sum_{i_{n+1}}|i_{n+1},i_1|^2
   \sum_{i_{n+1}}
   \left(\sum_{i_2,\ldots,i_n}|i_1,\ldots,i_n;i_{n+1}|\right)^2
 \right]^{1/2}\\
 &\leq
 \left(
   \sum_{i_1,i_{n+1}}|i_{n+1},i_1|^2\,
   \sigma^{2n-2}
   \sum_{i_1,i_{n+1}}|i_1,i_{n+1}|^2
 \right)^{1/2}\\
 &\leq\sigma^{n+1}.
\end{align*}
% [Source printed page 38]
Similarly,
\begin{align*}
 &\sum_j\left(\sum_{i_2,\ldots,i_{n+1}}
 |i_1,\ldots,i_{n+1};j|\right)^2\\
 &=\sum_j\left[
   \sum_{i_2}\left(
     |i_1,i_2|
     \sum_{i_3,\ldots,i_{n+1}}|i_2,\ldots,i_{n+1};j|
   \right)
 \right]^2\\
 &\quad\leq
 \sum_{i_2}|i_1,i_2|^2
 \sum_{j,i_2}\left(\sum_{i_3,\ldots,i_{n+1}}
 |i_2,\ldots,i_{n+1};j|\right)^2\\
 &\quad\leq
 \sum_{i_2}|i_1,i_2|^2\,
 \sigma^{2n-2}\sum_{i_2,j}|i_2,j|^2\\
 &=\sigma^{2n}\sum_k|i_1,k|^2.
\end{align*}
Thus, if~\eqref{eq:29:16} and~\eqref{eq:29:17} hold for $n$, they hold for $n+1$; since they hold for $n=2$, they hold for all $n$.

By~\eqref{eq:29:16}, the sum of the moduli of all cyclic products of order greater than $1$ is bounded by the convergent series
\[
 \sigma^2+\sigma^3+\cdots=\frac{\sigma^2}{1-\sigma};
\]
hence $\Delta$ converges absolutely.

We assumed $\sigma<1$. The general case reduces to this special one. When $\sum_{i,k}|a_{ik}|^2$ converges, choose $m$ so that
\[
 \sum_{i=m+1}^{\infty}\sum_{k=1}^{\infty}|a_{ik}|^2<\frac12,
\]
and then choose a sufficiently small positive $t$ so that\footnote{Compare, in addition to von Koch's memoir, the Hungarian thesis of Sz\'asz, \emph{A v\'egtelen determin\'ansok elm\'elet\'ehez}, Budapest, 1911, where these systems and their determinants are studied independently of the general theory of absolutely convergent determinants.}
\[
 t^2\sum_{i=1}^{m}\sum_{k=1}^{\infty}|a_{ik}|^2<\frac12.
\]

% [Source printed page 39]
Multiply by $t$ the elements of the first $m$ rows of $\Delta$. It follows immediately from the definition in \hyperref[sec:26]{no.~26} that $\Delta$ converges absolutely if and only if the modified determinant does. For the latter $\sigma<1$, so it converges absolutely; therefore $\Delta$ also converges absolutely.
\end{proof}

\numberedparagraph{30}
The preceding considerations show that the results obtained for systems~\eqref{eq:23:10} and~\eqref{eq:23:12} may be extended by requiring of the unknowns $x_k$ and the data $a_{ik},c_i$ that the series
\[
 \sum_k|x_k|^2,\qquad \sum_i|a_{ii}|,\qquad
 \sum_{i,k}|a_{ik}|^2,\qquad \sum_i|c_i|^2
\]
 converge. The most important application of such systems is to integral equations. We shall discuss it later and shall see that all the present conditions arise quite naturally except convergence of $\sum|a_{ii}|$. This last condition can be removed by multiplying the equations of the system by suitable factors.

Indeed, suppose $\sum_i|c_i|^2$ and $\sum_{i,k}|a_{ik}|^2$, and hence in particular $\sum_i|a_{ii}|^2$, converge. Applying the inequality
\[
 |e^{-a}(1+a)-1|
 =\left|-\frac{a^2}{2}+\frac{a^3}{3}-\frac{a^4}{8}+\cdots\right|
 \leq\frac{|a|^2}{2}e^{|a|},
\]
and observing that $a_{ii}\to0$, one immediately obtains convergence of
\[
 \sum_i|e^{-a_{ii}}c_i|^2,\qquad
 \sum_{i,k}|e^{-a_{ii}}a_{ik}|^2,
 \qquad
 \sum_i|e^{-a_{ii}}(1+a_{ii})-1|.
\]
Thus, if each equation is multiplied by the corresponding quantity $e^{-a_{ii}}$, the modified system is of the type already considered.

\numberedparagraph{31}
Finally, suppose the data $a_{ik}$ are functions $a_{ik}(\lambda)$ of a parameter $\lambda$, holomorphic in a domain $D$, and that the conditions imposed above are satisfied uniformly with respect to $\lambda$. For example, for normal determinants we assume that $\sum_{i,k}|a_{ik}(\lambda)|$ converges uniformly. The expansions given for $\Delta$
% [Source printed page 40]
then converge uniformly, and the function $\Delta(\lambda)$ so defined is holomorphic in $\lambda$.

When $\lambda$ is not a zero of $\Delta(\lambda)$, the system formed with the $a_{ik}(\lambda)$ has the well-determined solution
\[
 x_k(\lambda)=\frac{\Delta^{(k)}(\lambda)}{\Delta(\lambda)}.
\]
As a function of $\lambda$, $x_k(\lambda)$ is meromorphic in the interior of the domain, and its poles are zeros of $\Delta(\lambda)$.

Expanding $\Delta(\lambda)$ as a series in products of the functions $a_{ik}(\lambda)$ gives an absolutely and uniformly convergent series; the general rules concerning series of analytic functions may therefore be applied to $\Delta(\lambda)$. In particular, successive derivatives with respect to $\lambda$ may be calculated, and it is easy to specify conditions under which they have the same form as for finite-order determinants. Suppose, for example, that the series
\[
 \sum_{i,k}\left|\frac{d a_{ik}(\lambda)}{d\lambda}\right|
\]
converges uniformly in $D$. Then the expansion of $\Delta(\lambda)$ converges absolutely and uniformly there, and $d\Delta(\lambda)/d\lambda$ is obtained by termwise differentiation. Since the preceding series also converges and the first-order minors of $\Delta(\lambda)$ remain bounded as a whole, the series
\begin{equation}
 \sum_{i,k}
 \left(\begin{matrix}i\\k\end{matrix}\right)
 \frac{d a_{ik}}{d\lambda}
 \tag{18}\label{eq:31:18}
\end{equation}
converges absolutely. The minors too may be expanded in products of the functions $a_{ik}(\lambda)$; even after every term is replaced by its absolute value, the resulting minors remain bounded as a whole. Thus, on expanding~\eqref{eq:31:18} in products of $a_{ik}(\lambda)$ and $d a_{ik}(\lambda)/d\lambda$, the resulting series converges absolutely. Its terms may therefore be rearranged, in particular so as to give the series obtained by differentiating the expansion of $\Delta(\lambda)$ term by term. Consequently, series~\eqref{eq:31:18} represents
% [Source printed page 41]
$d\Delta(\lambda)/d\lambda$. Finally, if
\[
 A_{ik}(\lambda)=
 \begin{cases}
 1+a_{ik}(\lambda),&i=k,\\
 a_{ik}(\lambda),&i\ne k,
 \end{cases}
\]
then $dA_{ik}(\lambda)/d\lambda=da_{ik}(\lambda)/d\lambda$. Replacing the latter derivatives by the former in~\eqref{eq:31:18} gives the customary formula for $d\Delta/d\lambda$.

% [Source printed page 42]
\chapter{Essay Toward a General Theory}\label{chap:3}

\section{Introduction}

\numberedparagraph{32}
To apply the classical method of determinants to infinite systems, it was necessary to impose more or less restrictive conditions on the data; and it must be admitted that these restrictions were required by the method, not by the problem. In what follows we shall adopt a much more general point of view, due in principle to E.~Schmidt.\footnote{E.~Schmidt, ``Ueber die Aufl\"osung linearer Gleichungen mit abz\"ahlbar unendlich vielen Unbekannten,'' \emph{Rendiconti del Circolo Matematico di Palermo}, vol.~XXV (1908), pp.~53--77.} We shall begin by specifying the nature of the admissible solutions $(x_k)$; as for the data, our sole hypothesis will be that, for every admissible system $(x_k)$, the series occurring on the left-hand sides of the equations must converge. The principal question will be the \emph{existence of solutions}; only secondarily shall we be concerned with the apparatus that produces them.

The most important case from the standpoint of applications is that in which one requires $\sum_k |x_k|^2$ to converge. To treat it by the method of determinants, we had to suppose that the double series $\sum_{i,k}|a_{ik}|^2$ and the series $\sum_i|c_i|^2$ converge. As we shall see, Schmidt's theory assumes only the convergence of the singly indexed series $\sum_k|a_{ik}|^2$. An intermediate theory was based by Hilbert on the study of quadratic forms and bilinear forms
% [Source printed page 43]
in infinitely many variables.\footnote{Hilbert, ``Grundz\"uge einer allgemeinen Theorie der linearen Integralgleichungen, 4.~Mitteilung,'' \emph{Nachrichten von der K\"oniglichen Gesellschaft der Wissenschaften zu G\"ottingen} (1906), pp.~157--227. Hilbert's six memoirs on integral equations have just been collected in a volume (Leipzig, 1912). Cf. also Hilbert, ``Wesen und Ziele einer Analysis der unendlichvielen unabh\"angigen Variabeln,'' \emph{Rendiconti del Circolo Matematico di Palermo}, vol.~XXVII (1909), pp.~59--74.} It is much more effective than the method of determinants, but less general than Schmidt's. As regards historical order, it should be observed that Hilbert's theory predates Schmidt's; it also predates the study of determinants for which $\sum_{i,k}|a_{ik}|^2$ converges, but it is later than the general theory of infinite and absolutely convergent determinants.

\hyperref[chap:4]{Chapters IV} and \hyperref[chap:5]{V} are devoted to Hilbert's theory. In the present chapter we shall set forth Schmidt's theory in a slightly generalized form.

\section{The Fundamental Inequalities}

\numberedparagraph{33}
We shall have occasion to start from the inequality
\begin{equation}
 \left|\sum_{k=1}^{n}a_kb_k\right|
 \leq
 \left(\sum_{k=1}^{n}|a_k|^{\frac{p}{p-1}}\right)^{\frac{p-1}{p}}
 \left(\sum_{k=1}^{n}|b_k|^p\right)^{\frac1p}
 \qquad (p>1),
 \tag{1}\label{eq:33:1}
\end{equation}
which concerns arbitrary quantities $a_k,b_k$, real or complex. Already established and used by Cauchy for $p=2$,\footnote{Cf. \hyperref[sec:29]{no.~29}.} it was extended by H\"older to all $p>1$.\footnote{H\"older, ``Ueber einen Mittelwertsatz,'' \emph{Nachrichten von der K\"oniglichen Gesellschaft der Wissenschaften zu G\"ottingen} (1889), pp.~38--47.}

It is enough to prove inequality~\eqref{eq:33:1} for positive quantities $a_k,b_k$; it will then hold, \emph{a fortiori}, for arbitrary $a$ and $b$. We may also suppose the $a$'s and $b$'s strictly positive, regarding the cases in which one or another of these quantities vanishes as limiting cases. Finally, we may suppose $n=2$, for once the inequality has been proved in this particular case, it extends to every $n$ by induction. It therefore remains to prove the inequality
% [Source printed page 44]
\begin{equation}
 a_1b_1+a_2b_2
 -\left(a_1^{\frac{p}{p-1}}+a_2^{\frac{p}{p-1}}\right)^{\frac{p-1}{p}}
  \left(b_1^p+b_2^p\right)^{\frac1p}
 \leq0,
 \tag{2}\label{eq:33:2}
\end{equation}
the quantities $a_1,b_1,a_2,b_2$ being supposed strictly positive.

The proof is readily carried out by the usual methods of differential calculus. Keeping $a_1,a_2,b_1$ fixed, let $b_2$ vary. The derivative with respect to $b_2$ of the expression on the left vanishes only when
\begin{equation}
 b_2=\left(\frac{a_2}{a_1}\right)^{\frac1{p-1}}b_1;
 \tag{3}\label{eq:33:3}
\end{equation}
it is positive for every smaller value of $b_2$ and negative for every larger value. Hence the indicated value of $b_2$ makes our expression a maximum; substitution shows that this maximum is zero. Inequality~\eqref{eq:33:2} is therefore true; moreover, equality holds only in case~\eqref{eq:33:3}.

For positive quantities $a_k,b_k$, Jensen has given a very interesting geometric interpretation of inequality~\eqref{eq:33:1}. He considers the curve $y=x^p$ $(x>0)$. This curve is convex, that is, its chords lie above it. Let $P_1,\ldots,P_n$ be the points of the curve whose abscissae are
\[
 a_1^{\frac1{1-p}}b_1,\ldots,a_n^{\frac1{1-p}}b_n.
\]
Assign to these points the positive masses
\[
 a_1^{\frac{p}{p-1}},\ldots,a_n^{\frac{p}{p-1}}.
\]
Then inequality~\eqref{eq:33:1} expresses the fact that the center of gravity of the system lies above the curve; equality can hold only when the points $P_k$ coincide.\footnote{Jensen, ``Sur les fonctions convexes et les in\'egalit\'es entre les valeurs moyennes,'' \emph{Acta Mathematica}, vol.~XXX (1906), pp.~175--193.}

Inequality~\eqref{eq:33:1} extends immediately to infinite series. Indeed, suppose that the two series
\[
 A=\sum|a_k|^{\frac{p}{p-1}},
 \qquad B=\sum|b_k|^p
\]
converge. Then, by~\eqref{eq:33:1}, the partial sums of the series $\sum|a_kb_k|$ do not exceed the bound $A^{(p-1)/p}B^{1/p}$. Consequently,
% [Source printed page 45]
the series $\sum a_kb_k$ converges absolutely, and
\begin{equation}
 \left|\sum_{k=1}^{\infty}a_kb_k\right|
 \leq
 \left(\sum_{k=1}^{\infty}|a_k|^{\frac{p}{p-1}}\right)^{\frac{p-1}{p}}
 \left(\sum_{k=1}^{\infty}|b_k|^p\right)^{\frac1p}
 \qquad(p>1).
 \tag{4}\label{eq:33:4}
\end{equation}

\numberedparagraph{34}
We shall also use the inequality
\begin{equation}
 \left(\sum_{k=1}^{\infty}|a_k+b_k|^p\right)^{\frac1p}
 \leq
 \left(\sum_{k=1}^{\infty}|a_k|^p\right)^{\frac1p}
 +\left(\sum_{k=1}^{\infty}|b_k|^p\right)^{\frac1p}
 \qquad(p>1).
 \tag{5}\label{eq:34:5}
\end{equation}
It is obtained by an analysis entirely analogous to the preceding one, but it also follows immediately from inequality~\eqref{eq:33:4}. We shall merely present the case of four positive quantities $a_1,a_2,b_1,b_2$, since one passes from this particular case to the general case by reasoning exactly as above.\footnote{For a finite number of terms, inequality~\eqref{eq:34:5} was established by Minkowski: \emph{Diophantische Approximationen}, 1907, p.~95. That work treats the important theory which its author calls the geometry of numbers. In the order of ideas of that theory, inequality~\eqref{eq:34:5} expresses the convexity of the $n$-dimensional domain $\sum_{k=1}^{n}|x_k|^p\leq1$. Notice that, in the arguments to follow, inequality~\eqref{eq:34:5} could in most cases be replaced by the nearly evident inequality
\[
 \sum_{k=1}^{\infty}|a_k+b_k|^p
 \leq2^p\sum_{k=1}^{\infty}|a_k|^p
      +2^p\sum_{k=1}^{\infty}|b_k|^p.
\]}

Applying~\eqref{eq:33:2}, we write
\begin{equation}
\begin{aligned}
 &(a_1+b_1)^p+(a_2+b_2)^p\\
 &\quad=(a_1+b_1)^{p-1}a_1+(a_2+b_2)^{p-1}a_2
       +(a_1+b_1)^{p-1}b_1+(a_2+b_2)^{p-1}b_2\\
 &\quad\leq
 \bigl[(a_1+b_1)^p+(a_2+b_2)^p\bigr]^{\frac{p-1}{p}}
 \left[(a_1^p+a_2^p)^{\frac1p}+(b_1^p+b_2^p)^{\frac1p}\right];
\end{aligned}
\tag{6}\label{eq:34:6}
\end{equation}
equality is possible only if
\[
 a_2=\frac{a_2+b_2}{a_1+b_1}a_1,
 \qquad
 b_2=\frac{a_2+b_2}{a_1+b_1}b_1,
\]
that is, only if $a_1b_2=a_2b_1$.

% [Source printed page 46]
Dividing by
\[
 \bigl[(a_1+b_1)^p+(a_2+b_2)^p\bigr]^{\frac{p-1}{p}},
\]
we obtain the inequality to be proved:
\[
 \bigl[(a_1+b_1)^p+(a_2+b_2)^p\bigr]^{\frac1p}
 \leq(a_1^p+a_2^p)^{\frac1p}+(b_1^p+b_2^p)^{\frac1p};
\]
equality holds only in the case just described.

Here is a second form of our inequality, obtained by replacing $a_k$ in it by $a_k-b_k$:
\begin{equation}
 \left(\sum_{k=1}^{\infty}|a_k-b_k|^p\right)^{\frac1p}
 \geq
 \left(\sum_{k=1}^{\infty}|a_k|^p\right)^{\frac1p}
 -\left(\sum_{k=1}^{\infty}|b_k|^p\right)^{\frac1p}.
 \tag{7}\label{eq:34:7}
\end{equation}

Let us further specify the conditions under which equality holds in~\eqref{eq:34:5}. It is necessary and sufficient that the $b$'s be proportional to the $a$'s,
\[
 a_1:b_1=a_2:b_2=\cdots,
\]
and that these ratios be positive.

These conditions are plainly sufficient. Let us show that they are also necessary. When the second condition is not fulfilled, one can increase the values of the terms $|a_k+b_k|^p$ without changing the absolute values of the $a$'s and $b$'s. Thus this condition must be fulfilled; and when it is, we may suppose without loss of generality that the $a$'s and $b$'s are positive. But in that case equality can hold only when the $b$'s are proportional to the $a$'s. Indeed, suppose, to fix ideas, that $a_1:b_1\ne a_2:b_2$, that is, $a_1b_2\ne a_2b_1$. We have just seen that in this case the strict inequality sign applies in~\eqref{eq:34:6}; hence one can increase
\[
 (a_1+b_1)^p+(a_2+b_2)^p
\]
without changing either $a_1^p+a_2^p$ or $b_1^p+b_2^p$. The left-hand side of~\eqref{eq:34:5} will thus have been increased without changing the right-hand side; consequently equality is excluded in the case under consideration.

\section{The Problem. Landau's Theorem}

\numberedparagraph{35}
Consider the system formed by a finite or countably infinite number of equations
\begin{equation}
 \sum_{k=1}^{\infty}a_{ik}x_k=c_i
 \qquad(i=1,2,\ldots).
 \tag{8}\label{eq:35:8}
\end{equation}

% [Source printed page 47]
Given a number $p>1$, chosen arbitrarily once and for all, we propose to determine a solution $(x_k)$ of the system in such a way that the series
\begin{equation}
 \sum_{k=1}^{\infty}|x_k|^p
 \tag{9}\label{eq:35:9}
\end{equation}
converges.

As for the data, we make only one hypothesis. Under this hypothesis the equations~\eqref{eq:35:8} must have a meaning. But before carrying out the solution, we know nothing about the quantities $x_k$. We are therefore led to suppose that the left-hand sides of~\eqref{eq:35:8} converge for every system $(x_k)$ for which the series~\eqref{eq:35:9} converges. This is equivalent to assuming that, for every $i$, the sums
\[
 \sum_{k=1}^{\infty}|a_{ik}|^{\frac{p}{p-1}}
\]
converge. Indeed, Hellinger and Toeplitz proved for $p=2$,\footnote{Hellinger and Toeplitz, ``Grundlagen f\"ur eine Theorie der unendlichen Matrizen,'' \emph{Nachrichten von der K\"oniglichen Gesellschaft der Wissenschaften zu G\"ottingen} (1906), pp.~351--355.} and Landau for arbitrary $p>1$,\footnote{Landau, ``Ueber einen Konvergenzsatz,'' \emph{Nachrichten von der K\"oniglichen Gesellschaft der Wissenschaften zu G\"ottingen} (1907), pp.~25--27.} that the series
\begin{equation}
 \sum_{k=1}^{\infty}a_kx_k
 \tag{10}\label{eq:35:10}
\end{equation}
cannot converge for every system $(x_k)$ for which the series~\eqref{eq:35:9} converges unless
\begin{equation}
 \sum_{k=1}^{\infty}|a_k|^{\frac{p}{p-1}}
 \tag{11}\label{eq:35:11}
\end{equation}
converges. Conversely, when this last condition is fulfilled, inequality~\eqref{eq:33:1} establishes the absolute convergence of the series~\eqref{eq:35:10}.

\begin{landautheorem}
If the series~\eqref{eq:35:10} converges for every system $(x_k)$ for which the series~\eqref{eq:35:9} converges, then the series~\eqref{eq:35:11} converges.
\end{landautheorem}
\begin{proof}
Suppose that the latter series diverges. We shall define a system $(x_k)$ for which~\eqref{eq:35:9}
% [Source printed page 48]
converges and~\eqref{eq:35:10} diverges. Put
\[
 s_0=0,
 \qquad
 s_n=\sum_{k=1}^{n}|a_k|^{\frac{p}{p-1}};
\]
then $s_n\to\infty$. Divide the sequence of indices $0,1,2,\ldots$ into groups
\[
 0;\quad 1,2,\ldots,n_1;\quad n_1+1,\ldots,n_2;\quad\ldots,
\]
in such a way that
\[
 s_{n_\nu}-s_{n_{\nu-1}}\geq2^{\frac{\nu}{p-1}}.
\]
For an index $k$ belonging to the group ending with $n_\nu$, put
\[
 x_k=
 \frac{|a_k|^{\frac1{p-1}}\operatorname{sgn}^{*}a_k}
 {s_{n_\nu}-s_{n_{\nu-1}}}.\footnote{For the notation $\operatorname{sgn}^{*}$, cf. the remark in \hyperref[sec:21]{no.~21}.}
\]
It follows that
\[
 \sum_{k=1}^{n_\nu}a_kx_k=\nu,
\]
and series~\eqref{eq:35:10} diverges. On the other hand,
\[
 \sum_{k=1}^{n_\nu}|x_k|^p
 =\sum_{\mu=1}^{\nu}
   \bigl(s_{n_\mu}-s_{n_{\mu-1}}\bigr)^{1-p}
 \leq\sum_{\mu=1}^{\nu}\frac1{2^\mu};
\]
therefore series~\eqref{eq:35:9} converges.
\end{proof}

\section{A Necessary Condition}

\numberedparagraph{36}
Suppose that system~\eqref{eq:35:8} admits a solution such that
\begin{equation}
 \sum_{k=1}^{\infty}|x_k|^p\leq M^p.
 \tag{12}\label{eq:36:12}
\end{equation}
Let $\mu_1,\ldots,\mu_n$ be arbitrary numbers, real or complex. The values $(x_k)$ also satisfy the equation
\[
 \sum_{k=1}^{\infty}\left(\sum_{i=1}^{n}\mu_i a_{ik}\right)x_k
 =\sum_{i=1}^{n}\mu_i c_i,
\]
which is simply a linear combination of the first $n$ equations of~\eqref{eq:35:8}.

% [Source printed page 49]
Applying inequality~\eqref{eq:33:4}, we obtain
\begin{equation}
\begin{aligned}
 \left|\sum_{i=1}^{n}\mu_i c_i\right|
 &=\left|\sum_{k=1}^{\infty}
       \left(\sum_{i=1}^{n}\mu_i a_{ik}\right)x_k\right|\\
 &\leq
 \left(\sum_{k=1}^{\infty}
       \left|\sum_{i=1}^{n}\mu_i a_{ik}\right|^{\frac{p}{p-1}}
 \right)^{\frac{p-1}{p}}
 \left(\sum_{k=1}^{\infty}|x_k|^p\right)^{\frac1p}.
\end{aligned}
\tag{13}\label{eq:36:13}
\end{equation}
By~\eqref{eq:36:12}, it follows that
\begin{equation}
 \left|\sum_{i=1}^{n}\mu_i c_i\right|
 \leq M
 \left(\sum_{k=1}^{\infty}
       \left|\sum_{i=1}^{n}\mu_i a_{ik}\right|^{\frac{p}{p-1}}
 \right)^{\frac{p-1}{p}}.
 \tag{14}\label{eq:36:14}
\end{equation}

Thus we have established a necessary condition for system~\eqref{eq:35:8} to admit a solution satisfying~\eqref{eq:36:12}: inequality~\eqref{eq:36:14} must hold for every $n$ and all real or complex numbers $\mu_i$. (When the number of equations~\eqref{eq:35:8} is finite, $n$ may be fixed and taken equal to the number of equations.)

We shall see that this condition is also sufficient.

\section{The Condition Is Also Sufficient: The Case of Finitely Many Equations}

\numberedparagraph{37}
First consider the particular case of finitely many equations
\begin{equation}
 \sum_{k=1}^{\infty}a_{ik}x_k=c_i
 \qquad(i=1,2,\ldots,n).
 \tag{15}\label{eq:37:15}
\end{equation}
Suppose the condition is fulfilled. Put
\[
 A_k(\mu_1,\ldots,\mu_n)=\sum_{i=1}^{n}\mu_i a_{ik},
 \qquad
 A(\mu_1,\ldots,\mu_n)
 =\sum_{k=1}^{\infty}|A_k(\mu_1,\ldots,\mu_n)|^{\frac{p}{p-1}},
\]
and
\[
 C(\mu_1,\ldots,\mu_n)
 =\left|\sum_{i=1}^{n}\mu_i c_i\right|^{\frac{p}{p-1}}.
\]

% [Source printed page 50]
The function $C(\mu_1,\ldots,\mu_n)$ is continuous in the $n$ variables $\mu_1,\ldots,\mu_n$. It is also continuous as a function of the $2n$ real variables
\[
 \mu_1',\mu_1'',\ldots,\mu_n',\mu_n''
 \qquad(\mu_i=\mu_i'+\mu_i''\sqrt{-1}).
\]
Moreover, with respect to these $2n$ variables it has continuous first derivatives. To simplify the calculation, we shall consider only the combinations
\[
 \frac{\partial C}{\partial\mu_i'}
 -\sqrt{-1}\frac{\partial C}{\partial\mu_i''}
\]
of these derivatives. We have
\begin{equation}
 \frac{\partial C}{\partial\mu_i'}
 -\sqrt{-1}\frac{\partial C}{\partial\mu_i''}
 =\frac{p}{p-1}c_i C^{\frac1p}
  \operatorname{sgn}^{*}\!\left(\sum_{j=1}^{n}\mu_jc_j\right).
 \tag{16}\label{eq:37:16}
\end{equation}
Similarly,
\begin{equation}
 \frac{\partial |A_k|^{\frac{p}{p-1}}}{\partial\mu_i'}
 -\sqrt{-1}\frac{\partial |A_k|^{\frac{p}{p-1}}}{\partial\mu_i''}
 =\frac{p}{p-1}a_{ik}|A_k|^{\frac1{p-1}}
  \operatorname{sgn}^{*}A_k.
 \tag{17}\label{eq:37:17}
\end{equation}

As for the function $A(\mu_1,\ldots,\mu_n)$, first observe that the series defining it converges for every system of values $\mu_1,\ldots,\mu_n$. This follows by successively applying inequality~\eqref{eq:34:5} to the sequences $(\mu_1a_{1k}),\ldots,(\mu_na_{nk})$, with the exponent $p$ there replaced by $p/(p-1)$, which is permissible since $p/(p-1)>1$. The same inequality also proves that the series defining $A(\mu_1,\ldots,\mu_n)$ converges uniformly in every bounded domain of the $\mu$'s; hence $A(\mu_1,\ldots,\mu_n)$ is a continuous function of the variables $\mu_i',\mu_i''$.

Differentiate the series defining $A(\mu_1,\ldots,\mu_n)$ term by term with respect to $\mu_i'$ (or $\mu_i''$). The absolute value of the $k$th term of the resulting series does not exceed the right-hand side of~\eqref{eq:37:17}. Consequently,
\begin{align*}
 \frac{p}{p-1}\sum_{k=1}^{\infty}|a_{ik}|\,|A_k|^{\frac1{p-1}}
 &\leq\frac{p}{p-1}
 \left(\sum_{k=1}^{\infty}|a_{ik}|^{\frac{p}{p-1}}\right)^{\frac{p-1}{p}}
 \left(\sum_{k=1}^{\infty}|A_k|^{\frac{p}{p-1}}\right)^{\frac1p}\\
 &=\frac{p}{p-1}A(\mu_1,\ldots,\mu_n)^{\frac1p}
 \left(\sum_{k=1}^{\infty}|a_{ik}|^{\frac{p}{p-1}}\right)^{\frac{p-1}{p}}
\end{align*}
is a majorant for the differentiated series. Since the majorant converges uniformly in every bounded domain of the $\mu$'s, the differentiated series does likewise. It follows from a well-known theorem on differentiation of series that the
% [Source printed page 51]
function $A(\mu_1,\ldots,\mu_n)$ has continuous partial derivatives obtained by differentiating term by term. Hence
\begin{equation}
 \frac{\partial A}{\partial\mu_i'}
 -\sqrt{-1}\frac{\partial A}{\partial\mu_i''}
 =\frac{p}{p-1}\sum_{k=1}^{\infty}
 a_{ik}|A_k|^{\frac1{p-1}}\operatorname{sgn}^{*}A_k.
 \tag{18}\label{eq:37:18}
\end{equation}

\numberedparagraph{38}
These preliminaries established, suppose for the moment that the left-hand sides of equations~\eqref{eq:37:15} are independent. This hypothesis may be formulated by requiring that $A(\mu_1,\ldots,\mu_n)$ cannot vanish unless all the $\mu$'s vanish. The restriction is not essential; we shall remove it at the end of the argument. For the moment it allows us to conclude that, when $A$ is bounded, the quantities $\mu$ remain bounded as well. Indeed, otherwise there would be an infinite sequence of systems $(\mu_i)$ such that, for at least one index $i$, the values $(\mu_i)$ would grow beyond every bound, while the function $A$ remained bounded. The homogeneity of $A$ would then yield a sequence of systems $(\mu_i)$ for each of which the largest value $|\mu_i|$ was exactly $1$, while the corresponding values of $A$ tended to zero. By Weierstrass's theorem, such a sequence contains a subsequence converging to a definite system $(\mu_i^{(0)})$. The values $\mu_i^{(0)}$ do not all vanish, since the largest $|\mu_i^{(0)}|$ is exactly $1$. On the other hand, continuity of $A$ would imply
\[
 A(\mu_1^{(0)},\ldots,\mu_n^{(0)})=0,
\]
contrary to the hypothesis.

The functions $A$ and $C$ are continuous in the $\mu$'s, and when $A$ is bounded the $\mu$'s remain bounded; hence $C$ remains bounded as well. Consequently, when $C$ is allowed to vary under the condition $A=1$, it attains its maximum for certain values $\mu_1,\ldots,\mu_n$. The classical method of differential calculus gives, for these values of the $\mu$'s, the equations
\begin{equation}
\left\{
\begin{aligned}
 &c_i\,C(\mu_1,\ldots,\mu_n)^{\frac1p}
 \operatorname{sgn}^{*}\!\left(\sum_{j=1}^{n}\mu_jc_j\right)\\
 &\qquad=\lambda\sum_{k=1}^{\infty}a_{ik}
 |A_k(\mu_1,\ldots,\mu_n)|^{\frac1{p-1}}
 \operatorname{sgn}^{*}A_k(\mu_1,\ldots,\mu_n)
 \quad(i=1,2,\ldots,n),
\end{aligned}
\right.
\tag{19}\label{eq:38:19}
\end{equation}
and
\begin{equation}
 A(\mu_1,\ldots,\mu_n)=1.
 \tag{20}\label{eq:38:20}
\end{equation}

% [Source printed page 52]
Multiplying equations~\eqref{eq:38:19} by $\mu_i$ and adding, we obtain
\[
 \lambda=C(\mu_1,\ldots,\mu_n).
\]
Consequently, for the values of the $\mu$'s under consideration,
\[
 \left(\sum_{j=1}^{n}\mu_jc_j\right)
 \sum_{k=1}^{\infty}a_{ik}|A_k(\mu_1,\ldots,\mu_n)|^{\frac1{p-1}}
 \operatorname{sgn}^{*}A_k(\mu_1,\ldots,\mu_n)
 =c_i
 \quad(i=1,2,\ldots,n).
\]
But this is precisely the solution of system~\eqref{eq:37:15} which we have just obtained. Indeed, with $\mu_1,\ldots,\mu_n$ denoting the extremal values in question, put
\begin{equation}
 x_k^{(n)}=
 |A_k(\mu_1,\ldots,\mu_n)|^{\frac1{p-1}}
 \operatorname{sgn}^{*}A_k(\mu_1,\ldots,\mu_n)
 \sum_{i=1}^{n}\mu_i c_i
 \qquad(k=1,2,\ldots).
 \tag{21}\label{eq:38:21}
\end{equation}
The quantities $x_k^{(n)}$ so defined satisfy equations~\eqref{eq:37:15}. Moreover, it follows from~\eqref{eq:36:14}, \eqref{eq:38:20}, and~\eqref{eq:38:21} that
\begin{equation}
 \sum_{k=1}^{\infty}|x_k^{(n)}|^p
 =\left|\sum_{i=1}^{n}\mu_i c_i\right|^p
 \leq M^p,
 \tag{22}\label{eq:38:22}
\end{equation}
that is, under the stated conditions, system~\eqref{eq:37:15} admits a solution satisfying~\eqref{eq:36:12}.

Inequality~\eqref{eq:38:22} may be replaced by a more precise relation. Let $M^{(n)}$ be the least of the numbers $M$ satisfying inequality~\eqref{eq:36:14} for every choice of $\mu_1,\ldots,\mu_n$. Then $[M^{(n)}]^{p/(p-1)}$ is the maximum value of $C(\mu_1,\ldots,\mu_n)$ under condition~\eqref{eq:38:20}. This maximum is attained for the values of the $\mu$'s which furnished the solution $(x_k^{(n)})$. Hence
\begin{equation}
 \sum_{k=1}^{\infty}|x_k^{(n)}|^p
 =\left|\sum_{i=1}^{n}\mu_i c_i\right|^p
 =[M^{(n)}]^p.
 \tag{23}\label{eq:38:23}
\end{equation}
On the other hand, by inequality~\eqref{eq:36:13}, every solution $(x_k)$ of system~\eqref{eq:37:15} satisfies
\[
 \sum_{k=1}^{\infty}|x_k|^p
 \geq\left[\frac{C(\mu_1,\ldots,\mu_n)}{A(\mu_1,\ldots,\mu_n)}\right]^{p-1}
 =[M^{(n)}]^p.
\]

% [Source printed page 53]
Thus the solution found, $(x_k^{(n)})$, has a second extremal property: among all solutions of system~\eqref{eq:37:15}, it is the one for which the sum $\sum|x_k|^p$ is as small as possible.

I say that $(x_k^{(n)})$ is the only solution having the property just observed. Indeed, let $(x_k)$ be another such solution. Inequality~\eqref{eq:34:5} gives
\[
 \sum_{k=1}^{\infty}\left|\frac{x_k^{(n)}+x_k}{2}\right|^p
 \leq
 \left[
  \frac12\left(\sum_{k=1}^{\infty}|x_k^{(n)}|^p\right)^{\frac1p}
 +\frac12\left(\sum_{k=1}^{\infty}|x_k|^p\right)^{\frac1p}
 \right]^p
 =[M^{(n)}]^p.
\]
Since $((x_k^{(n)}+x_k)/2)$ also satisfies system~\eqref{eq:37:15}, equality must hold. In \hyperref[sec:34]{no.~34}, while discussing inequality~\eqref{eq:34:5}, we also specified the conditions for equality; apply the criteria found there to the present case. We must have
\[
 x_1^{(n)}:x_1=x_2^{(n)}:x_2=\cdots,
\]
and these ratios must be positive. But, on the other hand,
\[
 \sum_{k=1}^{\infty}|x_k^{(n)}|^p
 =\sum_{k=1}^{\infty}|x_k|^p.
\]
Thus the ratio in question is $1$, unless system~\eqref{eq:37:15} is homogeneous.\translatornote{The French original reads ``unless the system is not homogeneous''; the argument and the following sentence require ``unless the system is homogeneous.''} The homogeneous case is evident: its extremal solution is $x_1=x_2=\cdots=0$.

It remains to remove the restriction that the left-hand sides of equations~\eqref{eq:37:15} be independent. Inequality~\eqref{eq:36:14} assures us that, whenever a linear combination of the left-hand sides vanishes identically, the same combination of the right-hand sides also vanishes. We may therefore suppress certain equations so that the remaining system is equivalent to the original system and at the same time is of the type considered. It is clear that reducing the system in this way does not alter the minimum of the $M$'s.

\numberedparagraph{39}
In the foregoing, we connected the solution of system~\eqref{eq:37:15} with another problem: to maximize the expression $C$ while the $\mu$'s vary under the condition $A=1$. I shall try in a few
% [Source printed page 54]
words to explain the order of ideas leading from the first problem to the second. Let us forget for a moment the preceding developments and suppose only that system~\eqref{eq:37:15} admits solutions for which $\sum|x_k|^p$ converges, and that among these solutions there is one which makes this sum a minimum; that is, suppose there is a system $(x_k^{(n)})$ which minimizes the function $\sum|x_k|^p$ of the variables $x_k$, the $x$'s varying in accordance with conditions~\eqref{eq:37:15}. The classical method of differential calculus shows that this particular solution $x_k^{(n)}$ must have the form~\eqref{eq:38:21}, the values of the $\mu$'s still remaining to be determined. Equations for the unknowns $\mu$ are obtained by substituting expressions~\eqref{eq:38:21} in equations~\eqref{eq:37:15}.

The resulting equations are linear when $p=2$; we shall return later to this particular case. But in the general case they are quite complicated. In that case one must be content to prove the \emph{existence} of a solution; that will suffice to infer the existence of the minimum solution $(x_k^{(n)})$ of~\eqref{eq:37:15}. This is what we did by following an indirect route, one suggested by inequality~\eqref{eq:36:13}. Indeed, let $\mu_1^{(n)},\ldots,\mu_n^{(n)}$ be the values to be determined which serve to define the $x_k^{(n)}$. Substitution of the expressions for $x_k^{(n)}$ in~\eqref{eq:36:13} gives
\[
 \frac{C(\mu_1,\ldots,\mu_n)}{A(\mu_1,\ldots,\mu_n)}
 \leq
 \frac{C(\mu_1^{(n)},\ldots,\mu_n^{(n)})}
      {A(\mu_1^{(n)},\ldots,\mu_n^{(n)})},
\]
whatever the values $\mu_1,\ldots,\mu_n$. Since, on the other hand, by~\eqref{eq:38:20},\translatornote{The French original refers here to equation~(15); the asserted equality is equation~(20).}
\[
 A(\mu_1^{(n)},\ldots,\mu_n^{(n)})=1,
\]
the $\mu_i^{(n)}$ must solve the maximum problem from which we started.

Let us add that, when $c_1=1$ and the other $c$'s vanish, our method is approximately the one Schmidt used for $p=2$ in the work cited in \hyperref[sec:32]{no.~32}. In this case our constrained maximum problem is immediately transformed into an unconstrained minimum problem: minimize the expression
\[
 \left(
  \sum_{k=1}^{\infty}
  \left|a_{1k}-\sum_{i=2}^{n}\mu_i a_{ik}\right|^{\frac{p}{p-1}}
 \right)^{\frac{p-1}{p}}.
\]
The minimum value of this expression permits an interesting geometric interpretation. It may be regarded
% [Source printed page 55]
in infinite-dimensional space as the distance from the point with coordinates $a_{1k}$ $(k=1,2,\ldots)$ to the linear variety determined by the $n-1$ points $(a_{2k}),\ldots,(a_{nk})$, the distance between two points $(b_k)$ and $(c_k)$ being measured by
\[
 \left(\sum_{k=1}^{\infty}|b_k-c_k|^{\frac{p}{p-1}}\right)^{\frac{p-1}{p}}.
\]

\section{Theorems Concerning Sequences of Systems
\texorpdfstring{\((y_k)\)}{(y-k)}}

\numberedparagraph{40}
Before passing to the general case, let us establish several theorems which we shall use.

Consider the infinite sequence of quantities $y_1^{(n)},y_2^{(n)},\ldots$. We suppose that these quantities vary with $n$, but in such a way that constantly
\[
 \sum_{k=1}^{\infty}|y_k^{(n)}|^p\leq G^p,
\]
where $G$ denotes a positive number independent of $n$. Suppose further that, as $n\to\infty$, each $y_k^{(n)}$ tends to a definite value $y_k^*$. We shall then say that \emph{the sequence of systems $(y_k^{(n)})$ tends to the system $(y_k^*)$}.

Let us prove that one also has
\begin{equation}
 \sum_{k=1}^{\infty}|y_k^*|^p\leq G^p.
 \tag{24}\label{eq:40:24}
\end{equation}
Indeed, for arbitrary $m$,
\[
 \sum_{k=1}^{m}|y_k^*|^p
 =\lim_{n\to\infty}\sum_{k=1}^{m}|y_k^{(n)}|^p
 \leq G^p.
\]
Since the partial sums of the positive-term series~\eqref{eq:40:24} remain at most $G^p$, the series converges and its sum is likewise at most $G^p$.

Now consider the series
\begin{equation}
 \sum_{k=1}^{\infty}a_ky_k;
 \tag{25}\label{eq:40:25}
\end{equation}
% [Source printed page 56]
we have already seen that, when the series $\sum|a_k|^{p/(p-1)}$ and $\sum|y_k|^p$ converge, series~\eqref{eq:40:25} converges absolutely. Moreover, series~\eqref{eq:40:25} converges \emph{uniformly} over the set of systems $(y_k)$ such that
\begin{equation}
 \sum_{k=1}^{\infty}|y_k|^p\leq G^p.
 \tag{26}\label{eq:40:26}
\end{equation}
Indeed, inequality~\eqref{eq:33:4} gives, for every system satisfying~\eqref{eq:40:26},
\begin{equation}
 \left|\sum_{k=m+1}^{\infty}a_ky_k\right|
 \leq G
 \left(\sum_{k=m+1}^{\infty}|a_k|^{\frac{p}{p-1}}\right)^{\frac{p-1}{p}};
 \tag{27}\label{eq:40:27}
\end{equation}
the right-hand side no longer depends on the $y$'s and tends to zero as $1/m$ tends to zero.

In particular, apply inequality~\eqref{eq:40:27} to the systems $(y_k^{(n)})$ and $(y_k^*)$ just considered. We obtain, for arbitrary $m$,
\begin{align*}
 \lim_{n\to\infty}
 \left|\sum_{k=1}^{\infty}a_k(y_k^*-y_k^{(n)})\right|
 &\leq
 \lim_{n\to\infty}\left|\sum_{k=1}^{m}a_k(y_k^*-y_k^{(n)})\right|\\
 &\quad+\lim_{n\to\infty}\left|\sum_{k=m+1}^{\infty}a_ky_k^{(n)}\right|
 +\left|\sum_{k=m+1}^{\infty}a_ky_k^*\right|\\
 &\leq 0+2G
 \left(\sum_{k=m+1}^{\infty}|a_k|^{\frac{p}{p-1}}\right)^{\frac{p-1}{p}}.
\end{align*}
Since this holds for every $m$, the limit on the left is zero\footnotemark{} and
\begin{equation}
 \sum_{k=1}^{\infty}a_ky_k^*
 =\lim_{n\to\infty}\sum_{k=1}^{\infty}a_ky_k^{(n)}.
 \tag{28}\label{eq:40:28}
\end{equation}
\begingroup
\let\OriginalFootnoteText\footnotetext
\long\def\footnotetext#1{%
  \OriginalFootnoteText{\clubpenalty=10000\widowpenalty=10000 #1}%
}
\footnotetext{This result is a special case of the following theorem, which is likewise very easy to prove: \emph{If the sequence $(y_k^{(n)})$ tends to $(y_k^*)$ as $n\to\infty$, and if, moreover, the quantities $a_k^{(n)}$ tend to limits $a_k^*$ in such a way that
\[
 \lim_{n\to\infty}\sum_{k=1}^{\infty}|a_k^*-a_k^{(n)}|^{\frac{p}{p-1}}=0,
\]
then
\[
 \sum_{k=1}^{\infty}a_k^*y_k^*
 =\lim_{n\to\infty}\sum_{k=1}^{\infty}a_k^{(n)}y_k^{(n)}.
\]}}
\endgroup

% [Source printed page 57]
\numberedparagraph{41}
Besides inequality~\eqref{eq:40:24} and equality~\eqref{eq:40:28}, which we shall use below, we shall also need a \emph{selection principle}. For the sequences of systems $(y_k^{(n)})$ which we shall have to consider, it will allow us to infer the existence of subsequences tending to limiting systems.\footnote{For the history of this principle and its role in analysis, see Fr\'echet, \emph{Sur quelques points du Calcul fonctionnel}, thesis, Paris, 1906 (also printed in the \emph{Rendiconti del Circolo Matematico di Palermo}, vol.~XXII, 1906), Chapters IV--VII.}

\begin{selectionprinciple}
Given an infinite sequence of systems $(y_k^{(n)})$ satisfying inequality~\eqref{eq:40:26}, one can extract from it a subsequence of systems
\[
 (y_k^{(n_1)}),(y_k^{(n_2)}),\ldots
\]
which tends to a system $(y_k^*)$.
\end{selectionprinciple}
\begin{proof}
Consider the sequence of quantities $y_1^{(n)}$ $(n=1,2,\ldots)$. Their absolute values remain at most $G$. Hence, by Weierstrass's theorem, the sequence has one or more limiting values. Choose one of them and call it $y_1^*$, and let $n_1',n_2',\ldots$ be a sequence of indices such that\translatornote{The French text says to choose ``the least'' limiting value, although the quantities may be complex and hence are not ordered.  Any limiting value suffices.}
\[
 y_1^{(n_1')},y_1^{(n_2')},\ldots
\]
tends to $y_1^*$. Likewise, the sequence
\[
 y_2^{(n_1')},y_2^{(n_2')},\ldots
\]
has one or more limiting values. Let $y_2^*$ be one of them, and let $n_1'',n_2'',\ldots$ be a sequence of indices selected from the preceding sequence such that
\[
 y_2^{(n_1'')},y_2^{(n_2'')},\ldots
\]
tends to $y_2^*$. Continuing this procedure, we determine values $y_3^*,y_4^*,\ldots$ and corresponding sequences $(n_i'''),(n_i^{\mathrm{iv}}),\ldots$, each selected from its predecessor. Now put
\[
 n_1=n_1',\qquad n_2=n_2'',\qquad n_3=n_3''',\qquad\ldots.
\]
Except always for a finite number of terms, the sequence $(n_i)$ so defined is contained in each of the sequences $(n_i'),(n_i''),\ldots$. Consequently, for every $k$,
\[
 y_k^*=\lim_{i\to\infty}y_k^{(n_i)},
\]
as was to be proved.
\end{proof}

\numberedparagraph{42}
In certain cases our selection principle allows us to conclude that the sequence $(y_k^{(n)})$ itself already tends to a limiting system $(y_k^*)$. This always occurs when the system $(y_k^*)$ is uniquely
% [Source printed page 58]
determined by the problem from which one started. Indeed, suppose that for some $k$, say $k=k_1$, one did not have
\[
 y_{k_1}^*=\lim_{n\to\infty}y_{k_1}^{(n)}.
\]
Since the quantities $y_{k_1}^{(n)}$ are bounded as a whole, they would have a second limiting value
\begin{equation}
 y_{k_1}^{**}\ne y_{k_1}^*,
 \tag{29}\label{eq:42:29}
\end{equation}
and there would be a sequence of indices $m_1,m_2,\ldots$ such that
\[
 y_{k_1}^{**}=\lim_{i\to\infty}y_{k_1}^{(m_i)}.
\]
But our selection principle also applies to the sequence
\[
 (y_k^{(m_1)}),(y_k^{(m_2)}),\ldots;
\]
applied to this sequence, it would yield a limiting system $(y_k^{**})$, which by~\eqref{eq:42:29} would be distinct from $(y_k^*)$, contrary to the hypothesis.

\numberedparagraph{43}
So far, in considering sequences of systems $(y_k^{(n)})$ that tend term by term to a system $(y_k^*)$, we have made only one mildly restrictive hypothesis: that the sums
\begin{equation}
 \sum_{k=1}^{\infty}|y_k^{(n)}|^p
 \qquad(n=1,2,\ldots)
 \tag{30}\label{eq:43:30}
\end{equation}
remain bounded as a whole. Inequality~\eqref{eq:40:24} assured us that the sum
\begin{equation}
 \sum_{k=1}^{\infty}|y_k^*|^p
 \tag{31}\label{eq:43:31}
\end{equation}
could not exceed any limiting value of the sequence~\eqref{eq:43:30}.

Now consider the case in which precisely
\begin{equation}
 \sum_{k=1}^{\infty}|y_k^*|^p
 =\lim_{n\to\infty}\sum_{k=1}^{\infty}|y_k^{(n)}|^p.
 \tag{32}\label{eq:43:32}
\end{equation}
We shall show that in this case one also has
\begin{equation}
 \lim_{n\to\infty}\sum_{k=1}^{\infty}|y_k^*-y_k^{(n)}|^p=0.
 \tag{33}\label{eq:43:33}
\end{equation}

% [Source printed page 59]
Let $\varepsilon$ be an arbitrarily small positive quantity. Since series~\eqref{eq:43:31} converges, $m$ may be chosen so that
\[
 \sum_{k=m+1}^{\infty}|y_k^*|^p
 =\lim_{n\to\infty}\sum_{k=m+1}^{\infty}|y_k^{(n)}|^p
 <\varepsilon^p.
\]
Having chosen $m$ in this way, inequality~\eqref{eq:34:5} gives
\begin{align*}
 \lim_{n\to\infty}\sum_{k=m+1}^{\infty}|y_k^*-y_k^{(n)}|^p
 &\leq\lim_{n\to\infty}
 \left[
  \left(\sum_{k=m+1}^{\infty}|y_k^*|^p\right)^{\frac1p}
  +\left(\sum_{k=m+1}^{\infty}|y_k^{(n)}|^p\right)^{\frac1p}
 \right]^p\\
 &<(2\varepsilon)^p.
\end{align*}
On the other hand,
\[
 \lim_{n\to\infty}\sum_{k=1}^{m}|y_k^*-y_k^{(n)}|^p=0,
\]
since only finitely many terms are involved, each tending to zero. Consequently,
\[
 \lim_{n\to\infty}\sum_{k=1}^{\infty}|y_k^*-y_k^{(n)}|^p
 <(2\varepsilon)^p;
\]
and since the left-hand side does not depend on the value of $\varepsilon$, it must be zero.\footnote{Conversely, if equation~\eqref{eq:43:33} holds, it plainly entails the term-by-term convergence of the $y_k^{(n)}$ to the $y_k^*$; it also entails relation~\eqref{eq:43:32}, as follows immediately from inequality~\eqref{eq:34:7}. Let us add that, for $p=2$, this type of convergence (\emph{strong convergence}) was introduced by Hilbert (loc.~cit., 4th Communication, p.~177); it is also precisely the type used by Schmidt in his investigations. Let us state two further theorems concerning strong convergence; both follow readily from the preceding arguments. The first was established for $p=2$ by Schmidt, and the second,
% [Source printed page 60]
for the same case, by Fr\'echet (\emph{Comptes rendus}, June~24, 1907). Here is the first: \emph{Given the sequence $(y_k^{(n)})$, in order that there exist a system $(y_k^*)$ such that the systems $(y_k^{(n)})$ tend strongly to $(y_k^*)$, it is necessary and sufficient that
\[
 \lim_{m\to\infty,\,n\to\infty}
 \sum_{k=1}^{\infty}|y_k^{(m)}-y_k^{(n)}|^p=0.
\]}
The second theorem constitutes a kind of selection principle for strong convergence: \emph{Given a sequence $(y_k^{(n)})$, in order that one may extract from it a subsequence which converges strongly, it is necessary and sufficient that the series
\[
 \sum_{k=1}^{\infty}|y_k^{(n)}|^p
\]
converge uniformly for all $n$.}}

\section{The Case of Infinitely Many Equations}

\numberedparagraph{44}
Having established the necessary preliminaries, let us now study the case in which system~\eqref{eq:35:8} consists of a countably infinite number of equations. Suppose condition~\eqref{eq:36:14} fulfilled, and let $M^*$ be the least of the numbers $M$ for which~\eqref{eq:36:14} holds. Then the condition is also satisfied by the first $n$ equations of~\eqref{eq:35:8}, and $M^{(n)}\leq M^*$. Let $(x_k^{(n)})$ be the extremal solution (\hyperref[sec:38]{no.~38}) of these $n$ equations. We have
\begin{equation}
 \sum_{k=1}^{\infty}|x_k^{(n)}|^p
 =[M^{(n)}]^p\leq(M^*)^p.
 \tag{34}\label{eq:44:34}
\end{equation}

Thus the sequence of systems $(x_k^{(n)})$ $(n=1,2,\ldots)$ contains a subsequence
\[
 (x_k^{(n_1)}),(x_k^{(n_2)}),\ldots
\]
tending to a system $(x_k^*)$. Applying formula~\eqref{eq:40:28} to this system, we see immediately that the values $x_k^*$ satisfy each of equations~\eqref{eq:35:8}. Moreover, the solution of~\eqref{eq:35:8} so defined is an extremal solution in the sense of \hyperref[sec:38]{no.~38}, for inequality~\eqref{eq:40:24} gives
\[
 \sum_{k=1}^{\infty}|x_k^*|^p\leq(M^*)^p,
\]
while inequality~\eqref{eq:36:13} gives, for every solution $(x_k)$ of~\eqref{eq:35:8},
\begin{equation}
 \sum_{k=1}^{\infty}|x_k|^p\geq(M^*)^p.
 \tag{35}\label{eq:44:35}
\end{equation}

But there is more. System~\eqref{eq:35:8} has only one extremal solution. We have already proved this for the finite system~\eqref{eq:37:15};

% End inlined .parts/pages_001_060.tex
% Begin inlined .parts/pages_061_120.tex
% Translation of source printed pages 61--120.
% [Source printed page 61]
The same reasoning applies in the general case. Now, $(x_k^*)$ is such an
extremal solution; and the same would be true of every other system obtained,
in accordance with our principle of choice, from the sequence of systems
$(x_k^{(n)})$. Consequently, by \hyperref[sec:42]{no.~42},
\begin{equation}
  x_k^*=\lim_{n\to\infty}x_k^{(n)}.
  \tag{36}\label{eq:44:36}
\end{equation}

But the nature of the convergence of the solutions $(x_k^{(n)})$ to the
solution $(x_k^*)$ can be specified much more precisely. It follows from
\eqref{eq:44:34}, \eqref{eq:44:35}, and \eqref{eq:44:36} that
\[
  \sum_{k=1}^{\infty}|x_k^*|^p=M^{*p}
  =\lim_{n\to\infty}\sum_{k=1}^{\infty}|x_k^{(n)}|^p;
\]
hence, by \hyperref[sec:43]{no.~43},
\[
  \lim_{n\to\infty}\sum_{k=1}^{\infty}|x_k^*-x_k^{(n)}|^p=0.
\]

In summary, we may state the following theorem.

\begin{theorem}
For the system of a finite or countably infinite number of equations
\[
  \sum_{k=1}^{\infty}a_{ik}x_k=c_i\qquad(i=1,2,\ldots),
\]
whose coefficients $a_{ik}$ are such that the series
\[
  \sum_{k=1}^{\infty}|a_{ik}|^{p/(p-1)}\qquad(i=1,2,\ldots)
\]
converge, to admit at least one solution $(x_k)$ such that
\[
  \sum_{k=1}^{\infty}|x_k|^p\leq M^p,
\]
it is necessary and sufficient that the inequality
\[
  \left|\sum_{i=1}^{n}\mu_i c_i\right|
  \leq M\left(
    \sum_{k=1}^{\infty}
      \left|\sum_{i=1}^{n}\mu_i a_{ik}\right|^{p/(p-1)}
  \right)^{(p-1)/p}
\]
hold for every $n$ and every choice of the numbers $\mu_i$, real or complex.

% [Source printed page 62]
In particular, if $M$ denotes the least of the numbers for which the
condition is fulfilled, there is only one solution possessing the required
property. In the case of infinitely many equations, this extremal solution
$(x_k^*)$ is related to the extremal solutions $(x_k^{(n)})$ of the partial
systems~\eqref{eq:37:15} by
\[
  \lim_{n\to\infty}\sum_{k=1}^{\infty}|x_k^*-x_k^{(n)}|^p=0.
\]
\end{theorem}

\numberedparagraph{45}
Let us now make a remark which may perhaps suggest to the reader another
proof of the facts just established. Let $(x_k^*)$ be the extremal solution
of system~\eqref{eq:35:8}, and adjoin to that system the equation $x_1=x_1^*$; the
modified system clearly has the same extremal solution. Applying formula
\eqref{eq:36:14} to this new system leads to the inequality
\[
 \left|x_1^*-\sum_{i=1}^{n}\mu_i c_i\right|
 \leq M^*\left(
   \left|1-\sum_{i=1}^{n}\mu_i a_{i1}\right|^{p/(p-1)}
   +\sum_{k=2}^{\infty}\left|\sum_{i=1}^{n}\mu_i a_{ik}\right|^{p/(p-1)}
 \right)^{(p-1)/p}.
\]
This inequality says that the point $x_1^*$ in the complex plane must lie
inside, or on the circumference of, the circle whose centre is
$\sum\mu_i c_i$ and whose radius is given by the right-hand side of the
inequality. Thus $x_1^*$ belongs to every circle corresponding to the
different possible choices of $n$ and the $\mu_i$; and, since the extremal
solution is unique, these circles have no other common point.

The same reasoning applies to each of the points $x_k^*$.

The reader who wishes to base a proof upon the stated property of the
extremal solution should consult a note by Helly,\footnote{E. Helly,
\emph{Ueber lineare Funktionaloperationen}, \emph{Sitzungsberichte d.
k. Akademie d. Wiss., Wien}, vol.~CXXI, IIa, February 1912.} where a system
of integral equations is studied from the same point of view; by analogy it
corresponds to the case $p=1$, which we shall consider later. It should also
be observed that, when the data are real---as is also the case in the note
just cited---the circles may be replaced by segments of the real axis, and
% [Source printed page 63]
the necessary geometric considerations become much simpler.

\section{Homogeneous systems}

\numberedparagraph{46}
One may also ask whether system~\eqref{eq:35:8} possesses, besides $(x_k^*)$, other
admissible solutions, that is, solutions for which $\sum|x_k|^p$ converges.
This question is equivalent to asking whether the homogeneous system
\begin{equation}
  \sum_{k=1}^{\infty}a_{ik}x_k=0\qquad(i=1,2,\ldots)
  \tag{37}\label{eq:46:37}
\end{equation}
possesses admissible solutions other than $x_1=x_2=\cdots=0$. Indeed, from
every solution $(x_k^0)$ of \eqref{eq:46:37} one obtains a solution
$(x_k^{**})=(x_k^*+x_k^0)$ of~\eqref{eq:35:8}; conversely, every solution $(x_k^{**})$
of~\eqref{eq:35:8} gives a solution $(x_k^0)=(x_k^{**}-x_k^*)$; and, by~\eqref{eq:34:5},
$(x_k^0)$ and $(x_k^{**})$ are either both admissible or both inadmissible.

The discussion of the homogeneous system \eqref{eq:46:37} reduces to the
following observation. If the system has a solution $(x_k^0)$ distinct from
the evident solution $x_1=x_2=\cdots=0$, then $x_k^0\neq0$ for one or more
values of $k$; fixing one of these values, say $k=k_1$, we may suppose that
$x_{k_1}^0=1$. Consequently, the system obtained by adjoining to
\eqref{eq:46:37} the equation
\[
  x_{k_1}=1
\]
must satisfy, for some value of $M$, the condition for nonhomogeneous
systems. Thus the existence of an index $k_1$ and a number $M$ for which
that condition is fulfilled is necessary and sufficient for
\eqref{eq:46:37} to possess an admissible solution other than
$x_1=x_2=\cdots=0$.

For arbitrary $p$, the calculation of the solutions of the homogeneous
system has many drawbacks. This is readily seen by comparing the general
case with the case $p=2$. I shall spare the reader this comparison; we shall
return below to the case $p=2$. In any event, one can formulate a kind of
rule which displays all the solutions; this rule is purely theoretical in
the general case, but
% [Source printed page 64]
it lends itself well to effective calculation when $p=2$. Consider system
\eqref{eq:35:8}, keep the $a_{ik}$ fixed, and vary the $c_i$; the number $M^*$ may be
regarded as a function of the $c_i$. Plainly this function is well defined
for all systems $(c_i)$ obtained from~\eqref{eq:35:8} by assigning to the $x_k$ values
for which $\sum|x_k|^p$ converges. Hence, substituting in the function
\[
  M^*(c_1,c_2,\ldots)
\]
the left-hand sides of~\eqref{eq:35:8} in place of the $c_i$, one obtains a function
$F(x_1,x_2,\ldots)$ of infinitely many variables $x_k$; it satisfies
\begin{equation}
  F(x_1,x_2,\ldots)
  \leq\left[\sum_{k=1}^{\infty}|x_k|^p\right]^{1/p}.
  \tag{38}\label{eq:46:38}
\end{equation}
Equality holds when $(x_k)$ is the extremal solution of a system such as
\eqref{eq:35:8}, with suitable $c_i$; strict inequality holds when it is not. The
homogeneous system \eqref{eq:46:37} has no admissible solution other than
$x_1=x_2=\cdots=0$ if and only if every system $(x_k)$ belongs to the first
category, that is, if equality holds without exception in
\eqref{eq:46:38}. Otherwise \eqref{eq:46:37} also has solutions distinct
from $x_1=x_2=\cdots=0$. In every case, the totality of the solutions of
\eqref{eq:46:37} is given by the single equation
\[
  F(x_1,x_2,\ldots)=0.
\]

\section{The Case \texorpdfstring{\(p=2\)}{p = 2}. Schmidt's Theory}

\numberedparagraph{47}
We now pass to the most important special case, $p=2$, and see what can be
deduced for it from our general results.\footnote{Besides the works of
Hilbert and Schmidt already cited, see G. Kowalewski, \emph{Einf\"uhrung in
  die Determinantentheorie}, Leipzig, 1909, pp.~407--455; B\^ocher and Brand,
``On linear equations with an infinite number of variables,'' \emph{Annals
of Mathematics}, second series, vol.~XIII (1912), pp.~167--186.}

% [Source printed page 65]
By~\eqref{eq:38:21}, in the case under consideration the $x_k^{(n)}$ have the form
\[
  x_k^{(n)}
   =\left(\sum_{i=1}^{n}\mu_i c_i\right)
    \left(\sum_{i=1}^{n}\overline{\mu_i}\,\overline{a_{ik}}\right)
   =\sum_{i=1}^{n}\nu_i\overline{a_{ik}}.\footnote{Here and below,
   $\overline a$ denotes the conjugate of the complex number $a$.}
\]
The $\nu_i$ are determined by the equations
\begin{equation}
  \sum_{i=1}^{n}\alpha_{ij}\nu_i=c_j\qquad(j=1,2,\ldots,n),
  \qquad
  \alpha_{ij}=\sum_{k=1}^{\infty}\overline{a_{ik}}a_{jk}.
  \tag{39}\label{eq:47:39}
\end{equation}
These equations are obtained by replacing the $x_k$ in~\eqref{eq:35:8} by the
expressions $\sum\nu_i\overline{a_{ik}}$. Complete system
\eqref{eq:47:39} by adjoining the equation
\[
  \sum_{i=1}^{n}\overline{a_{ij}}\nu_i=x_j^{(n)},
\]
and eliminate the $\nu_i$; this gives
\begin{equation}
 x_j^{(n)}=-
 \frac{
 \begin{vmatrix}
  \alpha_{11}&\cdots&\alpha_{n1}&c_1\\
  \vdots&&\vdots&\vdots\\
  \alpha_{1n}&\cdots&\alpha_{nn}&c_n\\
  \overline{a_{1j}}&\cdots&\overline{a_{nj}}&0
 \end{vmatrix}}
 {\begin{vmatrix}
  \alpha_{11}&\cdots&\alpha_{n1}\\
  \vdots&&\vdots\\
  \alpha_{1n}&\cdots&\alpha_{nn}
 \end{vmatrix}}.
 \tag{40}\label{eq:47:40}
\end{equation}
Moreover,
\[
 M^{(n)2}=\sum_{k=1}^{\infty}|x_k^{(n)}|^2
 =\sum_{k=1}^{\infty}x_k^{(n)}\overline{x_k^{(n)}}
 =\sum_{i=1}^{n}\nu_i
   \left(\sum_{k=1}^{\infty}\overline{a_{ik}}\,\overline{x_k^{(n)}}\right)
 =\sum_{i=1}^{n}\nu_i\overline{c_i}.
\]
Thus the $\nu_i$ satisfy the equation
\[
  \sum_{i=1}^{n}\nu_i\overline{c_i}=M^{(n)2}.
\]

% [Source printed page 66]
Adjoining this equation to system \eqref{eq:47:39} and eliminating the
$\nu_i$, one obtains
\begin{equation}
 M^{(n)2}=-
 \frac{
 \begin{vmatrix}
  \alpha_{11}&\cdots&\alpha_{n1}&c_1\\
  \vdots&&\vdots&\vdots\\
  \alpha_{1n}&\cdots&\alpha_{nn}&c_n\\
  \overline{c_1}&\cdots&\overline{c_n}&0
 \end{vmatrix}}
 {\begin{vmatrix}
  \alpha_{11}&\cdots&\alpha_{n1}\\
  \vdots&&\vdots\\
  \alpha_{1n}&\cdots&\alpha_{nn}
 \end{vmatrix}}.
 \tag{41}\label{eq:47:41}
\end{equation}
Now the Hermitian form\translatornote{The French original omits the conjugation
on \(a_{ik}\) in the sum on the right; it is inserted here so that the displayed
identity agrees with the definition of \(\alpha_{ij}\).}
\[
  \sum_{i,j=1}^{n}\alpha_{ij}\nu_i\overline{\nu_j}
  =\sum_{k=1}^{\infty}\left|\sum_{i=1}^{n}\overline{a_{ik}}\nu_i\right|^2
\]
is manifestly positive; hence the determinant in the denominator is real
and positive. It is understood that the left-hand sides of equations~\eqref{eq:37:15}
are assumed independent, which is equivalent to assuming that
$\det(\alpha_{ik})_{i,k=1}^{n}\neq0$. Otherwise formulas \eqref{eq:47:40} and
\eqref{eq:47:41} have no meaning. The modifications required in that case
belong to the ordinary theory of linear equations, and we shall not dwell
on them.

As in the general case, the extremal solutions $(x_k^{(n)})$ of the partial
systems tend to the extremal solution $(x_k^*)$ of the complete system, and
do so in such a way that
\[
  \lim_{n\to\infty}\sum_{k=1}^{\infty}|x_k^*-x_k^{(n)}|^2=0.
\]

Formula \eqref{eq:47:41} also permits the solvability condition to be stated
in a form no longer containing the indeterminates $\mu_i$. Regard
$M^{(n)}$ as a function of the $c_i$. For each system $(c_i)$, the sequence
$M^{(1)}\leq M^{(2)}\leq M^{(3)}\leq\cdots$ tends to a definite positive
value, finite or infinite, denoted by $M^*$. Thus the formula
\begin{equation}
 M^*(c_1,c_2,\ldots)=\lim_{n\to\infty}
 \left\{-
 \frac{
 \begin{vmatrix}
  \alpha_{11}&\cdots&\alpha_{n1}&c_1\\
  \vdots&&\vdots&\vdots\\
  \alpha_{1n}&\cdots&\alpha_{nn}&c_n\\
  \overline{c_1}&\cdots&\overline{c_n}&0
 \end{vmatrix}}
 {\begin{vmatrix}
  \alpha_{11}&\cdots&\alpha_{n1}\\
  \vdots&&\vdots\\
  \alpha_{1n}&\cdots&\alpha_{nn}
 \end{vmatrix}}\right\}^{1/2}
 \tag{42}\label{eq:47:42}
\end{equation}
% [Source printed page 67]
defines a function of infinitely many variables $c_i$. Accordingly, the
necessary and sufficient condition for system~\eqref{eq:35:8} to possess a solution
$(x_k)$ such that
\begin{equation}
  \sum_{k=1}^{\infty}|x_k|^2\leq M^2
  \tag{43}\label{eq:47:43}
\end{equation}
is the inequality
\[
  M^*(c_1,c_2,\ldots)\leq M.
\]
All these details concern the general situation in which the left-hand
sides are assumed independent. But the function $M^*(c_1,c_2,\ldots)$
also exists in the contrary case; only slight changes to the preceding
method are then required to calculate it.

\numberedparagraph{48}
To carry the study of systems~\eqref{eq:35:8} further, substitute in
\eqref{eq:47:42}, in place of the $c_i$, the left-hand sides of~\eqref{eq:35:8}, and in
place of the $\overline{c_i}$ their conjugates. The result is the expression
denoted in the general case by $F(x_1,x_2,\ldots)$ in
\hyperref[sec:46]{no.~46}. This function $F$ is the square root of the
Hermitian form in infinitely many variables
\begin{equation}
  \sum_{j,k=1}^{\infty}\Lambda_{jk}x_j\overline{x_k},
  \qquad \Lambda_{kj}=\overline{\Lambda_{jk}},
  \tag{44}\label{eq:48:44}
\end{equation}
whose coefficients, when the left-hand sides are independent, are
calculated by
\begin{equation}
 \Lambda_{jk}=-\lim_{n\to\infty}
 \frac{
 \begin{vmatrix}
  \alpha_{11}&\cdots&\alpha_{n1}&a_{1j}\\
  \vdots&&\vdots&\vdots\\
  \alpha_{1n}&\cdots&\alpha_{nn}&a_{nj}\\
  \overline{a_{1k}}&\cdots&\overline{a_{nk}}&0
 \end{vmatrix}}
 {\begin{vmatrix}
  \alpha_{11}&\cdots&\alpha_{n1}\\
  \vdots&&\vdots\\
  \alpha_{1n}&\cdots&\alpha_{nn}
 \end{vmatrix}}.
 \tag{45}\label{eq:48:45}
\end{equation}
For every system of values $(x_k)$, the Hermitian form
\eqref{eq:48:44} remains at most $\sum|x_k|^2$. Equality holds only when
$(x_k)$ is the extremal solution of a system such as~\eqref{eq:35:8}. The solutions of
the homogeneous
% [Source printed page 68]
system
\[
  \sum_{k=1}^{\infty}a_{ik}x_k=0\qquad(i=1,2,\ldots)
\]
are characterized by the fact that the Hermitian form \eqref{eq:48:44}
vanishes for these systems $(x_k)$ and not for the others. The condition
that the homogeneous system have no solution except the evident one is
\[
  \Lambda_{jk}=\begin{cases}1,&j=k,\\0,&j\neq k.\end{cases}
\]

Here is another simple interpretation of the coefficients $\Lambda_{jk}$.
For a fixed index $j$, consider the system
\begin{equation}
  \sum_{k=1}^{\infty}a_{ik}x_k=a_{ij}\qquad(i=1,2,\ldots).
  \tag{46}\label{eq:48:46}
\end{equation}
This system plainly has the solution $x_k=1$ for $k=j$ and $x_k=0$ for
$k\neq j$. On the other hand, to calculate its extremal solution one need
only replace the $c_i$ in formula \eqref{eq:47:40} by the $a_{ij}$ and let
$n$ increase indefinitely; comparison with \eqref{eq:48:45} gives
\[
  x_k=\Lambda_{jk}\qquad(k=1,2,\ldots).
\]
This is the desired extremal solution of \eqref{eq:48:46}. The difference
between the two solutions of \eqref{eq:48:46}, namely
$\Lambda_{j1},\Lambda_{j2},\ldots,\Lambda_{jj}-1,\ldots$, therefore satisfies
the homogeneous system. Thus each value of $j$ gives one solution of the
homogeneous system.

More generally, let $(\xi_k)$ be any system of values for which
$\sum|\xi_k|^2$ converges. The system
\begin{equation}
  \sum_{k=1}^{\infty}a_{ik}x_k
  =\sum_{k=1}^{\infty}a_{ik}\xi_k\qquad(i=1,2,\ldots)
  \tag{47}\label{eq:48:47}
\end{equation}
plainly has the solution $x_k=\xi_k$. Let us now evaluate its extremal
solution. The first $n$ equations of \eqref{eq:48:47} have the extremal
solution
\[
  x_k^{(n)}=\sum_{j=1}^{\infty}\Lambda_{jk}^{(n)}\xi_j
  \qquad(k=1,2,\ldots),
\]
where $\Lambda_{jk}^{(n)}$ denotes the $n$th approximate value of
$\Lambda_{jk}$ occurring
% [Source printed page 69]
in \eqref{eq:48:45}. But
$x_1=\Lambda_{j1}^{(n)}=\overline{\Lambda_{1j}^{(n)}}$,
$x_2=\Lambda_{j2}^{(n)}=\overline{\Lambda_{2j}^{(n)}},\ldots$ is the
extremal solution of the first $n$ equations of \eqref{eq:48:46}; those
equations also have the solution $x_k=1$ for $k=j$ and $x_k=0$ for
$k\neq j$. Hence
\[
  \sum_{j=1}^{\infty}|\Lambda_{jk}^{(n)}|^2\leq1.
\]
Consequently formula~\eqref{eq:40:28} applies: there one need only replace the $a_k$
by the $\xi_j$, the $y_k^*$ by the $\Lambda_{jk}$, and the $y_k^{(n)}$ by
the $\Lambda_{jk}^{(n)}$. The desired extremal solution is therefore
\[
  x_k^*=\lim_{n\to\infty}x_k^{(n)}
  =\lim_{n\to\infty}\sum_{j=1}^{\infty}\Lambda_{jk}^{(n)}\xi_j
  =\sum_{j=1}^{\infty}\Lambda_{jk}\xi_j.
\]
The difference between the two solutions just found for
\eqref{eq:48:47} satisfies the homogeneous system, and in fact gives its
general solution. For if $(x_k^0)$ is any solution of the homogeneous
system, choose $\xi_k=-x_k^0$; then \eqref{eq:48:47} becomes homogeneous,
has $x_1=x_2=\cdots=0$ as its extremal solution, and the preceding procedure
gives $(x_k^0)$.

Thus the formulas\translatornote{The printed formula has \(+\xi_k\), but the
source copy carries a handwritten correction to \(-\xi_k\), as required by the
preceding choice \(\xi_k=-x_k^0\).}
\[
  x_k=-\xi_k+\sum_{j=1}^{\infty}\Lambda_{jk}\xi_j
  \qquad(k=1,2,\ldots)
\]
include all solutions of the homogeneous system. Equivalently, every
solution of the homogeneous system is a finite or infinite linear
combination of the particular solutions
$\Lambda_{j1},\ldots,\Lambda_{jj}-1,\ldots$ $(j=1,2,\ldots)$, the
multipliers $\xi_j$ being subject to the condition that
$\sum|\xi_j|^2$ converge.

\numberedparagraph{49}
All these calculations simplify greatly in the special case in which
\[
  \alpha_{ij}=\sum_{k=1}^{\infty}a_{ik}\overline{a_{jk}}
  =\begin{cases}1,&i=j,\\0,&i\neq j,\end{cases}
\]
which Hilbert and
% [Source printed page 70]
Schmidt express by saying that the linear forms
\[
  \sum_{k=1}^{\infty}a_{ik}x_k\qquad(i=1,2,\ldots)
\]
constitute an orthogonal and normalized system.

Formulas \eqref{eq:47:40} and \eqref{eq:47:41} give in this case
\[
  x_k^{(n)}=\sum_{i=1}^{n}\overline{a_{ik}}c_i,
  \qquad M^{(n)2}=\sum_{i=1}^{n}|c_i|^2.
\]
Consequently, the condition that the system possess a solution $(x_k)$
satisfying \eqref{eq:47:43} is
\[
  M^{*2}=\sum_{i=1}^{\infty}|c_i|^2\leq M^2,
\]
and the extremal solution $(x_k^*)$ is
\[
  x_k^*=\sum_{i=1}^{\infty}\overline{a_{ik}}c_i
  \qquad(k=1,2,\ldots).
\]
For this solution,
\[
  \sum_{k=1}^{\infty}|x_k^*|^2=\sum_{i=1}^{\infty}|c_i|^2.
\]
The coefficients of the Hermitian form \eqref{eq:48:44} are
\begin{equation}
  \Lambda_{jk}=\sum_{n=1}^{\infty}a_{nj}\overline{a_{nk}}.
  \tag{48}\label{eq:49:48}
\end{equation}
If, under the present hypothesis, the homogeneous system has no solution
other than $x_1=x_2=\cdots=0$, we say that the system of left-hand sides is
\emph{complete}; for this is equivalent to the nonexistence of a linear
form orthogonal to all the left-hand sides at once. As we have seen, this
holds if and only if $\Lambda_{jk}=1$ for $j=k$ and $0$ for $j\neq k$.
In view of \eqref{eq:49:48}, this condition may be stated by saying that
% [Source printed page 71]
the transposed forms
\[
  \sum_{i=1}^{\infty}a_{ik}x_i\qquad(k=1,2,\ldots)
\]
also constitute an orthogonal and normalized system. This formulation is
noteworthy because of its evident analogy with the familiar case of an
orthogonal and normalized system of $n$ linear forms in $n$ variables. In
that case completeness follows immediately because there are as many forms
as variables; with infinitely many variables there is no analogous reason.

All these results concerning orthogonal systems can also be established
directly and very simply. We shall not give the direct proof, which would
merely repeat rapidly, with obvious changes, the arguments used in the
general case.

\numberedparagraph{50}
The special case just considered was first discussed by Hilbert, who was
led to it by ideas to which we shall return later. Schmidt studied this
special case in greater detail and, moreover, succeeded in reducing the
general case $p=2$ to it.\footnote{See the memoir cited in
\hyperref[sec:32]{no.~32}.} In his method the special case under
consideration plays a very important role: he reduces the general case to
it by a very simple process, \emph{orthogonalization}. This process consists
in replacing system~\eqref{eq:35:8} by an equivalent system
\begin{equation}
  \sum_{k=1}^{\infty}b_{ik}x_k=d_i\qquad(i=1,2,\ldots),
  \tag{49}\label{eq:50:49}
\end{equation}
whose left-hand sides are, on the one hand, linear combinations of those
of equations~\eqref{eq:35:8}, and, on the other hand, form an orthogonal and normalized
system. More precisely, the $n$th equation in \eqref{eq:50:49} is a linear
combination of the first $n$ equations in~\eqref{eq:35:8}. For such a reduction to be
possible, equations~\eqref{eq:35:8} (or rather every finite number of them) must
% [Source printed page 72]
be linearly independent; but we have seen that this restriction is only
apparent. The multipliers, and hence the quantities $b_{ik}$ and $d_i$, are
easily found. To simplify the calculation we first form an intermediate
system
\begin{equation}
  \sum_{k=1}^{\infty}b'_{ik}x_k=d'_i\qquad(i=1,2,\ldots),
  \tag{50}\label{eq:50:50}
\end{equation}
requiring only that its left-hand sides be orthogonal; to pass from it to
the normalized system \eqref{eq:50:49}, each equation is divided by
\[
  \left[\sum_{k=1}^{\infty}|b'_{ik}|^2\right]^{1/2}.
\]
The data of such a system \eqref{eq:50:50} are obtained from the linear
equations expressing orthogonality by the formulas
\[
 b'_{nk}=
 \frac{
 \begin{vmatrix}
  \alpha_{11}&\cdots&\alpha_{1n}&0\\
  \vdots&&\vdots&\vdots\\
  \alpha_{n1}&\cdots&\alpha_{nn}&1\\
  a_{1k}&\cdots&a_{nk}&0
 \end{vmatrix}}
 {\begin{vmatrix}
  \alpha_{11}&\cdots&\alpha_{1n}\\
  \vdots&&\vdots\\
  \alpha_{n1}&\cdots&\alpha_{nn}
 \end{vmatrix}},
 \qquad
 d'_n=
 \frac{
 \begin{vmatrix}
  \alpha_{11}&\cdots&\alpha_{1n}&0\\
  \vdots&&\vdots&\vdots\\
  \alpha_{n1}&\cdots&\alpha_{nn}&1\\
  c_1&\cdots&c_n&0
 \end{vmatrix}}
 {\begin{vmatrix}
  \alpha_{11}&\cdots&\alpha_{1n}\\
  \vdots&&\vdots\\
  \alpha_{n1}&\cdots&\alpha_{nn}
 \end{vmatrix}}.
\]
The left-hand sides of equations \eqref{eq:50:49} may also be calculated
from
\[
  \left|\sum_{k=1}^{\infty}b_{nk}x_k\right|^2
  =\sum_{j,k=1}^{\infty}
   \bigl(\Lambda_{jk}^{(n)}-\Lambda_{jk}^{(n-1)}\bigr)x_j\overline{x_k}
  \qquad(n=1,2,\ldots).
\]
This is simply the evident fact that two equivalent systems always lead to
the same Hermitian form \eqref{eq:48:44}.

Having transformed system~\eqref{eq:35:8} in this manner into the special system
\eqref{eq:50:49}, its discussion and solution proceed, as we have seen, by
very simple formulas. Substitution in those formulas of the expressions for
$b$ and $d$ in terms of $a$ and $c$ recovers the formulas established above
by another method.

But this is only the principal idea of Schmidt's work
% [Source printed page 73]
which we have just described; for the details we must refer the reader to
the original memoir. One of its most interesting features is that it sees
and speaks as though the subject were geometry. To fix ideas, suppose that
all quantities are real. Interpret systems $(a_k)$ or $(x_k)$, whenever
$\sum a_k^2$ or $\sum x_k^2$ converges, as vectors in a space of infinitely
many dimensions. Given finitely or infinitely many vectors, the vectors
orthogonal to each of them constitute a certain linear manifold, and the
set of vectors orthogonal to this latter manifold is the smallest linear
manifold containing the given vectors. Thus, to solve a homogeneous system,
one need only regard its coefficients as the coordinates of vectors; the
manifold orthogonal to these vectors represents the totality of the
solutions. Nonhomogeneous systems may be made homogeneous by introducing an
auxiliary unknown. Finally, orthogonalization amounts in principle to
replacing the given system of vectors by a system of mutually orthogonal
vectors in such a way that the two systems determine the same linear
manifold.

\section{The Cases \texorpdfstring{\(p=1\)}{p = 1} and
\texorpdfstring{\(p\to\infty\)}{p tending to infinity}}

\numberedparagraph{51}
Consider first the case $p=1$. It was not included in the preceding
arguments, which applied only to $p>1$, but it enters as a limiting case.

Here one seeks a solution $(x_k)$ of system~\eqref{eq:35:8} for which the series
\begin{equation}
  \sum_{k=1}^{\infty}|x_k|
  \tag{51}\label{eq:51:51}
\end{equation}
converges. As regards the data, we assume that, for every $i$,
\[
  \lim_{k\to\infty}a_{ik}=0.
\]
The reason is as follows. If the series
\begin{equation}
  \sum_{k=1}^{\infty}a_kx_k
  \tag{52}\label{eq:51:52}
\end{equation}
% [Source printed page 74]
were merely required to converge for every system $(x_k)$ for which
\eqref{eq:51:51} exists, it would suffice to suppose the quantities
$|a_k|$ bounded as a whole. But, just as above, we also need this series to
converge uniformly for all systems $(x_k)$ such that
\[
  \sum_{k=1}^{\infty}|x_k|\leq G.
\]
In particular, convergence must be uniform for all systems in which one
$x_k$ equals $G$ and all the others are zero. Substitute these values in
turn in \eqref{eq:51:52}; the $n$th remainders are respectively
\[
  0,\ldots,0,Ga_{n+1},Ga_{n+2},\ldots.
\]
Uniform convergence holds if and only if, as $n\to\infty$, these remainders
tend uniformly to zero. This is equivalent to assuming, as above, that
\[
  a_k\to0.
\]
With this restriction, all the auxiliary arguments used in the general
case extend, \emph{mutatis mutandis}, to $p=1$.

\numberedparagraph{52}
\begin{theorem}
A necessary and sufficient condition for system~\eqref{eq:35:8} to possess a solution
$(x_k)$ such that
\begin{equation}
  \sum_{k=1}^{\infty}|x_k|\leq M
  \tag{53}\label{eq:52:53}
\end{equation}
is that
\begin{equation}
  \left|\sum_{i=1}^{n}\mu_i c_i\right|
  \leq M\operatorname*{sup}_{k=1,2,\ldots}
    \left|\sum_{i=1}^{n}\mu_i a_{ik}\right|
  \tag{54}\label{eq:52:54}
\end{equation}
for every $n$ and every choice of the $\mu_i$.
\end{theorem}

The condition is plainly necessary. To prove it sufficient, it is enough
to consider $n$ equations, since passage to infinitely many equations is
made exactly as for $p>1$.

% [Source printed page 75]
From the hypothesis $a_{ik}\to0$ it follows at once that there is a number
$m$ such that the value of
\[
  \operatorname*{sup}_{k=1,2,\ldots}
    \left|\sum_{i=1}^{n}\mu_i a_{ik}\right|
\]
depends only on indices $k\leq m$. This permits us to replace the system
with infinitely many unknowns by one with $m$ unknowns,
\begin{equation}
  \sum_{k=1}^{m}a_{ik}x_k=c_i\qquad(i=1,2,\ldots,n).
  \tag{55}\label{eq:52:55}
\end{equation}
Indeed, if \eqref{eq:52:55} has a solution $(x_k)$ such that
\[
  \sum_{k=1}^{m}|x_k|\leq M,
\]
then, putting $x_{m+1}=x_{m+2}=\cdots=0$, we obtain a solution of the
original system satisfying \eqref{eq:52:53}. By the stated property of
$m$, condition \eqref{eq:52:54} holds for the reduced system
\eqref{eq:52:55}.\translatornote{The French text refers here to system~(56),
evidently a misprint for the reduced system~(55).} On the other hand, for every $p>1$,
\[
 \operatorname*{sup}_{k=1,\ldots,m}
 \left|\sum_{i=1}^{n}\mu_i a_{ik}\right|
 \leq
 \left[
   \sum_{k=1}^{m}
   \left|\sum_{i=1}^{n}\mu_i a_{ik}\right|^{p/(p-1)}
 \right]^{(p-1)/p}.
\]
Consequently, for every $p>1$ the data of \eqref{eq:52:55} satisfy
\[
 \left|\sum_{i=1}^{n}\mu_i c_i\right|
 \leq M\left[
   \sum_{k=1}^{m}
   \left|\sum_{i=1}^{n}\mu_i a_{ik}\right|^{p/(p-1)}
 \right]^{(p-1)/p}.
\]
But this is precisely the condition that \eqref{eq:52:55} possess a
solution $(x_k^{(p)})$ such that
\begin{equation}
  \sum_{k=1}^{m}|x_k^{(p)}|^p\leq M^p.
  \tag{56}\label{eq:52:56}
\end{equation}
Thus, for each $p>1$, there exists a solution
% [Source printed page 76]
$(x_k^{(p)})$ satisfying \eqref{eq:52:56}; for definiteness one may take,
for example, the extremal solutions.

The quantities $|x_k^{(p)}|\leq M$ are bounded as a whole. Hence there is a
sequence of exponents $p_1,p_2,\ldots$ tending to $1$ such that the limits
\[
  x_k^*=\lim_{p_\nu\to1}x_k^{(p_\nu)}\qquad(k=1,\ldots,m)
\]
exist. Since every system $(x_k^{(p_\nu)})$ satisfies
\eqref{eq:52:55}, the limiting system $(x_k^*)$ does also. Moreover,
\[
  \sum_{k=1}^{m}|x_k^*|
  =\lim_{p_\nu\to1}\sum_{k=1}^{m}|x_k^{(p_\nu)}|^{p_\nu}
  \leq\lim_{p\to1}M^p=M.
\]
The sufficiency of the condition is therefore proved.

\numberedparagraph{53}
A second limiting case corresponds to $p\to\infty$. Much less is required
of the solution $(x_k)$ than in the cases already treated: it is required
only that the values $(x_k)$ be bounded as a whole. As for the data, the
only assumption is that the left-hand sides of equations~\eqref{eq:35:8} converge for
every bounded system $(x_k)$. This amounts to supposing that all the series
$\sum_k|a_{ik}|$ $(i=1,2,\ldots)$ converge.

\begin{theorem}
The necessary and sufficient condition for system~\eqref{eq:35:8} to possess a
solution $(x_k)$ such that $|x_k|\leq M$ for every $k$ is that
\begin{equation}
  \left|\sum_{i=1}^{n}\mu_i c_i\right|
  \leq M\sum_{k=1}^{\infty}
    \left|\sum_{i=1}^{n}\mu_i a_{ik}\right|
  \tag{57}\label{eq:53:57}
\end{equation}
for every $n$ and every choice of the $\mu_i$.
\end{theorem}

The passage from finitely many equations to infinitely many being almost
the same as in the other cases, it is enough to study $n$ equations.

The condition is plainly necessary. To prove sufficiency, suppose that it
holds and also suppose (as we know this entails no loss of generality) that
there is no linear relation among the left-hand sides of the equations.
% [Source printed page 77]
Then the ratio
\[
 \frac{
  \displaystyle\sum_{k=1}^{\infty}
    \left|\sum_{i=1}^{n}\mu_i a_{ik}\right|}
 {\displaystyle\left[
  \sum_{k=1}^{\infty}
    \left|\sum_{i=1}^{n}\mu_i a_{ik}\right|^{p/(p-1)}
 \right]^{(p-1)/p}}
\]
tends to $1$ as $p\to\infty$, uniformly for every choice of the $\mu_i$.
Inequality \eqref{eq:53:57} may therefore be replaced by
\[
 \left|\sum_{i=1}^{n}\mu_i c_i\right|
 \leq (M+\varepsilon)\left[
  \sum_{k=1}^{\infty}
    \left|\sum_{i=1}^{n}\mu_i a_{ik}\right|^{p/(p-1)}
 \right]^{(p-1)/p},
\]
where $\varepsilon$ is an arbitrarily small positive quantity and $p$ is
sufficiently large. This modified inequality is precisely the condition
that our system of equations possess a solution $(x_k)$ such that
\[
  \sum_{k=1}^{\infty}|x_k|^p\leq(M+\varepsilon)^p.
\]
Let $(x_k^{\varepsilon})$ be such a solution; plainly
$|x_k^{\varepsilon}|\leq M+\varepsilon$ for every $k$. Letting
$\varepsilon$ tend to zero yields an infinite sequence of solutions
$(x_k^{(\varepsilon)})$. Then, by applying a procedure analogous to that of
\hyperref[sec:41]{no.~41}, one selects a subsequence such that, for every
$k$, $x_k^{(\varepsilon)}$ tends to a limit $x_k^*$. The system $(x_k^*)$
is then the required solution.

It remains to observe that, whereas in the general case $p>1$ the extremal
solution was always unique, this is no longer so in the two limiting cases
just considered. The reason is readily seen. For example, in the case
$p=1$, it is the fact that in
$\sum|a_k+b_k|\leq\sum|a_k|+\sum|b_k|$, equality may hold without all the
ratios $a_k:b_k$ being equal.

% [Source printed page 78]
\chapter{Theory of Linear Substitutions in Infinitely Many Variables}\label{chap:4}

\section{Linear Substitutions}\label{sec:linear-substitutions}

\numberedparagraph{54}
The criteria established in the \hyperref[chap:3]{preceding chapter} relate to very broadly
conceived systems of equations. We shall now restrict ourselves to the case
in which $p=p/(p-1)=2$. Most of the arguments, however, extend without
essential alteration to every $p>1$, and even to the two limiting cases.

We shall call the aggregate of all systems $(x_k)$ for which
$\sum_k|x_k|^2$ converges the \emph{Hilbert space}. A correspondence which
assigns to every element $(x_k)$ of this space one and only one element
$(x'_k)$ of the same space will be called a \emph{linear substitution} if it
has the following properties. First, it is distributive: the elements
corresponding to $(c x_k)$ and $(x_k+y_k)$ are respectively $(c x'_k)$ and
$(x'_k+y'_k)$. In algebra, where only finitely many variables occur,
distributivity also entails continuity. Here continuity must be assumed
explicitly. Secondly, it is continuous: whenever a sequence of elements
tends to a limiting element, the corresponding sequence of transformed
elements tends to the element corresponding to that limit. Such
distributive and continuous correspondences will henceforth be called
linear substitutions.

% [Source printed page 79]
\numberedparagraph{55}
Let us recall that the convergence of $(x_k^{(n)})$ to $(x_k)$ means that
$x_k^{(n)}\to x_k$ for every fixed $k$ and that the sums
$\sum_k|x_k^{(n)}|^2$ remain bounded uniformly in $n$
(\hyperref[sec:40]{no.~40}). Continuity then implies the existence of a
finite constant $M$ such that
\begin{equation}
  \sum_{k=1}^{\infty}|x'_k|^2
  \leq M^2\sum_{k=1}^{\infty}|x_k|^2.
  \tag{1}\label{eq:55:1}
\end{equation}
Indeed, otherwise there would be a sequence of elements $(x_k^{(n)})$ and the
corresponding sequence $(x_k^{(n)\prime})$ such that
$\sum_k|x_k^{(n)}|^2$ remained bounded, while
$\sum_k|x_k^{(n)\prime}|^2$ increased beyond every bound as $n\to\infty$.
By the selection principle of \hyperref[sec:41]{no.~41}, one could extract
from the former a subsequence converging to an element; the corresponding
sequence, lacking property~\(2^\circ\), could converge to no element.
Continuity would therefore fail.

The least admissible value of $M$ will be called the \emph{bound} of the
substitution $A$ and will be denoted by $M_A$.

\numberedparagraph{56}
We shall say that $(x_k^{(n)})$ tends \emph{strongly} to $(x_k)$ if
\begin{equation}
  \sum_{k=1}^{\infty}|x_k-x_k^{(n)}|^2\longrightarrow0.
  \tag{2}\label{eq:56:2}
\end{equation}
Since a linear substitution is distributive, the element corresponding to
$(x_k-x_k^{(n)})$ is $(x'_k-x_k^{(n)\prime})$; hence
% [Source printed page 80]
\[
  \sum_{k=1}^{\infty}|x'_k-x_k^{(n)\prime}|^2
  \leq M^2\sum_{k=1}^{\infty}|x_k-x_k^{(n)}|^2\longrightarrow0.
\]
Thus a linear substitution also preserves strong convergence.

The reduced elements
\[
  (x_1,x_2,\ldots,x_n,0,0,\ldots)
\]
tend strongly to $(x_k)$. If $(a_{ik})_{i=1}^{\infty}$ denotes the element
corresponding to the element whose $k$th coordinate is $1$ and whose other
coordinates vanish,\translatornote{The French text says that the $j$th
coordinate is $1$ when $j=i$; the column $(a_{ik})_i$ and formula~(4) show
that $j=k$ is intended.} distributivity and the preceding observation give
\begin{equation}
  \sum_{i=1}^{\infty}
  \left|x'_i-\sum_{k=1}^{n}a_{ik}x_k\right|^2\longrightarrow0,
  \tag{3}\label{eq:56:3}
\end{equation}
and consequently
\begin{equation}
  x'_i=\sum_{k=1}^{\infty}a_{ik}x_k
  \qquad(i=1,2,\ldots).
  \tag{4}\label{eq:56:4}
\end{equation}
Every linear substitution therefore corresponds to a matrix $(a_{ik})$
for which the series in \eqref{eq:56:4} converge and
\begin{equation}
  \sum_{i=1}^{\infty}
    \left|\sum_{k=1}^{\infty}a_{ik}x_k\right|^2
  \leq M^2\sum_{k=1}^{\infty}|x_k|^2.
  \tag{5}\label{eq:56:5}
\end{equation}
Conversely, a matrix satisfying \eqref{eq:56:5} for every element of
Hilbert space defines by \eqref{eq:56:4} a bounded distributive, and hence
continuous, correspondence; the continuity of the individual linear forms
follows from the results of the \hyperref[chap:3]{preceding chapter}.

% [Source printed page 81]
Inequality \eqref{eq:56:5}, required for every element $(x_k)$, is plainly
equivalent to
\[
  \sum_{i=1}^{m}\left|\sum_{k=1}^{n}a_{ik}x_k\right|^2
  \leq M^2\sum_{k=1}^{n}|x_k|^2,
\]
required for every choice of the integers $m,n$ and of the quantities
$x_1,\ldots,x_n$. We thus have the following rule.

\begin{theorem}
A matrix $(a_{ik})$ defines a linear substitution $A$ with $M_A\leq M$ if
and only if
\[
  \sum_{i=1}^{m}\left|\sum_{k=1}^{n}a_{ik}x_k\right|^2
  \leq M^2\sum_{k=1}^{n}|x_k|^2
\]
for every choice of the integers $m,n$ and of the quantities
$x_1,\ldots,x_n$.
\end{theorem}

\numberedparagraph{57}
We should mention here a very remarkable fact discovered by Hellinger and
Toeplitz. Instead of assuming inequality \eqref{eq:56:5}, it is enough to
suppose that its left-hand side converges for every $(x_k)$. The existence
of a finite constant $M_A$ follows from this last hypothesis by a very
delicate argument. The authors stated their theorem in a slightly
different form, connecting it with the notion of a bilinear form which we
shall meet in the next section.\footnote{Hellinger and Toeplitz,
\emph{Grundlagen f\"ur eine Theorie der unendlichen Matrizen},
\emph{Mathematische Annalen}, vol.~69, pp.~289--330. They consider the
expression
\[
  \sum_{i=1}^{\infty}\left(\sum_{k=1}^{\infty}a_{ik}x_k\right)y_i
\]
and suppose that it converges for all elements $(x_k),(y_k)$ of Hilbert
space. By a much simpler theorem of the same authors, discussed in the
\hyperref[chap:3]{preceding chapter} (\hyperref[sec:35]{no.~35}), convergence of $\sum_i b_i y_i$ for every
element $(y_i)$ entails convergence of $\sum_i|b_i|^2$. This reduces their
hypothesis to the one in the text.

Here, in a few words, is the idea of the proof. Contrary to what is to be
proved, suppose that
\[
 \frac{\displaystyle\sum_{i=1}^{\infty}
       \left|\sum_{k=1}^{\infty}a_{ik}x_k\right|^2}
      {\displaystyle\sum_{k=1}^{\infty}|x_k|^2}
\]
can exceed every finite bound. The analogous expressions in which the
summations begin only with $i=m$, $k=n$ have the same property. On the
other hand, for a fixed element $(x_k)$, restricting the summations to
sufficiently large finite indices $m',n'$ does not appreciably alter the
value. These facts permit one to choose successively indices
$i_1,k_1,i_2,k_2,\ldots$ and quantities
$x_1,\ldots,x_{k_1},\ldots,x_{k_2},\ldots$ so that
$\sum_k|x_k|^2$ converges, while
\[
 \sum_{i=i_\nu+1}^{i_{\nu+1}}
   \left|\sum_{k=k_\nu+1}^{k_{\nu+1}}a_{ik}x_k\right|^2
\]
increases rapidly with $\nu$, and at the same time
\[
 \sum_{i=i_\nu+1}^{i_{\nu+1}}
   \left|\sum_{k=1}^{\infty}a_{ik}x_k\right|^2
\]
depends appreciably only upon the terms whose indices lie between
$k_\nu+1$ and $k_{\nu+1}$. For this particular element the left-hand side
of \eqref{eq:56:5} therefore diverges.

This reasoning is closely related to a still little-known argument of
Lebesgue, who connects the existence of continuous $2\pi$-periodic
functions whose Fourier series diverges or fails to converge uniformly
with the fact that the supremum of the absolute value of the $n$th partial
sums, over continuous $2\pi$-periodic functions satisfying $|f|\leq1$,
increases without bound with $n$; see Lebesgue, \emph{Sur la divergence et
convergence non uniforme des s\'eries de Fourier}, \emph{Comptes rendus},
1905, second semester, pp.~875--877, and \emph{Le\c{c}ons sur les s\'eries
trigonom\'etriques}, p.~86. Fej\'er has recently carried Lebesgue's idea
further, and Haar and Lebesgue himself extended it to analogous problems;
see Haar, \emph{Zur Theorie der orthogonalen Funktionensysteme}, G\"ottingen
thesis, 1909 (reprinted in \emph{Mathematische Annalen}, vol.~69,
pp.~331--371); Lebesgue, \emph{Sur les int\'egrales singuli\`eres},
\emph{Annales de la Facult\'e des Sciences de Toulouse}, 3rd series,
vol.~1 (1910), pp.~25--117; and Fej\'er, \emph{Sur les singularit\'es de la
s\'erie de Fourier des fonctions continues}, \emph{Annales de l'\'Ecole
Normale}, 3rd series, vol.~28 (1911), pp.~63--104, where earlier work is
also listed. These investigations do not belong to the subject of this
book; they are mentioned in the hope that their study may suggest applying
the Hellinger--Toeplitz method to many questions of the present theory and,
in particular, extending their theorem to every case $p\geq1$.}

% [Source printed page 82]
Let us add that, in the definition of linear substitutions, the first
definition of continuity used above could be replaced by one based upon
strong convergence. The correspondence would be called continuous when,
whenever $(x_k^{(n)})$ tends strongly to $(x_k)$ in the sense of
\eqref{eq:56:2}, the corresponding $(x_k^{(n)\prime})$ tends strongly to
$(x'_k)$. Only the definition would change: the two definitions being
equivalent, the notion itself remains completely unaltered.

\section{Bilinear Forms and Transposed Substitutions}
\label{sec:bilinear-transpose}

\numberedparagraph{58}
The theory of matrices $(a_{ik})$ of the kind considered here is due to
Hilbert. His formal point of departure differs from ours. He considers the
infinite bilinear form
% [Source printed page 83]
\begin{equation}
  A(x,y)=\sum_{i,k=1}^{\infty}a_{ik}x_k y_i,
  \tag{6}\label{eq:58:6}
\end{equation}
and assumes that the absolute values of the reductions
\[
  [A(x,y)]_n=\sum_{i,k=1}^{n}a_{ik}x_k y_i,
\]
formed for every $n$ and all elements $(x_k),(y_k)$ satisfying
\begin{equation}
  \sum_{k=1}^{\infty}|x_k|^2\leq1,
  \qquad
  \sum_{k=1}^{\infty}|y_k|^2\leq1,
  \tag{7}\label{eq:58:7}
\end{equation}
remain below a positive constant $M$. The convergence of the series
\eqref{eq:56:4} and inequality \eqref{eq:56:5} follow at once; hence
Hilbert's matrices define linear substitutions. Conversely, if our point
of departure is adopted, every linear substitution leads, by the analysis
of \hyperref[sec:56]{no.~56}, to a uniquely determined matrix $(a_{ik})$;
inequality \eqref{eq:56:5}, together with the Cauchy--Lagrange inequality,
places us under Hilbert's hypothesis.
% [Source printed page 84]
These two inequalities assure at the same time the convergence of
\[
  \sum_{i=1}^{\infty}
    \left(\sum_{k=1}^{\infty}a_{ik}x_k\right)y_i
\]
and show that its absolute value is at most $M$. Formula
\eqref{eq:58:6} and the Cauchy--Lagrange inequality give
\begin{equation}
\begin{split}
  \sum_{i=1}^{\infty}
    \left(\sum_{k=1}^{\infty}a_{ik}x_k\right)y_i
  &=\lim_{n\to\infty}\sum_{k=1}^{n}
    \left(\sum_{i=1}^{\infty}a_{ik}y_i\right)x_k\\
  &=\sum_{k=1}^{\infty}
    \left(\sum_{i=1}^{\infty}a_{ik}y_i\right)x_k.
\end{split}
  \tag{8}\label{eq:58:8}
\end{equation}
More generally,
\begin{align*}
&\left|
 \sum_{i=1}^{m+p}\sum_{k=1}^{n+q}a_{ik}x_k y_i
 -\sum_{i=1}^{m}\sum_{k=1}^{n}a_{ik}x_k y_i
 \right|\\
&\quad\leq
 \left|\sum_{i=1}^{m}\sum_{k=n+1}^{n+q}a_{ik}x_k y_i\right|
 +\left|\sum_{i=m+1}^{m+p}\sum_{k=1}^{n+q}a_{ik}x_k y_i\right|\\
&\quad\leq M\left(\sum_{k=n+1}^{n+q}|x_k|^2\right)^{1/2}
 +M\left(\sum_{i=m+1}^{m+p}|y_i|^2\right)^{1/2}.
\end{align*}
It follows that the limit
\begin{equation}
  \lim_{m,n\to\infty}
  \sum_{i=1}^{m}\sum_{k=1}^{n}a_{ik}x_k y_i
  \tag{9}\label{eq:58:9}
\end{equation}
exists independently of the manner in which $m$ and $n$ tend to infinity;
under \eqref{eq:58:7} its numerical value does not exceed $M$.

Returning to the more special formula \eqref{eq:58:8}, we may interpret it
as follows. Besides the substitution \eqref{eq:56:4}, the matrix $(a_{ik})$
defines a second linear substitution
\begin{equation}
  y''_k=\sum_{i=1}^{\infty}a_{ik}y_i,
  \tag{10}\label{eq:58:10}
\end{equation}
and the two substitutions are related by
\begin{equation}
  \sum_{i=1}^{\infty}x'_i y_i
  =\sum_{k=1}^{\infty}x_k y''_k.
  \tag{11}\label{eq:58:11}
\end{equation}

% [Source printed page 85]
We shall say that the substitutions \eqref{eq:56:4} and
\eqref{eq:58:10} are transposes of one another. For brevity, the former
will be denoted by $A$ and its transpose by $\mathfrak A$. Their bounds
satisfy
\[
  M_A=M_{\mathfrak A}.
\]

\section{Inversion of Linear Substitutions}\label{sec:inversion-linear-substitutions}

\numberedparagraph{59}
As in algebra, the \emph{product} of two linear substitutions $A,B$ is the
linear substitution obtained by performing them successively:
\[
  AB(x_k)=A[B(x_k)].
\]
In general this product depends upon the order of its factors. If
$(a_{ik})$, $(b_{ik})$, and $(c_{ik})$ are the matrices corresponding to
$A$, $B$, and $C=AB$, respectively, then
\[
  c_{ik}=\sum_{j=1}^{\infty}a_{ij}b_{jk}
  \qquad(i,k=1,2,\ldots).
\]
If $\mathfrak A$ and $\mathfrak B$ are the transposes of $A$ and $B$, the
transpose of $AB$ is $\mathfrak B\mathfrak A$.

We denote by $E$ the identity substitution
$x'_k=x_k$ $(k=1,2,\ldots)$; its matrix $(e_{ik})$ has $e_{ik}=1$ when
$i=k$ and $e_{ik}=0$ when $i\neq k$.

\numberedparagraph{60}
Given a linear substitution $A$, we seek its reciprocal, or inverse,
$A^{-1}$, satisfying
\begin{equation}
  A^{-1}A=AA^{-1}=E.
  \tag{12}\label{eq:60:12}
\end{equation}
When it exists, $A^{-1}$ is uniquely determined by \eqref{eq:60:12}; in
fact either one of the two relations alone already determines it, provided
an inverse is known to exist.
% [Source printed page 86]
For if $D A=E$, then
$D=DE=DAA^{-1}=EA^{-1}=A^{-1}$. Suppose therefore that $A^{-1}$ exists,
and let $\mathfrak A'$ temporarily denote its transpose. The substitutions
$\mathfrak A\mathfrak A'$ and $\mathfrak A'\mathfrak A$ are the transposes
of $A^{-1}A=E$ and $AA^{-1}=E$; hence
\[
  \mathfrak A\mathfrak A'=\mathfrak A'\mathfrak A=E.
\]
Thus $\mathfrak A'$ is the inverse of $\mathfrak A$ and may be denoted by
$\mathfrak A^{-1}$.

Let $1/m$ be a common bound for $A^{-1}$ and $\mathfrak A^{-1}$. If $A$
maps $(x_k)$ to $(x'_i)$ and $\mathfrak A$ maps $(y_i)$ to $(y''_k)$, then
\begin{equation}
  m^2\sum_{k=1}^{\infty}|x_k|^2
  \leq\sum_{i=1}^{\infty}|x'_i|^2
  =\sum_{i=1}^{\infty}
    \left|\sum_{k=1}^{\infty}a_{ik}x_k\right|^2,
  \tag{13}\label{eq:60:13}
\end{equation}
\begin{equation}
  m^2\sum_{i=1}^{\infty}|y_i|^2
  \leq\sum_{k=1}^{\infty}|y''_k|^2
  =\sum_{k=1}^{\infty}
    \left|\sum_{i=1}^{\infty}a_{ik}y_i\right|^2.
  \tag{14}\label{eq:60:14}
\end{equation}
These inequalities give a necessary condition for the existence of
$A^{-1}$: the two quotients
\[
  \frac{\displaystyle\sum_i\left|\sum_k a_{ik}x_k\right|^2}
       {\displaystyle\sum_k|x_k|^2},
  \qquad
  \frac{\displaystyle\sum_k\left|\sum_i a_{ik}y_i\right|^2}
       {\displaystyle\sum_i|y_i|^2}
\]
must remain above some fixed number $m^2>0$. We shall show that this condition
is also sufficient.

Suppose it is satisfied, and consider systems \eqref{eq:56:4} and
\eqref{eq:58:10}, regarding $a_{ik}$, $x'_i$, and $y''_k$ as given. In
\eqref{eq:58:11} put $\mu_1,\ldots,\mu_n,0,0,\ldots$ in place of the
$y_i$. Applying the Cauchy--Lagrange inequality gives
% [Source printed page 87]
\[
 \left|\sum_{i=1}^{n}\mu_i x'_i\right|^2
 \leq\sum_{i=1}^{n}|\mu_i|^2\sum_{i=1}^{\infty}|x'_i|^2
 \leq\frac1{m^2}\sum_{i=1}^{\infty}|x'_i|^2
   \sum_{k=1}^{\infty}
     \left|\sum_{i=1}^{n}\mu_i a_{ik}\right|^2.
\]
This is precisely the condition for system \eqref{eq:56:4} to possess a
solution $(x_k)$ such that
\begin{equation}
  \sum_{k=1}^{\infty}|x_k|^2
  \leq\frac1{m^2}\sum_{i=1}^{\infty}|x'_i|^2.
  \tag{15}\label{eq:60:15}
\end{equation}
Consequently, for every given element $(x'_i)$, system
\eqref{eq:56:4} has at least one solution satisfying \eqref{eq:60:15}.
Likewise, for every given $(y''_k)$, system \eqref{eq:58:10} has at least
one solution $(y_i)$ such that
\begin{equation}
  \sum_{i=1}^{\infty}|y_i|^2
  \leq\frac1{m^2}\sum_{k=1}^{\infty}|y''_k|^2.
  \tag{16}\label{eq:60:16}
\end{equation}

The solutions of \eqref{eq:56:4} and \eqref{eq:58:10} are unique. For if,
for example, \eqref{eq:56:4} had two distinct solutions, its homogeneous
system would have a nonzero solution; but for the homogeneous system the
right-hand side of \eqref{eq:60:15} vanishes, a contradiction. Thus each
given $(x'_i)$ determines one element $(x_k)$, and each given $(y''_k)$ one
element $(y_i)$. These correspondences are distributive, by uniqueness,
and bounded by \eqref{eq:60:15} and \eqref{eq:60:16}. They are therefore
linear substitutions, and by their definition they are precisely
$A^{-1}$ and $\mathfrak A^{-1}$.

\numberedparagraph{61}
The entries $a_{ik}^{(-1)}$ of the matrix corresponding to $A^{-1}$ and
$\mathfrak A^{-1}$ may be calculated from system \eqref{eq:56:4} (or from
\eqref{eq:58:10}): $a_{ik}^{(-1)}$ is the value of $x_k$ obtained after
putting $x'_i=1$ and $x'_j=0$ for $j\neq i$.
% [Source printed page 88]
Once the $a_{ik}^{(-1)}$ have been found in this way, $A^{-1}$ and
$\mathfrak A^{-1}$ are represented by
\[
  x_k=\sum_{i=1}^{\infty}a_{ik}^{(-1)}x'_i,
  \qquad
  y_i=\sum_{k=1}^{\infty}a_{ik}^{(-1)}y''_k.
\]
Formula~\eqref{eq:47:40} of the \hyperref[chap:3]{preceding chapter} gives
\translatornote{In the French original the row involving \(a_{ik}\) is placed
last, its entries carry \(i\) instead of \(k\), and the quotient has a prefixed
minus sign.  Expanding formula~\eqref{eq:47:40} gives the corrected row position, index,
and sign used here.}
\[
 a_{ik}^{(-1)}=\lim_{n\to\infty}
 \frac{
 \begin{vmatrix}
 \alpha_{11}&\cdots&\alpha_{n1}\\
 \vdots&&\vdots\\
 \alpha_{1,i-1}&\cdots&\alpha_{n,i-1}\\
 \overline{a_{1k}}&\cdots&\overline{a_{nk}}\\
 \alpha_{1,i+1}&\cdots&\alpha_{n,i+1}\\
 \vdots&&\vdots\\
 \alpha_{1n}&\cdots&\alpha_{nn}
 \end{vmatrix}}
 {\begin{vmatrix}
 \alpha_{11}&\cdots&\alpha_{n1}\\
 \vdots&&\vdots\\
 \alpha_{1n}&\cdots&\alpha_{nn}
 \end{vmatrix}},
 \qquad
 \alpha_{ij}=\sum_{k=1}^{\infty}\overline{a_{ik}}a_{jk}.
\]

The criterion just established may also be expressed as follows. Consider
the two infinite Hermitian forms
\begin{equation}
  \sum_{i=1}^{\infty}\left|\sum_{k=1}^{\infty}a_{ik}x_k\right|^2,
  \qquad
  \sum_{k=1}^{\infty}\left|\sum_{i=1}^{\infty}a_{ik}x_i\right|^2.
  \tag{17}\label{eq:61:17}
\end{equation}
The inequalities \eqref{eq:60:13}--\eqref{eq:60:14} state that both forms
remain at least $m^2>0$ for every element $(x_k)$ satisfying
\begin{equation}
  \sum_{k=1}^{\infty}|x_k|^2=1.
  \tag{18}\label{eq:61:18}
\end{equation}
When this is so, $A$ has an inverse $A^{-1}$; the lower limits of the two
Hermitian forms under condition \eqref{eq:61:18} are equal, and the
reciprocal of their common square root is the bound $M_{A^{-1}}$.

% [Source printed page 89]
\numberedparagraph{62}
The criterion just obtained appears here as a particular consequence of
the general results of the \hyperref[chap:3]{preceding chapter}. Historically, however, it
was discovered by Toeplitz independently of that general theory, and even
before the latter had been developed.\footnote{Toeplitz,
\emph{Die Jacobische Transformation der quadratischen Formen von unendlich
vielen Ver\"anderlichen}, \emph{Nachrichten von der Gesellschaft der
Wissenschaften zu G\"ottingen}, 1907, pp.~101--109.} Another very simple and
interesting proof is due to Hilb.\footnote{Hilb, \emph{Ueber die Aufl\"osung
von Gleichungen mit unendlich vielen Unbekannten}, \emph{Sitzungsberichte
der Physikalisch-Medizinischen Soziet\"at zu Erlangen}, 1908, pp.~84--89.}
Those two proofs do not appear susceptible of extension to more general
cases, whereas ours extends immediately to every $p\geq1$.

To approach Toeplitz's line of thought, let $b_{ik}$ and $c_{ik}$ be the
coefficients of $x_i\overline{x_k}$ in the two Hermitian forms
\eqref{eq:61:17}, and let $B,C$ denote the corresponding substitutions.
Writing $\overline{\mathfrak A}$ for the substitution corresponding to the
conjugate-transposed matrix, we have
\[
  B=\overline{\mathfrak A}A,
  \qquad
  C=A\overline{\mathfrak A}.
\]
If $A^{-1}$ exists, then clearly
\[
  B^{-1}=A^{-1}\overline{\mathfrak A}^{\,-1},
  \qquad
  C^{-1}=\overline{\mathfrak A}^{\,-1}A^{-1}.
\]
Conversely, if $B$ and $C$ have inverses, then $A$ has the left inverse
$B^{-1}\overline{\mathfrak A}$ and the right inverse
$\overline{\mathfrak A}C^{-1}$; these are equal, since
\[
 (B^{-1}\overline{\mathfrak A})A
 =B^{-1}(\overline{\mathfrak A}A)=B^{-1}B=E,
  \qquad
 A(\overline{\mathfrak A}C^{-1})
 =(A\overline{\mathfrak A})C^{-1}=CC^{-1}=E,
\]
and
\[
 B^{-1}\overline{\mathfrak A}
 =B^{-1}\overline{\mathfrak A}CC^{-1}
 =B^{-1}\overline{\mathfrak A}A\overline{\mathfrak A}C^{-1}
 =E\overline{\mathfrak A}C^{-1}
 =\overline{\mathfrak A}C^{-1}.
\]
Everything is thus reduced to proving the following proposition.

% [Source printed page 90]
\begin{theorem}
If, as $(x_k)$ varies under condition \eqref{eq:61:18}, the positive
Hermitian form
\begin{equation}
  B(x,\overline{x})=
  \sum_{i,k=1}^{\infty}b_{ik}x_i\overline{x_k}
  \tag{19}\label{eq:62:19}
\end{equation}
has a lower limit $J>0$, then the linear substitution $B=(b_{ik})$ has an
inverse $B^{-1}$.
\end{theorem}

For this purpose consider the reductions
\begin{equation}
  [B(x,\overline{x})]_n
  =\sum_{i,k=1}^{n}b_{ik}x_i\overline{x_k}.
  \tag{20}\label{eq:62:20}
\end{equation}
Under the condition $\sum_{k=1}^{n}|x_k|^2=1$, these have lower limits
$J_n$. Plainly $J_n\geq J$; moreover, the $J_n$ form a decreasing sequence
and in fact tend to $J$, though this last fact is not needed below. What is
essential is that all $J_n$ remain above a positive lower bound.

Let $B_n$ be the linear substitution in $n$ variables with coefficients
$b_{ik}$ $(i,k=1,\ldots,n)$. The finite-dimensional analogue of the
proposition follows at once from determinant theory. Since $J_n>0$, the
homogeneous system
\[
  \sum_{i=1}^{n}b_{ik}x_i=0
  \qquad(k=1,\ldots,n)
\]
has only the solution $x_1=\cdots=x_n=0$; hence the determinant
$|b_{ik}|_n$ does not vanish and $B_n$ has a well-defined inverse
$B_n^{-1}$. Its bound is $1/J_n$.\footnote{For the purpose at hand it is
enough to know that this bound does not exceed $1/J_n$, since we need only
show that the upper bound does not increase without limit with $n$. If one
is content with this inequality rather than exact equality, the
 Cauchy--Lagrange inequality gives, with
 $x'_i=\sum_{k=1}^{n}b_{ik}\overline{x_k}$,
\[
 J_n^2\left(\sum_{i=1}^{n}|x_i|^2\right)^2
 \leq\left(\sum_{i,k=1}^{n}b_{ik}x_i\overline{x_k}\right)^2
 \leq\sum_{i=1}^{n}|x_i|^2
      \sum_{i=1}^{n}\left|\sum_{k=1}^{n}b_{ik}\overline{x_k}\right|^2
 =\sum_{i=1}^{n}|x_i|^2\sum_{i=1}^{n}|x'_i|^2,
\]
whence
\[
 \frac{\sum_i|x_i|^2}{\sum_i|x'_i|^2}\leq\frac1{J_n^2}.
\]}

% [Source printed page 91]
Let $n$ increase without limit. The coefficients of $B_n^{-1}$ tend to
well-defined limits $b_{ik}^{(-1)}$, and the matrix
$(b_{ik}^{(-1)})$ defines a linear substitution $B^{-1}$, which is the
inverse of $B$. In short, in the present case the principle of reductions
applies.

All this follows immediately from the principles of the \hyperref[chap:3]{preceding chapter},
especially the theorems of \hyperref[sec:40]{nos.~40} and
\hyperref[sec:42]{42}. A direct proof conforming to those principles would
be just as easy, but its details would take us far from Toeplitz's line of
thought. Toeplitz places himself at the standpoint of pure algebra, draws
from it everything possible, and uses unlimited algorithms only when the
power of algebraic methods has been exhausted.

There is in fact a considerable class of linear substitutions whose
inversion, as regards the computation of coefficients, is accomplished by
finite algorithms. Consider a substitution $D$ such that, for every $n$,
$x'_n$ depends only on $x_1,\ldots,x_n$. In its matrix $(d_{ik})$ all
entries to the right of the diagonal vanish. When the inverse $D^{-1}$
exists, its coefficients require only the solution of linear systems with
finitely many unknowns; the coefficients of $D_n^{-1}$ are at the same time
those of $D^{-1}$.

There is another problem requiring only finite algorithms: decompose the
substitution
% [Source printed page 92]
$B$ in the form
\[
  B=DD^*,
\]
where $D$ is of the triangular type just considered and $D^*$ is its
conjugate transpose. The coefficients $d_{ik}$ are obtained successively by
making the analogous decompositions of the substitutions $B_n$.

From this standpoint, inversion of $B$ consists first in forming the
decomposition $B=DD^*$, then determining $D^{-1}$, and finally---only here
are unlimited algorithms used---forming the product $(D^*)^{-1}D^{-1}$.
This product is plainly $B^{-1}$. It is understood that these operations
lead to bounded, hence linear, substitutions; this is easily verified.

Toeplitz calculates the coefficients of $D^{-1}$ by decomposing the
$B_n^{-1}$ rather than the $B_n$. Let $\Delta_n$ be the determinant of the
first $n$ rows and columns of $(b_{ik})$, and let
$\left(\begin{smallmatrix}i\\k\end{smallmatrix}\right)_n$ denote its
first-order minors. Put\translatornote{The French original gives the condition
\(i\ge k>1\), which omits every \(d_{i1}^{(-1)}\) with \(i>1\).  The identity
immediately below requires the corrected condition \(i>1\), \(1\le k\le i\).}
\[
 d_{11}^{(-1)}=\frac1{\sqrt{b_{11}}},\qquad
 d_{ik}^{(-1)}=
 \frac{\left(\begin{smallmatrix}i\\k\end{smallmatrix}\right)_i}
      {\sqrt{\Delta_{i-1}\Delta_i}}
 \quad(i>1,\ 1\leq k\leq i),
 \qquad d_{ik}^{(-1)}=0\quad(i<k).
\]
If $D_1^{-1},D_2^{-1},\ldots$ are the substitutions in
$1,2,\ldots$ variables formed with these coefficients, an identity going
back to Lagrange gives
\[
 \sum_{i=1}^{n}\left|\sum_{k=1}^{i}d_{ik}^{(-1)}x_k\right|^2
 =\frac1{\Delta_n}\sum_{i,k=1}^{n}
   \left(\begin{smallmatrix}i\\k\end{smallmatrix}\right)_n
   x_i\overline{x_k};
\]
that is, $(D_n^*)^{-1}D_n^{-1}=B_n^{-1}$.\footnote{The systems of
equations corresponding to substitutions of the type $D^*$ have the
special form considered in \hyperref[sec:12]{no.~12}. The reader may examine whether the
method indicated there---the principle of reductions---can be applied
here.}

\numberedparagraph{63}
Unlike the method just sketched, Hilb's method belongs to analysis. It is
based on the study of the series
\begin{equation}
  E+A+A^2+A^3+\cdots,
  \tag{21}\label{eq:63:21}
\end{equation}
% [Source printed page 93]
where $A^2,A^3,\ldots$ denote the iterated substitutions
\[
 A^2=AA,\qquad A^3=AA^2=A^2A=AAA,\quad\ldots.
\]
At first sight \eqref{eq:63:21} is merely a formal expansion of
$(E-A)^{-1}$, analogous to the power-series expansion of $(1-z)^{-1}$.
Hilb uses the fact that, as long as $M_A<1$, the series converges and gives
the inverse of $E-A$. We shall return later to the precise meaning of this
statement. For the moment it is enough to note that it implies, among other
things, that the coefficients of
\[
  S_n=E+A+\cdots+A^{n-1}
\]
tend to definite limits $s_{ik}$, and that these are exactly the
coefficients of the inverse sought.

Recall the theorem to be proved: if the positive Hermitian form
\eqref{eq:62:19} has lower limit $J>0$, then the corresponding substitution
$B$ has an inverse. For a positive numerical constant $\alpha$, write
\[
 B=\frac1\alpha[E-(E-\alpha B)].
\]
This gives the formal expansion
\[
 B^{-1}=\alpha\{E+(E-\alpha B)+(E-\alpha B)^2+\cdots\}.
\]
It remains to choose $\alpha$ so that $M_{E-\alpha B}<1$. The square of
this bound is the upper limit, under \eqref{eq:61:18}, of
\begin{align*}
 \sum_{i=1}^{\infty}\left|x_i-\alpha\sum_{k=1}^{\infty}b_{ik}x_k\right|^2
 &=\sum_{i=1}^{\infty}|x_i|^2
 -2\alpha\sum_{i,k=1}^{\infty}b_{ik}x_i\overline{x_k}\\
 &\quad+\alpha^2\sum_{i=1}^{\infty}
    \left|\sum_{k=1}^{\infty}b_{ik}x_k\right|^2.
\end{align*}
Consequently
\[
 M_{E-\alpha B}^2\leq1-2\alpha J+\alpha^2M_B^2,
\]
and, since $J>0$, the right-hand side is less than $1$ whenever
\[
 0<\alpha<\frac{2J}{M_B^2}.
\]
Every such $\alpha$ therefore determines $B^{-1}$.\footnote{A somewhat
more exact discussion gives the bound $2/M_B$. The expansion converges for
every smaller value of $\alpha$ and ceases to converge for
$\alpha=2/M_B$; convergence is fastest for
$\alpha=2/(J+M_B)$.}

% [Source printed page 94]
\numberedparagraph{64}
Let us emphasize once more the fact from which Hilb started and which will
be useful later:
\begin{theorem}
If $M_A<1$, the substitution $E-A$ has an inverse.
\end{theorem}
We have seen how Toeplitz's criteria follow from it. Conversely those
general criteria include the special hypothesis $M_A<1$. Indeed, the two
Hermitian forms involved are
\[
 \sum_{i=1}^{\infty}
   \left|x_i-\sum_{k=1}^{\infty}a_{ik}x_k\right|^2,
 \qquad
 \sum_{k=1}^{\infty}
   \left|x_k-\sum_{i=1}^{\infty}a_{ik}x_i\right|^2.
\]
The triangle inequality gives
\[
 \sum_{i=1}^{\infty}
   \left|x_i-\sum_{k=1}^{\infty}a_{ik}x_k\right|^2
 \geq(1-M_A)^2\sum_{k=1}^{\infty}|x_k|^2,
\]
and the same inequality for the transposed form. Thus Toeplitz's conditions
are satisfied when $M_A<1$; $E-A$ has an inverse and
\[
  M_{(E-A)^{-1}}\leq\frac1{1-M_A}.
\]

\section{Completely Continuous Substitutions}
\label{sec:completely-continuous-substitutions}

\numberedparagraph{65}
We have just obtained criteria for a substitution $A$ to possess an inverse
$A^{-1}$. While these conditions hold, systems \eqref{eq:56:4} and
\eqref{eq:58:10} have unique solutions, obtained by applying $A^{-1}$ and
$\mathfrak A^{-1}$ to the given elements. Up to this point the substitutions
behave exactly as if there were only finitely many variables. What happens
when the inverse does not exist?

% [Source printed page 95]
Recall the finite-dimensional result. When the inverse does not exist,
each of the homogeneous systems
\[
 \sum_{k=1}^{n}a_{ik}x_k=0\quad(i=1,\ldots,n),
 \qquad
 \sum_{i=1}^{n}a_{ik}x_i=0\quad(k=1,\ldots,n)
\]
has one or more solutions. If $(x_i^0)$ is a solution of the second, then
the nonhomogeneous system
\[
  \sum_{k=1}^{n}a_{ik}x_k=x'_i
\]
can have a solution only if
\[
  \sum_{i=1}^{n}x'_i x_i^0=0.
\]
Conversely, if this relation holds for every solution of the transposed
homogeneous system, the nonhomogeneous system has solutions; and the two
homogeneous systems have the same number of independent solutions.

These simple results fail completely for infinitely many variables. For
example, the substitution
\[
  x'_k=x_{k+1}\qquad(k=1,2,\ldots)
\]
has no inverse. Yet its two homogeneous systems are
\[
  x_k=0\quad(k=2,3,\ldots),
  \qquad
  x_i=0\quad(i=1,2,\ldots),
\]
which are not of the same type: in the first $x_1$ is arbitrary, whereas
the second has only the trivial solution.
% [Source printed page 96]
The substitution $x'_k=x_k/k$ also has no inverse, since an inverse would
map $x'_k=1/k$ to $x_k=1$, which is not an element of Hilbert space; but
both associated homogeneous systems have only the trivial solution. On
the other hand, the two homogeneous systems corresponding to
$x'_k=x_1/k$ both have infinitely many independent solutions.

\numberedparagraph{66}
There is nevertheless a broad class of substitutions for which the
finite-dimensional results remain valid. If, for example, the double
series $\sum_{i,k}|a_{ik}|^2$ converges, the substitution $E-A$ has the
required behavior, as follows from the method of infinite determinants.
This is only a special case of the one we shall study, in which $A$ is
assumed \emph{completely continuous}.

Every linear substitution is continuous: ordinary convergence of elements
implies ordinary convergence of their transforms, and strong convergence
implies strong convergence. A linear substitution $A$ will be called
\emph{completely continuous} when it maps every ordinarily convergent
sequence $(x_k^{(n)})$ to a sequence $(x_k^{(n)\prime})$ which converges
strongly to the corresponding element $(x'_k)$.\footnote{In Hilbert's
theory the bilinear form $A(x,y)$ is called completely continuous
(\emph{vollstetig}) when
\[
 (x_k^{(n)})\to(x_k),\qquad(y_k^{(n)})\to(y_k)
\]
imply $A(x^{(n)},y^{(n)})\to A(x,y)$. The definition in the text and
Hilbert's are related by the following fact. Given $(u_k)$ and a sequence
$(u_k^{(n)})$ in Hilbert space, the relation
\[
 \sum_{k=1}^{\infty}u_k^{(n)}x_k^{(n)}
 \longrightarrow\sum_{k=1}^{\infty}u_kx_k
\]
for every element $(x_k)$ and every sequence $(x_k^{(n)})\to(x_k)$ holds
if and only if $(u_k^{(n)})$ tends strongly to $(u_k)$. Compare the remark
following formula~\eqref{eq:40:28} of \hyperref[chap:3]{Chapter III}; the fact stated there proves the
sufficiency. This correspondence between completely continuous
substitutions and forms also shows that the transpose $\mathfrak A$ of a
completely continuous substitution is completely continuous.}

% [Source printed page 97]
For very simple examples showing the nature of this restriction, consider
\[
  x'_k=a_kx_k\qquad(k=1,2,\ldots).
\]
These formulas define a bounded, hence continuous, substitution exactly
when the $a_k$ are bounded; it is completely continuous exactly when
$a_k\to0$.

In particular, $E$ is not completely continuous. It follows directly from
the definition that if at least one of $A,B$ is completely continuous, then
$AB$ is also. A completely continuous substitution therefore cannot have
an inverse, for otherwise $E=AA^{-1}$ would be completely continuous.

If, with the exception of a single row or a single column, all coefficients
of $A$ vanish, then $A$ is completely continuous. Sums of finitely many
completely continuous substitutions are also completely continuous.
Consequently, if some $m$ satisfies $a_{ik}=0$ whenever both indices exceed
$m$, then $A$ is completely continuous.

\numberedparagraph{67}
We shall study the systems corresponding to $E-A$, assuming $A$ completely
continuous. This study could be connected with the general arguments of
the \hyperref[chap:3]{preceding chapter}, especially Schmidt's theory. I prefer to give a
special method which is a kind of generalization of the reasoning of
\hyperref[sec:25]{no.~25}. It seems to have been invented by A.~C. Dixon,\footnote{Dixon, \emph{On a class
of matrices of infinite order and on the existence of ``matricial''
functions on a Riemann surface}, received May~15, 1901,
\emph{Transactions of the Cambridge Philosophical Society}, vol.~19,
pp.~190--233.} who applied it to substitutions on bounded systems
$(x_k)$.\footnote{Dixon assumes constants $\alpha_k$ such that
$|a_{ik}|\leq\alpha_k$ and $\sum_k\alpha_k$ converges. The weaker condition
$A_i=\sum_k|a_{ik}|\to0$ as $i\to\infty$ would suffice.

For substitutions acting on bounded systems $(x_k)$, define convergence by
$x_k^{(n)}\to x_k$ for every $k$ together with a common bound on all
$|x_k^{(n)}|$, and strong convergence by uniform convergence in $k$. Then
a matrix defines a linear substitution exactly when
$A_i=\sum_k|a_{ik}|$ remains bounded, and it is completely continuous when
$A_i\to0$.} With a slight modification, Dixon's method serves our
present substitutions.

% [Source printed page 98]
Given $A=(a_{ik})$, let $R_n$ be obtained by replacing by zeros all entries
in its first $n$ rows and first $n$ columns, and let $M_{R_n}$ be its bound.
As $n$ increases these bounds do not increase. Complete continuity implies
the much stronger assertion
\[
  M_{R_n}\longrightarrow0.
\]
Indeed, for each $n$ choose an element $(x_k^{(n)})$ with
$x_1^{(n)}=\cdots=x_n^{(n)}=0$ and $\sum_k|x_k^{(n)}|^2=1$, for which
\[
 \sum_{i=n+1}^{\infty}
 \left|\sum_{k=n+1}^{\infty}a_{ik}x_k^{(n)}\right|^2
\]
comes within $1/n$ of its upper bound $M_{R_n}^2$.
% [Source printed page 99]
The sequence $(x_k^{(n)})$ tends ordinarily to zero; therefore its images
under $A$ tend strongly to zero:
\[
 \sum_{i=1}^{\infty}
 \left|\sum_{k=n+1}^{\infty}a_{ik}x_k^{(n)}\right|^2\longrightarrow0.
\]
Together with
\[
 M_{R_n}^2\leq\frac1n+
 \sum_{i=n+1}^{\infty}
 \left|\sum_{k=n+1}^{\infty}a_{ik}x_k^{(n)}\right|^2
 \leq\frac1n+
 \sum_{i=1}^{\infty}
 \left|\sum_{k=n+1}^{\infty}a_{ik}x_k^{(n)}\right|^2,
\]
this proves $M_{R_n}\to0$. In particular, for sufficiently large $n$,
$M_{R_n}<1$. Choose such an $n$. By \hyperref[sec:64]{no.~64}, $E-R_n$
then has an inverse; write
\[
  (E-R_n)^{-1}=E+B.
\]
The first $n$ rows and first $n$ columns of $(b_{ik})$ consist entirely of
zeros.

\numberedparagraph{68}
For simplicity decompose
\[
  A=A_n+P_n+Q_n+R_n,
\]
where $A_n$ has the coefficients of $A$ for $i\leq n$, $k\leq n$;
$P_n$ those for $i\leq n$, $k>n$; and $Q_n$ those for $i>n$, $k\leq n$,
all remaining coefficients being zero. Let $E_n$ be the projection
$x'_k=x_k$ for $k\leq n$, $x'_k=0$ for $k>n$. Multiplying
$E=(E+B)(E-R_n)$ by $P_n$ gives
\[
  P_n=P_n(E+B)(E-R_n).
\]
Adding $E_n-A_n-P_n-P_n(E+B)Q_n$ to both sides yields
% [Source printed page 100]
\begin{equation}
 E_n-A_n-P_n(E+B)Q_n
 =E_n-A_n-P_n+P_n(E+B)(E-Q_n-R_n).
 \tag{22}\label{eq:68:22}
\end{equation}
Let $(x'_k)=(E-A)(x_k)$. Applied to $(x_k)$, the right-hand side of
\eqref{eq:68:22} gives the same result as
$E_n+P_n(E+B)$ applied to $(x'_k)$. Put
\[
 C=E_n-A_n-P_n(E+B)Q_n,
 \qquad
 D=E_n+P_n(E+B).
\]
Then $C(x_k)=D(x'_k)$. The substitution $C$ may be regarded as one in $n$
variables, since its coefficients vanish if $i>n$ or $k>n$.

This symbolic computation reduces the infinite system
\begin{equation}
 x_i-\sum_{k=1}^{\infty}a_{ik}x_k=x'_i
 \qquad(i=1,2,\ldots)
 \tag{23}\label{eq:68:23}
\end{equation}
to one involving only $x_1,\ldots,x_n$. Indeed, set aside the first $n$
equations and write the others as
\[
 x_i-\sum_{k=n+1}^{\infty}a_{ik}x_k
 =x'_i+\sum_{k=1}^{n}a_{ik}x_k
 \qquad(i=n+1,n+2,\ldots).
\]
Solving for $x_{n+1},x_{n+2},\ldots$ gives
\[
 x_i=x'_i+\sum_{k=1}^{n}a_{ik}x_k
 +\sum_{j=n+1}^{\infty}b_{ij}
   \left(x'_j+\sum_{k=1}^{n}a_{jk}x_k\right).
\]
Substitution in the equations set aside eliminates the remaining infinite
tail and gives
% [Source printed page 101]
\begin{equation}
 \sum_{k=1}^{n}c_{ik}x_k
 =\sum_{j=1}^{\infty}d_{ij}x'_j
 \qquad(i=1,2,\ldots,n),
 \tag{24}\label{eq:68:24}
\end{equation}
where $c_{ik}$ and $d_{ij}$ are the coefficients of $C$ and $D$.

An analogous process transforms the transposed system
\begin{equation}
 x_k-\sum_{i=1}^{\infty}a_{ik}x_i=x''_k
 \tag{25}\label{eq:68:25}
\end{equation}
into
\begin{equation}
 \sum_{i=1}^{n}c_{ik}x_i
 =\sum_{j=1}^{\infty}d'_{jk}x''_j
 \qquad(k=1,2,\ldots,n).
 \tag{26}\label{eq:68:26}
\end{equation}
The $d'_{jk}$, coefficients of $E_n+(E+B)Q_n$, will not concern us for
the moment. The important point is that the coefficients on the left are
the same as in \eqref{eq:68:24}, since they form the transpose of $C$.
Thus the study of the infinite systems \eqref{eq:68:23} and
\eqref{eq:68:25} is reduced to that of the finite systems
\eqref{eq:68:24} and \eqref{eq:68:26}.

\numberedparagraph{69}
Two cases must be distinguished.

\emph{First case.} The determinant $|c_{ik}|_n$ does not vanish. Then the
finite substitution $C_n$ has an inverse $C_n^{-1}$, regarded also as an
infinite substitution by filling missing coefficients with zeros. Applying
$C_n^{-1}D$ to $(x'_i)$ yields the true values of $x_1,\ldots,x_n$ and
zeros for the remaining coordinates.
% [Source printed page 102]
Substituting those first $n$ values in the preceding expression for the
tail shows that
\[
 E-E_n+B+Q_nC_n^{-1}D+BQ_nC_n^{-1}D
\]
gives the true values of $x_{n+1},x_{n+2},\ldots$ and zeros for the first
$n$ coordinates. Hence
\[
 E-E_n+B+C_n^{-1}D+Q_nC_n^{-1}D+BQ_nC_n^{-1}D
\]
gives the solution of \eqref{eq:68:23}; equivalently,
\begin{equation}
\begin{split}
 (E-A)\bigl(&E-E_n+B+C_n^{-1}D\\
 &{}+Q_nC_n^{-1}D+BQ_nC_n^{-1}D\bigr)=E.
\end{split}
 \tag{27}\label{eq:69:27}
\end{equation}
An analogous argument for \eqref{eq:68:25} supplies a right inverse of
$E-A$. The two one-sided inverses must coincide; therefore the substitution
displayed in \eqref{eq:69:27} is $(E-A)^{-1}$.

\emph{Second case.} The determinant $|c_{ik}|_n$ vanishes.

% [Source printed page 103]
Then the two finite homogeneous systems
\begin{equation}
 \sum_{k=1}^{n}c_{ik}x_k=0\quad(i=1,\ldots,n),
 \qquad
 \sum_{i=1}^{n}c_{ik}x_i=0\quad(k=1,\ldots,n)
 \tag{28}\label{eq:69:28}
\end{equation}
have the same number $m$ of independent solutions. Denote bases for the
first and second systems by
\[
 (\alpha_1^{(1)},\ldots,\alpha_n^{(1)}),\ldots,
 (\alpha_1^{(m)},\ldots,\alpha_n^{(m)})
\]
and
\[
 (\beta_1^{(1)},\ldots,\beta_n^{(1)}),\ldots,
 (\beta_1^{(m)},\ldots,\beta_n^{(m)}).
\]

Pass to the infinite homogeneous systems obtained from
\eqref{eq:68:23} and \eqref{eq:68:25} by putting the given element equal
to zero. The first $n$ coordinates of any solution form a solution of the
corresponding finite system \eqref{eq:69:28}. Conversely, every solution
of either finite system extends uniquely to a solution of the corresponding
infinite system. For example, for the homogeneous system
\eqref{eq:68:23}, the remaining unknowns satisfy
\[
 x_i-\sum_{k=n+1}^{\infty}a_{ik}x_k
 =\sum_{k=1}^{n}a_{ik}x_k^0
 \qquad(i=n+1,n+2,\ldots),
\]
and are obtained by applying $(E+B)Q_n$ to
$(x_1^0,\ldots,x_n^0,0,0,\ldots)$. Thus every infinite homogeneous
solution is obtained by applying
\[
 E_n+(E+B)Q_n
\]
to a corresponding finite solution.
% [Source printed page 104]
This is precisely the substitution whose matrix is $(d'_{ik})$; hence the
general infinite solution is
\[
 x_i^0=\sum_{k=1}^{n}d'_{ik}x_k^0
 \qquad(i=1,2,\ldots).
\]
Applying this process to the $m$ solutions $\alpha$ gives $m$ linearly
independent solutions $(\alpha_k^{(j)})$ $(j=1,\ldots,m)$ of the infinite
homogeneous system \eqref{eq:68:23}, and every solution is their linear
combination. The same statement holds, with the same number $m$, for the
homogeneous system \eqref{eq:68:25} and the extended $\beta$ solutions.

Now consider the nonhomogeneous systems, say \eqref{eq:68:25}. Reduction
to \eqref{eq:68:26} is always possible. The finite system
\eqref{eq:68:26} is solvable if and only if
\[
 \sum_{k=1}^{n}\left(\sum_{j=1}^{\infty}d'_{jk}x''_j\right)x_k^0=0
\]
for every solution $(x_1^0,\ldots,x_n^0)$ of the first homogeneous system
\eqref{eq:69:28}. Equivalently,
\[
 \sum_{j=1}^{\infty}x''_j
   \left(\sum_{k=1}^{n}d'_{jk}x_k^0\right)
 =\sum_{j=1}^{\infty}x''_j x_j^0=0.
\]
But $(x_j^0)$ is the general solution of the infinite homogeneous system
\eqref{eq:68:23}. Thus \eqref{eq:68:25} is solvable precisely when
\[
  \sum_{k=1}^{\infty}x''_k x_k^0=0
\]
for every solution $(x_k^0)$ of the homogeneous system
\eqref{eq:68:23}.
% [Source printed page 105]
Whenever $(x''_k)$ is so chosen, \eqref{eq:68:25} has several solutions;
the general solution is obtained from any particular solution by adding an
arbitrary linear combination of the $m$ elements $(\beta_k^{(j)})$.
Analogous results hold for the nonhomogeneous system \eqref{eq:68:23}.

In summary:
\begin{theorem}
If $A$ is completely continuous, the infinite systems
\eqref{eq:68:23} and \eqref{eq:68:25} retain the essential properties of
finite systems containing the same number of equations and unknowns.
\end{theorem}

Observe that complete continuity was used only to infer that
$M_{R_n}<1$ for all sufficiently large $n$. All the preceding conclusions
remain valid if this last property alone is assumed.

\numberedparagraph{70}
The preceding considerations also permit a variable parameter $\lambda$ to
be introduced and the inverse $(E-\lambda A)^{-1}$ to be studied as a
function of it. Suppose $\lambda$ varies in the disk $|\lambda|\leq r$.
For completely continuous $A$, choose $n$ so that $M_{R_n}<1/r$. Then
$E-\lambda R_n$ has, for every such $\lambda$, an inverse $E+B(\lambda)$;
the coefficients of $B(\lambda)$ are holomorphic in the disk. The same is
true of the coefficients of the corresponding substitutions
$C(\lambda)$, $C_n(\lambda)$, and $D(\lambda)$.

The inverse $C_n^{-1}(\lambda)$ is computed by finite determinants, which
are holomorphic functions. Since a nonzero holomorphic function has only
finitely many zeros in a bounded closed domain,
$C_n^{-1}(\lambda)$ exists except possibly at finitely many values. If it
fails at $\lambda_0$, it can be written
\[
  C_n^{-1}(\lambda)=\frac{G(\lambda)}{(\lambda-\lambda_0)^m},
\]
with the coefficients of $G$ holomorphic at $\lambda_0$. Finally,
\[
\begin{split}
 (E-\lambda A)^{-1}
 &=E-E_n+B(\lambda)+C_n^{-1}(\lambda)D(\lambda)\\
 &\quad+\lambda Q_nC_n^{-1}(\lambda)D(\lambda)
 +\lambda B(\lambda)Q_nC_n^{-1}(\lambda)D(\lambda).
\end{split}
\]
% [Source printed page 106]
This proves the same assertions for $(E-\lambda A)^{-1}$: in brief, the
inverse is a meromorphic function of $\lambda$. The conclusion holds for
all $\lambda$, since the radius $r$ may be chosen arbitrarily.

In cases amenable to infinite determinants---for example when
$\sum_{i,k}|a_{ik}|^2$ converges---the result follows immediately from the
fact that the infinite determinant with entries $e_{ik}-\lambda a_{ik}$ is
holomorphic in $\lambda$.

\section{Sequences and Series of Substitutions}
\label{sec:sequences-series-substitutions}

\numberedparagraph{71}
In \hyperref[sec:63]{no.~63}, while explaining Hilb's method, we said that
the series
\[
  E+A+A^2+\cdots
\]
converges, when $M_A<1$, to $(E-A)^{-1}$; but we neither proved the
assertion nor specified its meaning. We shall fill this gap and at the same
time develop a few general facts about sequences of substitutions. The
arguments extend immediately to other limiting processes depending, for
example, on a continuous parameter rather than an integer index.

Let $(A_n)$ be an infinite sequence of linear substitutions, and let
$(a_{ik}^{(n)})$ be their matrices. Suppose:

1. all bounds $M_{A_n}$ are below one finite number $M$;

2. for every $i,k$, $a_{ik}^{(n)}\to a_{ik}$ as $n\to\infty$.

Then the $a_{ik}$ are the coefficients of a linear substitution $A$ such
that
\begin{equation}
  M_A\leq M.
  \tag{29}\label{eq:71:29}
\end{equation}
% [Source printed page 107]
Indeed, for every choice of the integers $l,m$ and the quantities
$x_1,\ldots,x_m$,
\[
 \sum_{i=1}^{l}\left|\sum_{k=1}^{m}a_{ik}x_k\right|^2
 =\lim_{n\to\infty}\sum_{i=1}^{l}
   \left|\sum_{k=1}^{m}a_{ik}^{(n)}x_k\right|^2
 \leq M^2\sum_{k=1}^{m}|x_k|^2.
\]
By \hyperref[sec:56]{no.~56}, this is precisely the condition for
$(a_{ik})$ to define a substitution $A$ satisfying \eqref{eq:71:29}.

Let $(x_k)$ now be arbitrary, and let $(x_k^{(n)\prime})$ and $(x'_k)$ be
the elements obtained by applying $A_n$ and $A$, respectively. Hypothesis
1 gives, for every $i,n$,
\[
  \sum_{k=1}^{\infty}|a_{ik}^{(n)}|^2\leq M^2,
\]
and therefore, by \hyperref[sec:40]{no.~40},
\[
 x_i^{(n)\prime}=\sum_{k=1}^{\infty}a_{ik}^{(n)}x_k
 \longrightarrow\sum_{k=1}^{\infty}a_{ik}x_k=x'_i.
\]
Moreover, all sums $\sum_k|x_k^{(n)\prime}|^2$ remain at most
$M^2\sum_k|x_k|^2$. Thus $(x_k^{(n)\prime})$ tends ordinarily to
$(x'_k)$, whatever the initial element $(x_k)$. Under these hypotheses we
shall say that the substitutions $A_n$ \emph{tend to} $A$.

\numberedparagraph{72}
Series \eqref{eq:63:21} belongs to a still more special type. Given an
infinite sequence $(A_n)$, we shall say that it tends \emph{uniformly} to
$A$ when
\[
  M_{A-A_n}\longrightarrow0.
\]
The two hypotheses of \hyperref[sec:71]{no.~71} then hold, since
\[
  M_{A_n}\leq M_A+M_{A-A_n},
  \qquad
  |a_{ik}-a_{ik}^{(n)}|\leq M_{A-A_n}.
\]
Consequently $A_n$ tends to $A$ also in the earlier sense.
% [Source printed page 108]
In particular, $A_n(x_k)$ tends ordinarily to $A(x_k)$. More is true:
\[
 \sum_{k=1}^{\infty}|x'_k-x_k^{(n)\prime}|^2
 \leq M_{A-A_n}^2\sum_{k=1}^{\infty}|x_k|^2,
\]
so the convergence is strong, and is uniform throughout every bounded
domain of Hilbert space (that is, every domain on which
$\sum_k|x_k|^2\leq G^2$).

Uniform convergence has the further remarkable property
\[
  M_{A_n}\longrightarrow M_A,
\]
which follows immediately from
\[
  |M_A-M_{A_n}|\leq M_{A-A_n}.
\]

We now ask: given a sequence $(A_n)$, how can one determine whether there
exists a substitution $A$ to which it converges uniformly? A necessary
condition follows from
\[
  M_{A_n-A_m}\leq M_{A-A_n}+M_{A-A_m};
\]
namely,
\begin{equation}
  M_{A_n-A_m}\longrightarrow0
  \quad\text{as $m,n\to\infty$ independently.}
  \tag{30}\label{eq:72:30}
\end{equation}
This condition is also sufficient. For if it holds, then
\[
 |M_{A_n}-M_{A_m}|\leq M_{A_n-A_m},
 \qquad
 |a_{ik}^{(n)}-a_{ik}^{(m)}|\leq M_{A_n-A_m}
\]
show that $M_{A_n}$ and every $a_{ik}^{(n)}$ have limits $M$ and $a_{ik}$.
Thus the hypotheses of \hyperref[sec:71]{no.~71} hold and $(a_{ik})$
defines a substitution $A$ to which $A_n$ tends. Likewise, for fixed $m$,
$A_n-A_m$ tends to $A-A_m$.
% [Source printed page 109]
Given $\varepsilon>0$, choose $m$ so large that
$M_{A_n-A_m}<\varepsilon$ for every $n>m$. By \eqref{eq:71:29},
$M_{A-A_m}\leq\varepsilon$, and hence $M_{A-A_m}\to0$.

\begin{theorem}
The sequence $(A_n)$ converges uniformly if and only if
$M_{A_n-A_m}\to0$ as $m,n\to\infty$ independently.
\end{theorem}

\numberedparagraph{73}
For the partial sums of \eqref{eq:63:21} this condition is fulfilled. Put
$S_n=E+A+\cdots+A^{n-1}$ and suppose $n>m$.\translatornote{The French text
reverses the two terms in the numerator of the final fraction below.  Their
order is corrected so that the estimate is positive for $n>m$ and $M_A<1$.}
Then
\[
 S_n-S_m=A^m+A^{m+1}+\cdots+A^{n-1},
\]
and
\[
 M_{S_n-S_m}
 \leq(M_A)^m+\cdots+(M_A)^{n-1}
 =\frac{(M_A)^m-(M_A)^n}{1-M_A}.
\]
If $M_A<1$, this tends to zero; hence $S_n$ converges uniformly to a
substitution $S$.

The following theorem, which follows immediately from
\[
 M_{BA-BA_n}\leq M_BM_{A-A_n},
 \qquad
 M_{AB-A_nB}\leq M_{A-A_n}M_B,
\]
shows that $S=(E-A)^{-1}$.

\begin{theorem}
If $A_n$ converges uniformly to $A^*$ and $B$ is any linear substitution,
then $BA_n$ and $A_nB$ converge uniformly to $BA^*$ and $A^*B$.
\end{theorem}

Apply this with $A_n=S_n$, $A^*=S$, and $B=E-A$. Since
\[
 (E-A)S_n=S_n(E-A)=E-A^n\longrightarrow E,
\]
we obtain
\[
  (E-A)S=S(E-A)=E,
  \qquad S=(E-A)^{-1}.
\]

% [Source printed page 110]
\numberedparagraph{74}
We record the more general proposition:
\begin{theorem}
If $A_n$ and $B_n$ converge uniformly to $A$ and $B$, respectively, then
$A_nB_n$ converges uniformly to $AB$.
\end{theorem}
Indeed,
\[
 M_{AB-A_nB_n}
 \leq M_A M_{B-B_n}+M_{A-A_n}M_{B_n},
\]
while $M_{A-A_n}\to0$, $M_{B-B_n}\to0$, and $M_{B_n}\to M_B$.

The analogous proposition for nonuniform convergence is false. For
example, let
\[
 A_n:\ x'_1=x_n,\quad x'_i=0\ (i>1),
 \qquad
 B_n:\ x'_n=x_1,\quad x'_i=0\ (i\neq n).
\]
Both $A_n$ and $B_n$ tend to substitutions all of whose coefficients
vanish. Yet every product $A_nB_n$ is the same substitution
\[
 C:\ x'_1=x_1,\quad x'_i=0\ (i>1),
\]
so $A_nB_n=C\to C$, and $C$ is not identically zero.

For completeness we add an easily proved result:
\begin{theorem}
If $A_n\to A$ and $B_n\to B$ and at least one of the two sequences
converges uniformly, then $A_nB_n\to AB$.\footnote{This includes the
corollary: if $A_n\to A$ and $B$ is any fixed
linear substitution, then $A_nB\to AB$ and $BA_n\to BA$.}
\end{theorem}

\numberedparagraph{75}
There is another important case in which $A_nB_n\to AB$. The simplest
sequences tending to $A$ are its reductions: in $(a_{ik})$ set equal to
zero every entry for which at least one index exceeds $n$. In general these
reductions do not converge uniformly to $A$; that occurs only when $A$ is
completely continuous. But the reductions possess an essential special
property.
% [Source printed page 111]
Formula \eqref{eq:56:3} shows that their successive application to any
element $(x_k)$ produces elements tending strongly to $(x'_k)$. We
therefore insert between convergence and uniform convergence the notion of
\emph{strong convergence}. If $A_n\to A$, we shall say that $A_n$ tends
strongly to $A$ when, for every $(x_k)$, the elements $A_n(x_k)$ tend
strongly to $A(x_k)$.

\begin{theorem}
If $A_n$ and $B_n$ tend strongly to $A$ and $B$, respectively, then
$A_nB_n$ tends strongly to $AB$.
\end{theorem}

The bounds $M_{A_n}$ and $M_{B_n}$ remain bounded, and therefore so do the
bounds of $A_nB_n$. It remains only to show that
$[AB-A_nB_n](x_k)$ tends strongly to zero. Write
\[
 AB-A_nB_n=(A-A_n)B+A_n(B-B_n).
\]
For the first term, $B$ maps $(x_k)$ to some $(y_k)$ and
$(A-A_n)(y_k)$ tends strongly to zero. For the second, put
\[
 (B-B_n)(x_k)=(y_k^{(n)}),
 \qquad
 A_n(y_k^{(n)})=(z_k^{(n)}).
\]
By hypothesis $(y_k^{(n)})$ tends strongly to zero, and
\[
 \sum_k|z_k^{(n)}|^2\leq M_{A_n}^2\sum_k|y_k^{(n)}|^2;
\]
hence $(z_k^{(n)})$ also tends strongly to zero.

In particular, if $A_n$ is the $n$th reduction of $A$, then for every
fixed $k$ the powers $A_n^k$ tend strongly to $A^k$.

\numberedparagraph{76}
Suppose $A_n$ tends uniformly to $A$ and each $A_n$ has an inverse. Can we
infer the existence of $A^{-1}$, and, if so, is it the limit of
$A_n^{-1}$? This question is closely connected with the principle of
reductions.

First suppose the bounds $M_{A_n^{-1}}$ remain bounded. Then
$A_n^{-1}$ converges uniformly and its limit is $A^{-1}$. Indeed,
% [Source printed page 112]
\[
 A_n^{-1}-A_m^{-1}=A_n^{-1}(A_m-A_n)A_m^{-1},
\]
so
\[
 M_{A_n^{-1}-A_m^{-1}}
 \leq M_{A_n^{-1}}M_{A_m^{-1}}M_{A_n-A_m}\longrightarrow0.
\]
Let the uniform limit be $A^*$. Since $A_n\to A$ and
$A_n^{-1}\to A^*$, the products $A_nA_n^{-1}$ and $A_n^{-1}A_n$ tend
uniformly to $AA^*$ and $A^*A$; but both products equal $E$. Thus
$AA^*=A^*A=E$, and $A^*=A^{-1}$.

Conversely, suppose $A^{-1}$ exists and only that $A_n\to A$ uniformly.
Then $A_n^{-1}$ exists for all sufficiently large $n$, namely whenever
\[
  M_{A-A_n}<\frac1{M_{A^{-1}}}.
\]
For in that case
\[
 M_{E-A^{-1}A_n}
 =M_{A^{-1}(A-A_n)}
 \leq M_{A^{-1}}M_{A-A_n}<1,
\]
so $A^{-1}A_n$ has an inverse $B$. From
$B A^{-1}A_n=E$ and $A^{-1}A_nB=E$, multiplication of the second equality
on the left by $A$ and on the right by $A^{-1}$ gives
$A_nBA^{-1}=E$; hence $BA^{-1}=A_n^{-1}$. Moreover,
\[
\begin{split}
 M_{A_n^{-1}}
 &\leq M_BM_{A^{-1}}\\
 &\leq
 \frac{M_{A^{-1}}}
 {1-M_{A^{-1}}M_{A-A_n}},
\end{split}
\]
which tends to $M_{A^{-1}}$. Thus the inverses exist eventually, their
bounds remain collectively bounded, and the preceding result gives
$A_n^{-1}\to A^{-1}$ uniformly.

% [Source printed page 113]
\begin{theorem}
Let $A_n$ converge uniformly to $A$. The inverse $A^{-1}$ exists if and
only if, from some index onward, the inverses $A_n^{-1}$ exist and their
bounds remain collectively bounded; in that case $A_n^{-1}$ converges
uniformly to $A^{-1}$.
\end{theorem}

The essential point of the argument, needed later, is this:
\begin{theorem}
If $A^{-1}$ exists and $M_B<1/M_{A^{-1}}$, then $(A+B)^{-1}$ also exists.
\end{theorem}
This is a generalization of \hyperref[sec:64]{no.~64}.

\numberedparagraph{77}
These results justify the principle of reductions in an important case. We
have said that the reductions $A_n$ of $A$ do not in general converge
uniformly to $A$; but if $A$ is completely continuous, they do.\footnote{The
converse is also true: if $M_{A-A_n}\to0$, then $A$ is completely
continuous. This is contained in the more general proposition that a
uniform limit of completely continuous substitutions is completely
continuous. It also covers the condition $M_{R_n}\to0$ of
\hyperref[sec:67]{no.~67}. Indeed, $A_n+P_n+Q_n$ is always completely
continuous, and $M_{R_n}\to0$ says that these substitutions converge
uniformly to $A$. Consequently, together with \hyperref[sec:67]{no.~67}, $M_{R_n}\to0$ is a
necessary and sufficient condition for complete continuity of $A$.}
\nobreak\hspace{0.12em}\translatornote{In this footnote the French original prints
\(A_n-P_n+Q_n\) instead of \(A_n+P_n+Q_n\), and concludes with \(B\) instead
of \(A\); both misprints are corrected above.}
Therefore $E-A_n$ converges uniformly to $E-A$. If $(E-A)^{-1}$ exists,
then from some index onward $(E-A_n)^{-1}$ exists and converges uniformly to
$(E-A)^{-1}$. Thus, for completely continuous $A$, whenever $(E-A)^{-1}$
exists its calculation may be based on the principle of reductions.

\section{Study of the Inverse
\texorpdfstring{\((E-\lambda A)^{-1}\)}{(E - lambda A) inverse} as a Function
of \texorpdfstring{\(\lambda\)}{lambda}}\label{sec:resolvent-function}

\numberedparagraph{78}
Let $\lambda$ be a complex parameter and study $E-\lambda A$ from the
standpoint of inversion.
% [Source printed page 114]
One advantage of this study is that it presents within a single object the
two principal cases. In general there are two classes of values of
$\lambda$: for some, $E-\lambda A$ has an inverse; for the others it does
not. We call the first values \emph{ordinary} and the second
\emph{singular}. These terms, borrowed from function theory, are justified
by the relations between that theory and the study of
$(E-\lambda A)^{-1}$. In studying completely continuous substitutions we
already touched this problem and observed facts concerning the analytic
nature of the inverse.

Some arguments below also apply to the inverse of
\[
 E-\sum_{k=1}^{n}\lambda^kA_k,
\]
or to other substitutions depending upon $\lambda$ by less simple laws.
We restrict ourselves to $E-\lambda A$.

\numberedparagraph{79}
Let $A$ be any linear substitution and consider $(E-\lambda A)^{-1}$ at
ordinary values of $\lambda$. It is convenient to introduce $A_\lambda$
by
\begin{equation}
  (E-\lambda A)^{-1}=E+\lambda A_\lambda.
  \tag{31}\label{eq:79:31}
\end{equation}
For every ordinary $\lambda\neq0$ this determines $A_\lambda$ uniquely; set
$A_0=A$, which, as we shall see, defines it by continuity. Equation
\eqref{eq:79:31} is equivalent to
\begin{equation}
  E-\lambda A=(E+\lambda A_\lambda)^{-1}.
  \tag{32}\label{eq:79:32}
\end{equation}

Let $\lambda$ and $\mu$ be ordinary values. From \eqref{eq:79:32},
\[
 (\mu-\lambda)E
 =\mu(E+\lambda A_\lambda)^{-1}
 -\lambda(E+\mu A_\mu)^{-1}.
\]
Multiplying on the left by $E+\lambda A_\lambda$ and on the right by
$E+\mu A_\mu$ gives
% [Source printed page 115]
\[
 (\mu-\lambda)(E+\lambda A_\lambda)(E+\mu A_\mu)
 =\mu(E+\mu A_\mu)-\lambda(E+\lambda A_\lambda),
\]
or
\[
 \lambda\mu\,[A_\lambda-A_\mu
 +(\mu-\lambda)A_\lambda A_\mu]=0.
\]
If $\lambda\mu\neq0$, it follows that
\begin{equation}
  A_\lambda-A_\mu+(\mu-\lambda)A_\lambda A_\mu=0.
  \tag{33}\label{eq:79:33}
\end{equation}
Interchanging $\lambda$ and $\mu$ gives
\begin{equation}
  A_\lambda-A_\mu+(\mu-\lambda)A_\mu A_\lambda=0.
  \tag{34}\label{eq:79:34}
\end{equation}
Thus multiplication of $A_\lambda$ and $A_\mu$ is commutative.

Equations \eqref{eq:79:33} and \eqref{eq:79:34} remain valid when
$\lambda\mu=0$. If, for example, $\mu=0$, they say
\[
 A-A_\lambda+\lambda AA_\lambda
 =A-A_\lambda+\lambda A_\lambda A=0.
\]
Indeed, writing \eqref{eq:79:32} more explicitly gives
\begin{equation}
 (E-\lambda A)(E+\lambda A_\lambda)
 =(E+\lambda A_\lambda)(E-\lambda A)=E,
 \tag{35}\label{eq:79:35}
\end{equation}
and hence
\[
 -\lambda(A-A_\lambda+\lambda AA_\lambda)
 =-\lambda(A-A_\lambda+\lambda A_\lambda A)=0.
\]
Division by $-\lambda$ is legitimate if $\lambda\neq0$; the case
$\lambda=\mu=0$ is evident. Thus \eqref{eq:79:33} and
\eqref{eq:79:34} hold in all cases.

\numberedparagraph{80}
Let ordinary values $\lambda_1,\lambda_2,\ldots$ tend to $\lambda^*$, and
suppose all $M_{A_{\lambda_n}}$ are below a number $M$. Either resolvent
identity gives
\[
  M_{A_{\lambda_n}-A_{\lambda_m}}
  \leq|\lambda_n-\lambda_m|M^2.
\]
Hence
% [Source printed page 116]
$A_{\lambda_n}$ converges uniformly to a substitution $A^*$. Then
$E+\lambda_nA_{\lambda_n}$ converges uniformly to
$E+\lambda^*A^*$. The two product sequences in \eqref{eq:79:35} converge
to
\[
 (E-\lambda^*A)(E+\lambda^*A^*),
 \qquad
 (E+\lambda^*A^*)(E-\lambda^*A),
\]
respectively. Since every term equals $E$, both limits equal $E$.
Therefore $\lambda^*$ is ordinary and $E+\lambda^*A^*$ is the inverse of
$E-\lambda^*A$.

\begin{theorem}
The limiting value $\lambda^*$ is ordinary, and
$M_{A_{\lambda^*}-A_{\lambda_n}}\to0$.
\end{theorem}

On the entire set of ordinary values, $M_{A_\lambda}$ is a continuous
function of $\lambda$. From \eqref{eq:79:33} and the inequalities for
bounds,
\[
 |M_{A_\lambda}-M_{A_\mu}|
 \leq|\lambda-\mu|M_{A_\lambda}M_{A_\mu}.
\]
When both bounds are nonzero this yields
\[
 \left|\frac1{M_{A_\lambda}}-\frac1{M_{A_\mu}}\right|
 \leq|\lambda-\mu|.
\]
Thus $1/M_{A_\lambda}$, and hence $M_{A_\lambda}$, is continuous wherever
the latter is nonzero. But $A_\lambda$ can vanish identically only if
$A=0$, by $A=A_\lambda+\lambda AA_\lambda$. Consequently
$M_{A_\lambda}$ is continuous throughout the set of ordinary values.

\numberedparagraph{81}
The set of ordinary values is open. Indeed, if $\mu$ is ordinary, every
$\lambda$ satisfying
\[
  |\lambda-\mu|<\frac1{M_{A_\mu}}
\]
is ordinary.
% [Source printed page 117]
For
\begin{align*}
 E-\lambda A
 &=E-\mu A-(\lambda-\mu)A\\
 &=E-\mu A-(\lambda-\mu)(E-\mu A)(E+\mu A_\mu)A\\
 &=E-\mu A-(\lambda-\mu)(E-\mu A)A_\mu\\
 &=(E-\mu A)[E-(\lambda-\mu)A_\mu].
\end{align*}
The first factor is invertible by hypothesis, and the second because
$|\lambda-\mu|M_{A_\mu}<1$. Its inverse is the uniformly convergent series
\[
 E+(\lambda-\mu)A_\mu+(\lambda-\mu)^2A_\mu^2+\cdots.
\]
Multiplication gives
\[
 (E-\lambda A)^{-1}
 =E+\lambda A_\mu
 +\lambda(\lambda-\mu)A_\mu^2
 +\lambda(\lambda-\mu)^2A_\mu^3+\cdots,
\]
and therefore
\[
 A_\lambda=A_\mu+(\lambda-\mu)A_\mu^2
 +(\lambda-\mu)^2A_\mu^3+\cdots.
\]
Equivalently,
\[
 \frac{A_\lambda^k-A_\mu^k}{\lambda-\mu}
 \longrightarrow kA_\mu^{k+1}
 \quad(\lambda\to\mu)
\]
uniformly. Thus $A_\mu$, as a function of $\mu$, has successive
derivatives
\[
 A_\mu^2,\quad 2A_\mu^3,\quad\ldots,\quad k!A_\mu^{k+1},\quad\ldots.
\]
In short, for ordinary $\lambda$, $A_\lambda$ and the other substitutions
entering the calculation have the character of holomorphic functions of
$\lambda$.

\numberedparagraph{82}
One may proceed further by applying the methods of function theory to
$A_\lambda$, especially the calculus of residues. We merely indicate some
results.

% [Source printed page 118]
Let $D$ be a domain bounded by finitely many closed curves, which may be
assumed analytic. The integral along these curves, with respect to
$\lambda$, of a substitution depending on $\lambda$ is defined without
ambiguity from the definition of convergent sequences of substitutions. If
the boundary of $D$ consists of ordinary values, these integrals exist and
the whole classical apparatus of Cauchy applies. Suppose this condition
holds and that $0$ lies outside $D$. For any integer $k$, positive,
negative, or zero, define
\[
  A^{(k)}=-\frac1{2\pi i}\oint_{\partial D}
  \lambda^{-k}A_\lambda\,d\lambda.
\]
Residue calculus applied to \eqref{eq:79:33}--\eqref{eq:79:34} gives
\[
 AA^{(k)}=A^{(k)}A=A^{(k+1)},
 \qquad
 A^{(k)}A^{(l)}=A^{(k+l)}.
\]

These relations contain many remarkable consequences. For $k=l=1$,
\[
 (A-A^{(1)})A^{(1)}=A^{(1)}(A-A^{(1)})=0;
\]
thus $A^{(1)}$ and $A-A^{(1)}$ are \emph{orthogonal}. More generally, for
any integer $k$ and positive integer $l$, $A^{(k)}$ and
$A^l-A^{(l)}$ are orthogonal.

Multiplying the relations of \hyperref[sec:79]{no.~79} by $A^{(0)}$ and
using the identities above gives
\[
 (E-\lambda A^{(1)})(E+\lambda A^{(0)}A_\lambda)
 =(E+\lambda A_\lambda A^{(0)})(E-\lambda A^{(1)})=E.
\]
Hence every value ordinary for $A$ is ordinary for $A^{(1)}$, with inverse
\[
 E+\lambda A^{(0)}A_\lambda
 =E+\lambda A_\lambda A^{(0)}.
\]
An analogous calculation shows that $\lambda$ is ordinary for
$A-A^{(1)}$ and that the corresponding inverse is
% [Source printed page 119]
\[
 E+\lambda A_\lambda-\lambda A^{(0)}A_\lambda
 =E+\lambda A_\lambda-\lambda A_\lambda A^{(0)}.
\]

For brevity put $B=A^{(1)}$ and $C=A-A^{(1)}$, and introduce $B_\lambda$
and $C_\lambda$ analogously to $A_\lambda$. Then
\[
 B=A^{(0)}A=AA^{(0)},
 \qquad C=(E-A^{(0)})A=A(E-A^{(0)}),
\]
\[
 B_\lambda=A^{(0)}A_\lambda=A_\lambda A^{(0)},
 \qquad
 C_\lambda=(E-A^{(0)})A_\lambda=A_\lambda(E-A^{(0)}).
\]
We already know $BC=CB=0$, and similarly
\[
 BC_\lambda=C_\lambda B=B_\lambda C=CB_\lambda
 =B_\lambda C_\lambda=C_\lambda B_\lambda=0.
\]

Whenever $A_\lambda$ exists, $B_\lambda$ and $C_\lambda$ exist as well.
More is true: $B_\lambda$ exists for every $\lambda$ outside $D$, without
exception, and $C_\lambda$ exists for every $\lambda$ inside $D$, even when
$\lambda$ is singular for $A$. Indeed, residue calculus shows that the
substitutions
\[
 \frac1{2\pi i}\oint_{\partial D}
   \frac{A_\mu\,d\mu}{\lambda-\mu}
 \quad(\lambda\text{ outside}),
 \qquad
 \frac1{2\pi i}\oint_{\partial D}
   \frac{A_\mu\,d\mu}{\mu-\lambda}
 \quad(\lambda\text{ inside})
\]
satisfy the defining equations for $B_\lambda$ and $C_\lambda$,
respectively. So extended, they remain orthogonal to $C$ and $B$.

The importance of the decomposition $A=B+C$ is seen as follows. Let
$\lambda$ be any value inside $D$, and consider the two infinite systems
corresponding to $E-\lambda A$ and $E-\lambda B$, with the same given
element $(x'_k)$. The identities
\[
 E-\lambda A=(E-\lambda B)(E-\lambda C),
 \qquad
 E-\lambda B=(E-\lambda A)(E+\lambda C_\lambda)
\]
show that applying $E-\lambda C$ to a solution of the first system gives a
solution of the second, and applying $E+\lambda C_\lambda$ to a solution of
the second gives one of the first. For the homogeneous systems, the
relations
\[
 E-\lambda B=(E+\lambda C_\lambda)(E-\lambda A),
 \qquad
 E-\lambda A=(E-\lambda C)(E-\lambda B)
\]
% [Source printed page 120]
show that the two systems have exactly the same solutions. These solutions
also satisfy the homogeneous system corresponding to $E-A^{(0)}$, since
\[
 (E-A^{(0)})(E-\lambda B)=E-A^{(0)},
\]
or equivalently $B=A^{(0)}B$. The solutions of this last system are exactly
the elements obtained by applying $A^{(0)}$ to all of Hilbert space, since
$A^{(0)}=A^{(0)}A^{(0)}$. Call them the \emph{principal elements}
corresponding to $D$. Thus, for an interior $\lambda$, the solutions of the
homogeneous system corresponding to $E-\lambda A$ lie among the principal
elements.

The principal elements have another remarkable property: restricted to
their set, $A$ defines a one-to-one correspondence. They are of the form
$A^{(0)}(x_k)$, and since $AA^{(0)}=A^{(0)}A$, their images are of the same
form. Moreover, on this set $A^{(-1)}$ is an inverse of $A$, because
\[
 A^{(-1)}AA^{(0)}=A^{(0)},
 \qquad
 AA^{(-1)}A^{(0)}=A^{(0)}.
\]

We leave this somewhat too general train of ideas, while observing that the
interesting questions it raises are far from exhausted, and add only a few
words about completely continuous $A$. Sums of completely continuous
substitutions, and products having at least one completely continuous
factor, are completely continuous; the relations above therefore show
successively that $A_\lambda$, $A^{(k)}$, $B_\lambda$, and $C_\lambda$ are
completely continuous. In this special case the singular values are
isolated, by \hyperref[sec:70]{no.~70}. If $\lambda_0$ is singular, choose
$D$ to be a small circle about it. Then $B_\lambda$ exists for every
$\lambda\neq\lambda_0$. The substitution $A^{(0)}$, applied to all of
Hilbert space, gives the principal elements corresponding to $\lambda_0$.
But $A^{(0)}$ is completely continuous.

% End inlined .parts/pages_061_120.tex
% Begin inlined .parts/pages_121_180.tex
% [Source printed page 121]
If a completely continuous substitution is applied to Hilbert space, the
resulting elements can be expressed as linear combinations of finitely many of them.
Thus the study of the one-to-one
correspondence defined by \(A\) on the set of principal elements---and hence the
study of \(A_\lambda\) in a neighborhood of \(\lambda_0\)---reduces to that of a
linear substitution in a finite number of variables and is completed by the
classical methods of the theory of linear substitutions.

The same arguments apply also to the transpose \(\mathfrak A\) of \(A\), and the
substitutions \(\mathfrak A_\lambda,\mathfrak A^{(k)},\mathfrak B_\lambda,
\mathfrak C_\lambda\) thereby obtained are the transposes of the substitutions
\(A_\lambda,A^{(k)},B_\lambda,C_\lambda\) which we have just considered.

The method I have just sketched is not confined to the theory presently under
discussion.  A reader versed in the theory of integral equations will have
recognized from the very beginning the parallelism between our arguments and
certain more or less recent investigations belonging to that theory.\footnote{Cf.,
for example, E.~Goursat, ``Recherches sur les équations intégrales linéaires,''
\emph{Annales de la Faculté des Sciences de Toulouse}, vol.~X (1908), pp.~5--98.}
This parallelism is immediately explained by observing that the two theories,
that of integral equations and that of linear substitutions in infinitely many
variables, are special cases of a much more general theory, namely the theory of
distributive operations.\footnote{Cf. on this subject S.~Pincherle's article in
the \emph{Encyclopédie des Sciences mathématiques}, vol.~II, part~5, fascicle~1.}

% [Source printed page 122]
\chapter{The Theory of Quadratic Forms in Infinitely Many Variables}\label{chap:5}

\section{Generalities. Quadratic Forms in a Finite Number of Variables}

\numberedparagraph{83}
In the \hyperref[chap:4]{preceding chapter} we considered, among other things, matrices
\((a_{ik})\) such that \(a_{ik}=\overline{a_{ki}}\).  The corresponding
substitutions displayed certain special properties which became apparent on
studying the Hermitian form
\[
  \sum_{i,k=1}^{\infty} a_{ik}x_i\overline{x_k}.
\]
But it was always the general theory of linear substitutions that we had in
view, and our special matrices played only an auxiliary role in it.  In what
follows we shall confine ourselves to such matrices.  To simplify the
statements, we shall further suppose the \(a_{ik}\) real,\footnote{Throughout
this chapter, the quantities in question will in general be real.  The contrary
hypothesis will always be stated expressly.} and consequently
\(a_{ik}=a_{ki}\); this will lead us to consider the real quadratic forms in
infinitely many variables
\[
  \sum_{i,k=1}^{\infty} a_{ik}x_ix_k.
\]
The results we shall obtain extend almost immediately to Hermitian forms and to
the matrices and substitutions corresponding to them.

% [Source printed page 123]
\numberedparagraph{84}
Our point of departure will be certain well-known facts concerning quadratic
forms in a finite number of variables.  All these facts are contained, in germ,
in the following fundamental theorem.\footnote{Cauchy, \emph{Exercices de
Mathématiques}, fourth year; Paris, 1829, p.~159.}

\begin{theorem}
Every quadratic form in \(n\) variables
\begin{equation}
  A(x,x)=\sum_{i,k=1}^{n}a_{ik}x_ix_k,
  \qquad a_{ik}=a_{ki},
  \tag{1}\label{eq:84:1}
\end{equation}
can be decomposed into \(n\) squares of linear forms,
\begin{equation}
  A(x,x)=\sum_{j=1}^{n}\mu_j
  \left(\sum_{k=1}^{n}l_{jk}x_k\right)^2,
  \tag{2}\label{eq:84:2}
\end{equation}
and this can be done in such a way that
\begin{equation}
  \sum_{j=1}^{n}\left(\sum_{k=1}^{n}l_{jk}x_k\right)^2
  =\sum_{k=1}^{n}x_k^2.
  \tag{3}\label{eq:84:3}
\end{equation}
\end{theorem}

Here is the most important consequence of this theorem.  Let \(A\) denote the
substitution in \(n\) variables corresponding to the matrix \((a_{ik})\), and
let \(E\) denote the identity substitution in \(n\) variables.  Consider the
quadratic form
\begin{equation}
  \sum_{j=1}^{n}
  \frac{\left(\sum_{k=1}^{n}l_{jk}x_k\right)^2}{\lambda-\mu_j}.
  \tag{4}\label{eq:84:4}
\end{equation}
It is well defined for every value of \(\lambda\), except the values
\(\lambda=\mu_j\).  The linear substitution corresponding to it is precisely
the inverse of \(\lambda E-A\).\footnote{Considering \(\lambda E-A\), instead
of \(E-\lambda A\), amounts to replacing \(\lambda\) by \(1/\lambda\).  This
modification is not essential; it merely conforms better to the order of ideas
we shall follow.  Among other things, it has the advantage that
\(\lambda=\infty\) becomes an ordinary value.}

Indeed, formulas \eqref{eq:84:2} and \eqref{eq:84:3} may be interpreted as
follows.  Let \(E_j\) be the substitution corresponding to the quadratic form
\(\bigl(\sum_{k=1}^{n}l_{jk}x_k\bigr)^2\).  Since, by
\eqref{eq:84:3}, the matrix \((l_{jk})\)
% [Source printed page 124]
is orthogonal and normalized, we plainly have \(E_j^2=E_j\) and
\(E_iE_j=0\) for \(i\ne j\); moreover,
\(E_1+E_2+\cdots+E_n=E\).  By \eqref{eq:84:2},
\[
  A=\mu_1E_1+\cdots+\mu_nE_n,
\]
and hence
\[
  \lambda E-A=(\lambda-\mu_1)E_1+\cdots+(\lambda-\mu_n)E_n.
\]
On the other hand, the substitution corresponding to the quadratic form
\eqref{eq:84:4} is
\[
  \frac{E_1}{\lambda-\mu_1}+\cdots+
  \frac{E_n}{\lambda-\mu_n}.
\]
The product of these two substitutions, which plainly commute, is
\(E_1+\cdots+E_n=E\).  They are therefore inverses of one another.

This argument does not suppose \(\lambda\) real.  Thus \eqref{eq:84:4} gives
the expression for \((\lambda E-A)^{-1}\), or more precisely for the
corresponding quadratic form, for every ordinary value of the complex parameter
\(\lambda\).  It also exhibits the singular values
\(\lambda=\mu_1,\mu_2,\ldots,\mu_n\), and at the same time shows how the
substitution \((\lambda E-A)^{-1}\) behaves in a neighborhood of those singular
values.

In short, in the case of finitely many variables, the decomposition
\eqref{eq:84:2} of the form \eqref{eq:84:1} gives at once the solution of every
problem concerning the inverse of \(\lambda E-A\).  Will the same be true for
infinitely many variables?  Yes.  In the memoir already mentioned many times,
Hilbert proved that every real bounded quadratic form in infinitely many
variables can be decomposed in a manner analogous to \eqref{eq:84:2} and
\eqref{eq:84:3}; only, in general, finite sums must be replaced by two kinds of
infinite processes.  On the one hand, one must sum a countable infinity of terms
analogous to those in \eqref{eq:84:2} and \eqref{eq:84:3}; on the other, one
must integrate with respect to a continuous parameter.  I observed elsewhere
that these two procedures may be united by using the notion of the Stieltjes
integral, and that this way of stating Hilbert's theorem offers some advantages
for the proof.\footnote{See my letter to Hilbert, ``Ueber quadratische Formen von
unendlich vielen Veränderlichen,'' \emph{Nachrichten der Gesellschaft der
Wissenschaften zu Göttingen} (1910), pp.~190--195.}

\numberedparagraph{85}
Before passing finally to quadratic forms in infinitely many variables, let us
return once more to the form \eqref{eq:84:1}
% [Source printed page 125]
and to the corresponding substitution \(A\).  Let \(A^2,A^3,\ldots\) be the
powers of this substitution.  In view of the formulas
\[
  E_iE_j=\begin{cases}E_i,&i=j,\\0,&i\ne j,
  \end{cases}
\]
the corresponding quadratic forms are plainly
\[
  \sum_{j=1}^{n}\mu_j^2\left(\sum_{k=1}^{n}l_{jk}x_k\right)^2,
  \qquad
  \sum_{j=1}^{n}\mu_j^3\left(\sum_{k=1}^{n}l_{jk}x_k\right)^2,
  \quad\ldots.
\]
Let, further, \(\nu_0,\nu_1,\ldots,\nu_r\) be arbitrary real numbers, and set
\(\nu_0+\nu_1\mu+\cdots+\nu_r\mu^r=f(\mu)\).  Consider the substitution
\(\nu_0E+\nu_1A+\cdots+\nu_rA^r\), which, by an obvious symbolic notation,
may be denoted by \(f(A)\).  The corresponding quadratic form is
\begin{equation}
  \sum_{j=1}^{n}f(\mu_j)
  \left(\sum_{k=1}^{n}l_{jk}x_k\right)^2.
  \tag{5}\label{eq:85:5}
\end{equation}
By \eqref{eq:84:3}, the value of this expression is a certain mean of the
values \(f(\mu_1),\ldots,f(\mu_n)\), multiplied by
\(\sum_{k=1}^{n}x_k^2\).  On the other hand, the smallest and largest of the
quantities \(\mu_j\) are respectively equal to the minimum \(m\) and maximum
\(M\) of the form \eqref{eq:84:1}, varied subject to
\(\sum_{k=1}^{n}x_k^2=1\); this follows immediately from our decomposition
formulas.  Consequently, \emph{the set of values of the form
\eqref{eq:85:5}, varied under the same condition, lies entirely between the two
 extreme values of \(f(\mu)\) on the interval \([m,M]\).\translatornote{The French source uses \((m,M)\), here and throughout this chapter, for an interval including its endpoints.  Modern bracket notation is used in the translation.}}

\section{Quadratic Forms in Infinitely Many Variables. Symbolic Sequences and Products}

\numberedparagraph{86}
Let \((a_{ik})\) be a doubly indexed matrix, real and symmetric in \(i\) and
\(k\).  Consider the quadratic form
\begin{equation}
  A(x,x)=\sum_{i,k=1}^{\infty}a_{ik}x_ix_k
  \tag{6}\label{eq:86:6}
\end{equation}
and its sections
\[
  [A(x,x)]_n=\sum_{i,k=1}^{n}a_{ik}x_ix_k.
\]

% [Source printed page 126]
I assume that there exists a positive number \(M\) such that
\[
  \lvert[A(x,x)]_n\rvert\le M\sum_{k=1}^{n}x_k^2,
\]
whatever \(n\) and the \(x_k\) may be.  When this hypothesis is satisfied, we
say, with Hilbert, that the form \eqref{eq:86:6} is \emph{bounded}.\footnote{It
should be observed that, in presenting his investigations, Hilbert begins from
the most general possible point of view.  Without making any assumption on the
orders of magnitude of the coefficients \(a_{ik}\), he regards the form
\eqref{eq:86:6} as a symbol for the collection of its sections.  Cf. also the
note by J.~Le Roux, ``Sur les formes quadratiques définies à une infinité de
variables,'' \emph{Comptes rendus}, January 10, 1910, whose point of view is less
general but more accessible to calculation.}

Associate with the quadratic form \eqref{eq:86:6} the bilinear form
\[
  A(x,y)=\sum_{i,k=1}^{\infty}a_{ik}x_iy_k
\]
and its sections
\[
  [A(x,y)]_n=\sum_{i,k=1}^{n}a_{ik}x_iy_k.
\]
Plainly,
\begin{align*}
  4\lvert[A(x,y)]_n\rvert
  &=\left\lvert[A(x+y,x+y)]_n-[A(x-y,x-y)]_n\right\rvert\\
  &\le M\left[\sum_{k=1}^{n}(x_k+y_k)^2+
                    \sum_{k=1}^{n}(x_k-y_k)^2\right]\\
  &=2M\sum_{k=1}^{n}(x_k^2+y_k^2).
\end{align*}
Consequently, when \(\sum_{k=1}^{n}x_k^2\le1\) and
\(\sum_{k=1}^{n}y_k^2\le1\), we also have
\[
  \lvert[A(x,y)]_n\rvert\le M.
\]
Hence, by homogeneity, in the general case,
\[
  \lvert[A(x,y)]_n\rvert\le
  M\left(\sum_{k=1}^{n}x_k^2\sum_{i=1}^{n}y_i^2\right)^{1/2}.
\]

% [Source printed page 127]
Put, in particular,
\[
  y_k=\sum_{i=1}^{n}a_{ik}x_i;
\]
then
\[
  \sum_{i=1}^{n}\left(\sum_{k=1}^{n}a_{ik}x_k\right)^2
  \le M^2\sum_{k=1}^{n}x_k^2.
\]
That is, under the hypothesis made, the matrix \((a_{ik})\) corresponds to a
linear substitution \(A\), and \(M_A\le M\).  It also follows, by
\hyperref[sec:58]{no.~58}, that for every element \((x_k)\) of Hilbert space the
doubly indexed series \eqref{eq:86:6} converges to a determinate limit,
independent of the manner in which the limit is taken.  In particular,
\[
  [A(x,x)]_n\longrightarrow A(x,x).
\]

Conversely, let \(A\) be a linear substitution in infinitely many variables,
with coefficients \(a_{ik}\) real and symmetric in \(i\) and \(k\).  Let \(M_A\)
be the bound of \(A\).  By definition,
\[
  \sum_{i=1}^{\infty}\left(\sum_{k=1}^{n}a_{ik}x_k\right)^2
  \le M_A^2\sum_{k=1}^{n}x_k^2,
\]
and the Cauchy--Lagrange inequality gives
\[
  \lvert[A(x,x)]_n\rvert
  \le\left[\sum_{i=1}^{n}\left(\sum_{k=1}^{n}a_{ik}x_k\right)^2\right]^{1/2}
      \left(\sum_{k=1}^{n}x_k^2\right)^{1/2}
  \le M_A\sum_{k=1}^{n}x_k^2.
\]
Thus the form \eqref{eq:86:6} is bounded.

In summary, for a given real symmetric matrix \((a_{ik})\), the following two
hypotheses are equivalent: (1) it corresponds to a bounded quadratic form
\(A(x,x)\); (2) it corresponds to a linear substitution \(A\).  The bound
\(M_A\) of the substitution is at the same time the upper bound of
\(\lvert A(x,x)\rvert\), varied subject to
\(\sum_{k=1}^{\infty}x_k^2=1\).

\numberedparagraph{87}
Consider now a sequence \(A_1,A_2,\ldots\) of bounded quadratic forms, and
suppose that the corresponding substitutions tend to a substitution \(A\)
(\hyperref[sec:71]{no.~71}).  Under these conditions, by
\hyperref[sec:40]{no.~40}, we shall also have
\begin{equation}
  A_n(x,y)\longrightarrow A(x,y),
  \tag{7}\label{eq:87:7}
\end{equation}
% [Source printed page 128]
and in particular
\[
  A_n(x,x)\longrightarrow A(x,x).
\]
Conversely, we may suppose that this last relation holds for every element
\((x_k)\).  Then, a fortiori, the sections, and hence the coefficients, of the
forms \(A_n\) also converge to those of \(A\).  Consequently, if the quantities
\(M_{A_1},M_{A_2},\ldots\) are also bounded as a whole, we may assert that
\emph{the substitutions \(A_1,A_2,\ldots\) converge to the substitution \(A\).}

Let \(A(x,x)\) and \(B(x,x)\) be two bounded quadratic forms, and let \(A\) and
\(B\) be their corresponding substitutions.  Their product \(AB\) need not
correspond to a quadratic form, since its coefficient matrix need not be
symmetric in \(i\) and \(k\).  It is symmetric if and only if \(AB=BA\).  In
that case we shall call the quadratic form
\[
  AB(x,x)=BA(x,x)=\sum_{i=1}^{\infty}
  \left(\sum_{k=1}^{\infty}a_{ik}x_k\right)
  \left(\sum_{k=1}^{\infty}b_{ik}x_k\right)
\]
the \emph{symbolic product} of the forms \(A\) and \(B\).

We shall need the following fact: if
\[
  A_n(x,x)\longrightarrow A(x,x)
\]
and, moreover, the quantities \(M_{A_1},M_{A_2},\ldots\) are bounded as a whole,
then
\begin{equation}
  A_nB(x,x)\longrightarrow AB(x,x).
  \tag{8}\label{eq:87:8}
\end{equation}
This fact could be connected with the arguments of
\hyperref[sec:71]{nos.~71} and following; it also follows from
\eqref{eq:87:7} on setting \((y_k)=B(x_k)\).

\section{Symbolic Functions of a Quadratic Form. The Correspondence Between Functions of One Variable and Symbolic Functions}

\numberedparagraph{88}
Let \(A\) be a linear substitution that is both real and symmetric.  Its powers
\(A^2,A^3,\ldots\) will have the same properties, as will the polynomials
\(f(A)=\nu_0E+\nu_1A+\cdots+\nu_rA^r\).

% [Source printed page 129]
Each of the substitutions \(E,A,A^2,\ldots,f(A)\) corresponds to a bounded
quadratic form \(E(x,x),A(x,x),A^2(x,x),\ldots,f(A)(x,x)\).  For brevity we
shall speak of the quadratic forms \(E,A,A^2,\ldots,f(A)\).

Let \(A_n\) be the \(n\)-th section of the substitution \(A\); it may be regarded
either as a substitution in infinitely many variables or as a substitution in
\(n\) variables.  It corresponds to a quadratic form in \(n\) variables which
is simply the \(n\)-th section of the form \(A(x,x)\).  Likewise the powers
\(A_n^2,A_n^3,\ldots\) and the polynomials \(f(A_n)\) each give a quadratic form
in \(n\) variables: \(A_n^2(x,x),A_n^3(x,x),\ldots,f(A_n)(x,x)\).  By
\hyperref[sec:75]{no.~75}, for every element \((x_k)\),
\begin{equation}
  f(A_n)(x,x)\longrightarrow f(A)(x,x).
  \tag{9}\label{eq:88:9}
\end{equation}
This relation enables us to extend the theorem of \hyperref[sec:85]{no.~85} to
quadratic forms in infinitely many variables.  Indeed, let \(m_n\) and \(M_n\)
be the minimum and maximum of the form \(A_n\), varied subject to
\(E_n=\sum_{k=1}^{n}x_k^2=1\); they plainly lie between the lower bound \(m\)
and upper bound \(M\) of the form \(A\), varied subject to \(E=1\).  By the
theorem to be extended, as \(\mu\) ranges from \(m_n\) to \(M_n\),
\(f(A_n)(x,x)\) lies between the extreme values of \(f(\mu)E_n(x,x)\).  Passing
to the limit and using \eqref{eq:88:9}, we see that:

\emph{The set of values of the form \(f(A)\), varied subject to \(E=1\), lies
between the two extreme values of \(f(\mu)\), with \(\mu\) ranging over the
 interval \([m,M]\).}

In particular, if \(f(\mu)>0\) throughout the interval \([m,M]\), then \(f(A)\)
is a positive quadratic form.

\numberedparagraph{89}
We have just made a uniquely determined quadratic form \(f(A)\) correspond to
each polynomial \(f(\mu)\).  This correspondence has the following properties:

\begin{enumerate}[label=\Roman*.]
\item It is \emph{distributive}: if \(\varphi=c_1f_1+c_2f_2\), then
\(\varphi(A)=c_1f_1(A)+c_2f_2(A)\).
\item The set of values of \(f(A)\), when \(E=1\), lies between the extreme
values of \(f(\mu)\).
% [Source printed page 130]
\item It is multiplicative: the ordinary product of two functions \(f_1,f_2\)
corresponds to the symbolic product of the forms \(f_1(A),f_2(A)\).
\end{enumerate}

We shall see that the correspondence may be extended to a much wider field of
functions in such a way that all three stated properties remain valid.

The passage from polynomials to continuous functions is immediate.  By the
well-known theorem of Weierstrass, every function continuous on \([m,M]\) can be
uniformly approximated there by a sequence of polynomials.  Every such sequence
corresponds to a sequence of quadratic forms; by what precedes, this sequence
also converges uniformly.  It therefore has a limiting form, which we make
correspond to the given continuous function.  It is at once clear that the
limiting form does not depend on the particular sequence chosen to approximate
the given function.  It is equally easy to see that the stated properties remain
valid for the correspondence so established.

For the arguments that follow, however, continuous functions are not enough.
We shall extend the correspondence to a field containing, in addition to
continuous functions, certain discontinuous functions.  It is clear that to do
so we must relinquish the exclusive use of uniformly convergent
sequences.\footnote{The analogous idea of considering the correspondence between
functions of one or several variables and symbolic functions of certain
functional operations was recently applied successfully by V.~Volterra to the
study of integral equations and of a mixed type of functional equations
(integro-differential equations), in a series of notes published in the
\emph{Atti della Reale Accademia dei Lincei}.  See also the same author's
\emph{Leçons sur les équations intégrales et les équations
intégro-différentielles}, recently published in this Collection.  It should be
observed, however, that Volterra confines himself to analytic functions; here,
on the contrary, discontinuous functions are essential.}

\numberedparagraph{90}
Consider a sequence of polynomials \([f_n(\mu)]\).  Suppose that
\(f_1(\mu)\le f_2(\mu)\le f_3(\mu)\le\cdots\) on the interval \([m,M]\); we
express this by saying that the sequence is increasing.  Suppose also that it is
bounded.  Then the sequence tends
% [Source printed page 131]
on \([m,M]\) to a bounded function \(f(\mu)\).  On the other hand, by I and II,
the forms \(f_n(A)\) also constitute an increasing bounded sequence; they remain
\(\le E\) times the upper bound of \(f(\mu)\).  Consequently, the \(f_n(A)\)
tend to a limiting form, which also remains \(\le E\) times that upper bound.

Let us denote the limiting form by \(f(A)\) and attach it to the function
\(f(\mu)\).  To justify this convention we must prove that the limiting form
depends only on the limiting function.  More explicitly, when two increasing
sequences of polynomials \([f_n(\mu)]\), \([g_n(\mu)]\) tend to the same limiting
function, we must prove that the forms \(f_n(A)\), \(g_n(A)\) also tend to the
same limiting form.

The fact to be proved is contained in the following slightly more comprehensive
result.  Let \([f_n(\mu)]\), \([g_n(\mu)]\) be two increasing sequences, with
limiting functions \(f(\mu)\) and \(g(\mu)\), and suppose moreover that
\(f(\mu)\le g(\mu)\).  I assert that \(f(A)\le g(A)\).  Indeed, let
\(f_m(\mu)\) be any polynomial in the first sequence and let \(\varepsilon>0\)
be arbitrary.  Hold \(m\) and \(\varepsilon\) fixed and run through the sequence
\([g_n(\mu)]\).  For sufficiently large \(n\) we shall have
\(g_n(\mu)>f_m(\mu)-\varepsilon\).  Otherwise, the values \(\mu\) satisfying
\[
 f_m(\mu)-\varepsilon\ge g_1(\mu),\quad
 f_m(\mu)-\varepsilon\ge g_2(\mu),\quad\ldots
\]
would form a sequence of closed sets, each containing those that follow, and at
least one point \(\mu^*\) would belong to all of them.  Then
\(f(\mu^*)>f_m(\mu^*)-\varepsilon\ge g(\mu^*)\), contrary to the hypothesis.
Thus for all sufficiently large \(n\), \(f_m-\varepsilon<g_n\), and therefore
\(f_m(A)-\varepsilon E<g_n(A)\).  This implies
\(f_m(A)-\varepsilon E<g(A)\), and since \(m\) and \(\varepsilon\) are
arbitrary, \(f(A)\le g(A)\), as was to be proved.

The fact that the limiting form depends only on the limiting function is the
case \(f=g\); it follows from the result just established by replacing the
equality \(f=g\) by the two inequalities \(f\le g\) and \(g\le f\).

Thus the correspondence between \(f(\mu)\) and \(f(A)\), first established for
polynomials, extends uniquely to every bounded function that is the limit of an
increasing sequence of polynomials.  Plainly it may also be extended to limits
of decreasing sequences; one need only
% [Source printed page 132]
change signs.  There is more.  Let \(f\) and \(g\) be two functions of the type
just studied.  Then \(h=f+g\) is of the same type and plainly
\(h(A)=f(A)+g(A)\).  In contrast, \(\varphi=f-g\) will in general be neither of
the same nor of the opposite type.  Define
\(\varphi(A)=f(A)-g(A)\).  To justify this convention, we must show that if
\(f-g=f^*-g^*\), then also
\[
  f(A)-g(A)=f^*(A)-g^*(A).
\]
The hypothesis may be written \(f+g^*=f^*+g\); both sides now belong to the type
just studied.  Therefore
\[
  f(A)+g^*(A)=f^*(A)+g(A),
\]
and consequently
\[
  f(A)-g(A)=f^*(A)-g^*(A).
\]
The convention is legitimate, and our correspondence is thus extended to
functions of the type \(f-g\).

The extended correspondence still has properties I, II, and III.  First, it is
plainly distributive.  It also has property II: \emph{the set of values of the
form \(\varphi(A)=f(A)-g(A)\), for \(E(x,x)=1\), lies between the lower and
upper bounds of \(\varphi(\mu)\).}  For example, if \(c\) is the lower bound,
then \(f(\mu)\ge g(\mu)+c\), and hence \(f(A)\ge g(A)+cE\), that is,
\(\varphi(A)\ge cE\).  The argument for the upper bound is the same.

Finally, the extended correspondence also has property III: it is
\emph{multiplicative}.  Without loss of generality it suffices to prove this for
two functions \(f(\mu)\), \(g(\mu)\) which are the limits of two increasing
sequences of polynomials \([f_n(\mu)]\), \([g_n(\mu)]\).  We may moreover suppose
that \(f,g,f_n,g_n\) are positive, which amounts to adding suitable constants.

\numberedparagraph{91}
Consider the sequence of functions \([f(\mu)g_n(\mu)]\); it increases to the
function \(f(\mu)g(\mu)\).  Since all these functions are limits of increasing
sequences of polynomials, there correspond to them uniquely determined forms,
which we
% [Source printed page 133]
denote by \(fg_n(A)\) and \(fg(A)\).  I assert that
\begin{equation}
  fg_n(A)\longrightarrow fg(A).
  \tag{9}\label{eq:91:9}
\end{equation}
Indeed, let \(p(\mu)\) be any polynomial with
\(p(\mu)<f(\mu)g(\mu)\).  The points where
\(f_m(\mu)g_n(\mu)\le p(\mu)\) form a closed set \(F_{m,n}\).  The set \(F_n\)
of points where \(f(\mu)g_n(\mu)\le p(\mu)\), being the limiting set of the
closed sets \(F_{1,n},F_{2,n},\ldots\), each of which contains the following
ones, is also closed.  Finally, each of the closed sets
\(F_1,F_2,\ldots\) contains those that follow; hence from some rank onward the
sets \(F_n\) must be empty.  Otherwise they would have at least one common point
\(\mu^*\), and we would have
\(f(\mu^*)g(\mu^*)\le p(\mu^*)\), contrary to the hypothesis.  Thus for
sufficiently large \(n\), \(f(\mu)g_n(\mu)>p(\mu)\), whence
\(fg_n(A)\ge p(A)\), and consequently
\begin{equation}
  \lim_{n\to\infty}fg_n(A)\ge p(A).
  \tag{10}\label{eq:91:10}
\end{equation}
Let \([p_n(\mu)]\) now be an increasing sequence of polynomials tending to
\(fg\) (for example, \(p_n=f_ng_n\)).  The preceding argument applies to all
polynomials \(p_n(\mu)-1/n\), and hence
\begin{equation}
  \lim_{n\to\infty}fg_n(A)
  \ge\lim_{n\to\infty}\left[p_n(A)-\frac1nE\right]=fg(A).
  \tag{11}\label{eq:91:11}
\end{equation}
On the other hand, for every \(n\),
\(fg_n(\mu)\le fg(\mu)\), and consequently
\(fg_n(A)\le fg(A)\); hence
\begin{equation}
  \lim_{n\to\infty}fg_n(A)\le fg(A).
  \tag{12}\label{eq:91:12}
\end{equation}
Comparing \eqref{eq:91:11} and \eqref{eq:91:12} gives the asserted result
\eqref{eq:91:9}.

With this result established, the argument is soon completed.  Since the
correspondence is multiplicative for polynomials,
\[
  f_m(A)g_n(A)=f_mg_n(A),
\]
and therefore, by \eqref{eq:87:8},
\[
  f(A)g_n(A)=fg_n(A).
\]

% [Source printed page 134]
The same formula \eqref{eq:87:8} gives
\[
  f(A)g_n(A)\longrightarrow f(A)g(A).
\]
Consequently,
\[
  fg_n(A)\longrightarrow f(A)g(A).
\]
Comparison with \eqref{eq:91:9} yields
\[
  fg(A)=f(A)g(A),
\]
that is, the correspondence is multiplicative.

Let us add one observation that will be used later.  The argument by which we
established \eqref{eq:91:9} also implies the more general proposition:
\emph{if a sequence of functions of type \(f\) (that is, bounded functions that
are limits of increasing sequences of polynomials) increases to a function of
the same type, the corresponding forms tend to the form corresponding to the
limiting function.}

\numberedparagraph{92}
It would remain to characterize the enlarged field of functions more precisely.
How can one decide whether a given function belongs to it?  We need not dwell on
this interesting question, which is closely connected with the study of the
so-called semicontinuous functions.  For the applications that follow, a few
observations suffice.  I assert that our field contains: (1) all continuous
functions; (2) all bounded functions that are limits of increasing or decreasing
sequences of continuous functions.  The first are simultaneously of type \(f\)
and type \(-f\); limits of increasing sequences are of type \(f\), and limits of
decreasing sequences are of type \(-f\).  Indeed, let \(f(\mu)\) be continuous,
and choose a polynomial \(p_n(\mu)\) which approximates
\(f(\mu)-2^{-n}\) to within \(2^{-n-2}\).  Then \([p_n(\mu)]\) increases to
\(f(\mu)\).  If a function is the limit of an increasing sequence of continuous
functions \(f_n(\mu)\), choose \(p_n(\mu)\) so that it approximates the
continuous function \(f_n(\mu)-2^{-n}\) to within \(2^{-n-2}\).

% [Source printed page 135]
\section{Application of the Correspondence to the Calculation of
\texorpdfstring{\((\lambda E-A)^{-1}\)}{(lambda E - A) inverse} and to the
Study of the Spectrum}

\numberedparagraph{93}
The correspondence between the functions \(f(\mu)\) and the quadratic forms
\(f(A)\) which we have just established immediately supplies the inverse of
\(\lambda E-A\), and also other related quadratic forms called \emph{spectral
forms}.

Indeed, let \(\lambda\) be a real number outside the interval \([m,M]\).  The
function \(f(\mu)=1/(\lambda-\mu)\) is continuous on \([m,M]\), and hence a form
\(f(A)\) corresponds to it.  The function \(g(\mu)=\lambda-\mu\) corresponds to
the form \(\lambda E-A\), and the product \(f(\mu)g(\mu)=1\) corresponds to the
form \(E\).  Consequently,
\[
  f(A)(\lambda E-A)=(\lambda E-A)f(A)=E;
\]
that is, the substitution \(f(A)\) is the inverse of \(\lambda E-A\).

For complex \(\lambda\), separate in \(f\) and \(g\) the real from the purely
imaginary parts.  If \(f_1(\mu)+if_2(\mu)\) is the decomposition of
\(f(\mu)=1/(\lambda-\mu)\), define
\(f(A)=f_1(A)+if_2(A)\).  The substitution so defined is the inverse of
\(\lambda E-A\).

In summary, the correspondence between continuous functions \(f(\mu)\) and
forms \(f(A)\) implies the existence of an inverse \((\lambda E-A)^{-1}\) for
every value of \(\lambda\), except possibly the real values between \(m\) and
\(M\).  Moreover, every sequence of continuous functions \(f_n(\mu)\) tending
to \(1/(\lambda-\mu)\) uniformly, or in such a way that
\(\sum\lvert f_{n+1}(\mu)-f_n(\mu)\rvert\) is bounded, furnishes a procedure for
calculating \((\lambda E-A)^{-1}\).  Thus, in particular, for sufficiently large
\(\lvert\lambda\rvert\),
\[
  (\lambda E-A)^{-1}=\frac1\lambda E+\frac1{\lambda^2}A+
  \frac1{\lambda^3}A^2+\cdots,
\]
in accordance with \hyperref[sec:73]{no.~73}.

\numberedparagraph{94}
The use of discontinuous functions will enable us to carry the study of our
correspondence farther.  Let \(f(\mu)\)
% [Source printed page 136]
be continuous on the interval \([m,M]\), endpoints included.  We may suppose it
defined from \(-\infty\) to \(+\infty\), completing the definition, for example,
by putting \(f(\mu)=f(m)\) for \(\mu<m\) and \(f(\mu)=f(M)\) for \(\mu>M\).
One could remain within \([m,M]\), but the notation would then have certain
disadvantages.  In any event it would suffice to use, in place of
\(( -\infty,+\infty)\), a finite interval \([m-\varepsilon,M+\varepsilon]\),
where \(\varepsilon>0\) is arbitrary.  For definiteness we shall use
\(( -\infty,+\infty)\).

Partition \(( -\infty,+\infty)\) into partial intervals \(I_k\), with \(k\)
ranging from \(-\infty\) to \(+\infty\), so that the oscillation of \(f(\mu)\)
on each interval is at most \(\omega\).  Let \(\mu_k\) be an arbitrary point of
\(I_k\).  Define the discontinuous function \(f'\) by setting
\(f'=f(\mu_k)\) in the interior and at the left endpoint \(\xi_k\) of \(I_k\),
for every \(k\).  Since \(f'\) is the difference of two functions, each of
which is the limit of an increasing sequence of continuous functions, the
extended correspondence attaches to it a uniquely determined form
\(f'(A)\).\translatornote{The French original says that \(f'\) itself is the
limit of an increasing sequence of continuous functions.  That need not hold
at an upward jump; the difference-of-two statement is the one justified by the
preceding extension of the correspondence.}  Introduce also the functions \(f_\xi(\mu)\), one for
each point \(\xi\), defined by \(f_\xi(\mu)=1\) for \(\mu<\xi\), and
\(f_\xi(\mu)=0\) for \(\mu\ge\xi\).  To simplify the notation put
\(f_\xi(A)=A_\xi\).  For \(\xi\le m\), \(A_\xi=0\); for \(\xi>M\),
\(A_\xi=E\).

Plainly,
\[
  f'(\mu)=\sum_k f(\mu_k)
  [f_{\xi_{k+1}}(\mu)-f_{\xi_k}(\mu)],
\]
and consequently
\begin{equation}
  f'(A)=\sum_k f(\mu_k)(A_{\xi_{k+1}}-A_{\xi_k}).
  \tag{13}\label{eq:94:13}
\end{equation}
Since \(f-f'\) lies between \(-\omega\) and \(\omega\), the form
\(f(A)-f'(A)\) lies between \(-\omega E\) and \(\omega E\).  Thus, in
calculating \(f'(A)\), we have also calculated \(f(A)\) to within \(\omega E\).
In other words, the right-hand side of \eqref{eq:94:13} is an approximation to
\(f(A)\), with error at most \(\omega E\).  Now \(\omega\) may be made
arbitrarily small by choosing the lengths of the intervals \(I_k\) sufficiently
small.  Passing to the limit gives
% [Source printed page 137]
\begin{equation}
  f(A)=\int_{-\infty}^{+\infty}f(\xi)\,dA_\xi.
  \tag{14}\label{eq:94:14}
\end{equation}
By this last symbol we mean the limit of \eqref{eq:94:13} as \(\omega\to0\).
It is not an integral in the ordinary sense, but the manner in which it has just
been defined is entirely analogous to Riemann's classical definition.  The form
\(A_\xi(x,x)\) is regarded as a function of the parameter \(\xi\), and our
integral \eqref{eq:94:14} is of the type \(\int f(\xi)\,d\alpha(\xi)\), first
considered by Stieltjes in his celebrated memoir \emph{Recherches sur les
fonctions discontinues}.\footnote{\emph{Académie des Sciences: Mémoires présentés
par divers savants}, vol.~XXXII, no.~2; \emph{Annales de la Faculté des Sciences
de Toulouse}, vol.~VIII (1894), pp.~1--122; vol.~IX (1895), pp.~1--47.}

A few remarks will suffice to show the nature of the Stieltjes integral.  If
\(\alpha(\xi)=\xi\) between \(a\) and \(b\), and is constant elsewhere, we
recover \(\int_a^b f(\xi)\,d\xi\).  If \(f(\xi)\) and \(\alpha(\xi)\) are
continuous and \(\alpha\) has a derivative \(\alpha'(\xi)\) of which it is an
indefinite integral, then \(\int f(\xi)\,d\alpha(\xi)\) immediately becomes
\(\int f(\xi)\alpha'(\xi)\,d\xi\).  The great advantage of Stieltjes's notion is
that it also applies to functions \(\alpha(\xi)\) that are not indefinite
integrals, and even accommodates certain discontinuous functions.  For example,
given a countable infinity of distinct values \(\xi_k\) and corresponding
values \(a_k\) such that \(\sum\lvert a_k\rvert\) converges, define
\[
  \alpha(\xi)=\sum a_k,
\]
where the sum extends over all \(k\) such that \(\xi_k<\xi\).  Under this
hypothesis the integral
\[
  \int_{-\infty}^{\infty}f(\xi)\,d\alpha(\xi)
\]
exists for every continuous bounded function \(f(\xi)\), and equals
% [Source printed page 138]
\[
  \sum_k a_kf(\xi_k).
\]

\numberedparagraph{95}
Return to formula \eqref{eq:94:14}.  It gives the analytic expression, analogous
to \eqref{eq:85:5}, for every form \(f(A)\) corresponding to a continuous
function \(f(\mu)\).\footnote{It should be observed here that the correspondence
under consideration belongs to the class of linear operations whose analytic
expression as limits of integrals was given by J.~Hadamard, ``Sur les opérations
fonctionnelles,'' \emph{Comptes rendus}, February 9, 1903, and which I succeeded
in expressing by a single integral by using the Stieltjes integral, ``Sur les
opérations fonctionnelles linéaires,'' \emph{Comptes rendus}, November 29, 1909.
Cf. also the more detailed memoir, ``Sur certains systèmes singuliers d'équations
intégrales,'' \emph{Annales de l'École Normale} (1911), pp.~33--62, and, in
particular for the correspondence considered here, my note cited on p.~124.

Formula \eqref{eq:94:14} could easily be extended to forms corresponding to
discontinuous functions; one need only modify Stieltjes's definition to a
greater or lesser extent.  Cf. also H.~Lebesgue, ``Sur les intégrales de
Stieltjes et les opérations fonctionnelles linéaires,'' \emph{Comptes rendus},
January 12, 1910.}\nobreak\hspace{0.12em}\translatornote{The French original prints pp.~31--62 for
Riesz's ``Sur certains systèmes singuliers d'équations intégrales''; the
article begins on p.~33.}  In particular, putting
\(f(\mu)=1/(\lambda-\mu)\), we obtain
\begin{equation}
  (\lambda E-A)^{-1}=\int_{-\infty}^{\infty}
  \frac{dA_\xi}{\lambda-\xi}.
  \tag{15}\label{eq:95:15}
\end{equation}
This formula applies to every value of \(\lambda\), except perhaps real values.
By means of the convention we shall now make, it also applies to every real
ordinary value, that is, every value for which the inverse
\((\lambda E-A)^{-1}\) exists.  In the approximating expression
\eqref{eq:94:13} for the integral \eqref{eq:94:14}, agree to omit terms for
which \(A_{\xi_{k+1}}-A_{\xi_k}\) vanishes.  For functions \(f(\mu)\) which do
not become infinite, the meaning of the integral is unchanged.  With this
convention, it is first easy to see that \eqref{eq:95:15} applies for every
\(\lambda\) outside \([m,M]\): near \(\xi=\lambda\), \(A_\xi\) remains constant.
I assert that \eqref{eq:95:15} holds for every ordinary value of \(\lambda\).
Suppose, indeed, that the nondecreasing function \(A_\xi\) does not remain
constant on any interval containing \(\lambda\); in particular it will not be
constant on \(I=(\lambda-h,\lambda+h)\).  Put \(f(\mu)=1\) in the interior and
at the left endpoint of \(I\), and zero elsewhere.  Plainly \(f^2(\mu)=f(\mu)\),
and hence
% [Source printed page 139]
\(f(A)f(A)=f(A)\).  By hypothesis, the nonnegative form
\(f(A)=A_{\lambda+h}-A_{\lambda-h}\) does not vanish identically; that is, it is
positive for at least one system \((x_k)\).  Let \((y_k)\) be the system obtained
from \((x_k)\) by applying the substitution \(f(A)\).  Then
\(E(y,y)=f(A)(x,x)>0\), so \((y_k)\ne(0)\).  Moreover, applying \(f(A)\) to
\((y_k)\) reproduces it.  Thus applying \((\lambda E-A)^2\) to \((y_k)\) is
equivalent to applying \((\lambda E-A)^2f(A)\).  The latter substitution
corresponds to the function \(g(\mu)=(\lambda-\mu)^2f(\mu)\), whose values lie
between zero and \(h^2\).  Hence, if \((y'_k)\) denotes what \((y_k)\) becomes
after applying \(\lambda E-A\), then
\begin{equation}
  \frac{\sum_{k=1}^{\infty}{y'_k}^{2}}
       {\sum_{k=1}^{\infty}y_k^2}
  =\frac{(\lambda E-A)^2(y,y)}{\sum_{k=1}^{\infty}y_k^2}
  \le\frac{h^2E(y,y)}{\sum_{k=1}^{\infty}y_k^2}=h^2.
  \tag{16}\label{eq:95:16}
\end{equation}
But \(h^2\) may be chosen arbitrarily small; under the hypothesis made there is
always a corresponding system \((y_k)\).  Thus, by a suitable choice of
\((y_k)\), the left-hand side of \eqref{eq:95:16} can be made arbitrarily small.
Consequently, under this hypothesis \(\lambda E-A\) has no inverse; the value
\(\lambda\) is singular.

Following Hilbert, call the \emph{spectrum} of the form \(A\) the set of values
\(\lambda\) such that \(A_\xi\) remains constant on no interval
\((\lambda-h,\lambda+h)\).  The result just established may then be stated as
follows: \emph{the points of the spectrum represent singular values for the
substitution \(A\); for every other value of \(\lambda\), the inverse
\((\lambda E-A)^{-1}\) exists and is given by the integral
\eqref{eq:95:15}.}  In particular, when \(\lambda=0\) does not belong to the
spectrum, \(A^{-1}\) exists and
\[
  A^{-1}=\int_{-\infty}^{\infty}\frac{dA_\xi}{\xi}.
\]

The spectrum is contained in the closed interval \([m,M]\).  Moreover,
\emph{the two endpoints \(m\) and \(M\) always belong to the spectrum}.  Indeed,
by its definition the spectrum is a closed set;
% [Source printed page 140]
let \(m'\) and \(M'\) be its two endpoints.  The form \(A_\xi\), regarded as a
function of \(\xi\), is constant for \(\xi<m'\) and for \(\xi>M'\); therefore in
\eqref{eq:94:14} we may replace the limits of integration by
\(m'-\varepsilon\) and \(M'+\varepsilon\), where \(\varepsilon>0\) is arbitrary.
Applied to the functions \(f(\xi)=1\) and \(f(\xi)=\xi\), the modified formula
gives
\[
  \int_{m'-\varepsilon}^{M'+\varepsilon}dA_\xi=E,
  \qquad
  \int_{m'-\varepsilon}^{M'+\varepsilon}\xi\,dA_\xi=A.
\]
Since \(A_\xi\) is monotone, these formulas imply
\[
  m'E\le A\le M'E.
\]
Thus \(m'\le m\) and \(M'\ge M\).  But there are no singular values outside
\([m,M]\); consequently \(m'=m\) and \(M'=M\).

\numberedparagraph{96}
To carry the study of the spectrum farther, introduce alongside \(A_\xi\) another
form \(P_\xi(x,x)\), likewise depending on the parameter \(\xi\).  For each
fixed \(\xi\), this means the form corresponding to the function
\(f_\xi(\mu)=0\) for \(\mu\ne\xi\) and \(f_\xi(\mu)=1\) for \(\mu=\xi\).
By II, \(P_\xi\ge0\); and by I and II, for every finite number of distinct points
\(\xi_1,\ldots,\xi_n\),
\(P_{\xi_1}+P_{\xi_2}+\cdots+P_{\xi_n}\le E\).  It follows readily that the set
of points \(\xi\) for which \(P_\xi\) does not vanish identically is, if it
exists, finite or countable.  This set is called the \emph{point spectrum} of
the form \(A\).

Here are some further properties of \(P_\xi\) which follow immediately from its
definition.  The relation \(f_\xi^2=f_\xi\) gives \(P_\xi^2=P_\xi\), and for
\(\xi_1\ne\xi_2\), the relation \(f_{\xi_1}f_{\xi_2}=0\) gives
\(P_{\xi_1}P_{\xi_2}=0\).  The relations
\(f_\xi\mu=\mu f_\xi=\xi f_\xi\) give
\(P_\xi A=AP_\xi=\xi P_\xi\).

From the last relation follows the fundamental identity
\((\lambda E-A)P_\lambda=0\).  It tells us that every system \((y_k)\) obtained
from a system \((x_k)\) by applying \(P_\lambda\) is a solution of the
homogeneous system
\begin{equation}
  \lambda y_i-\sum_{k=1}^{\infty}a_{ik}y_k=0
  \qquad(i=1,2,\ldots).
  \tag{17}\label{eq:96:17}
\end{equation}
Suppose now that \(\lambda\) belongs to the point spectrum.  In that
% [Source printed page 141]
case, by the definition of the point spectrum, there are systems \((x_k)\)
such that
\[
  E(y,y)=P_\lambda^2(x,x)=P_\lambda(x,x)\ne0.
\]
Consequently, for every value \(\lambda\) belonging to the point spectrum,
the homogeneous system \eqref{eq:96:17} has nonzero solutions.

Conversely, suppose that the homogeneous system \eqref{eq:96:17} has a
nonzero solution \((y_k)\), that is, one for which \(E(y,y)\ne0\).  By
\eqref{eq:96:17}, applying the substitution \(A\) to \((y_k)\) amounts to
multiplying the \(y_k\) by \(\lambda\).  Consequently, applying \(f(A)\) to
\((y_k)\) amounts to multiplying the \(y_k\) by \(f(\lambda)\).  In
particular, the substitution \(P_\xi\) gives \((0)\) when \(\xi\ne\lambda\),
whereas for \(\xi=\lambda\) it carries \((y_k)\) into itself.  Finally,
\(P_\lambda(y,y)=E(y,y)\ne0\) by hypothesis; hence \(\lambda\) belongs to the
point spectrum.  We have therefore proved the following theorem.

\begin{theorem}
The system \eqref{eq:96:17} has a nonzero solution if and only if
\(\lambda\) belongs to the point spectrum.
\end{theorem}

It should also be observed that the formula
\(P_{\xi_1}P_{\xi_2}=0\) for \(\xi_1\ne\xi_2\) implies a result which may
equally easily be established more directly: \emph{two solutions of
\eqref{eq:96:17} corresponding to two distinct values of \(\lambda\) are
orthogonal to one another.}

\numberedparagraph{97}
There is a very intimate relation between the forms \(A_\xi\) and \(P_\xi\),
which we shall now establish.  To display it clearly, introduce also the form
\(A_\xi^{\pm}(x,x)\); by this we mean the form corresponding to the function
\(f(\mu)=1\) for \(\mu\le\xi\) and \(f(\mu)=0\) for \(\mu>\xi\).  Plainly,
\[
  P_\xi=A_\xi^{\pm}-A_\xi.
\]
Regard \(A_\xi\) and \(A_\xi^{\pm}\) as functions of \(\xi\).  They are monotone:
as \(\xi\) increases, they too increase, or at least do not decrease.  Thus
the four limits
\(A_{\xi-0},A_{\xi+0},A^{\pm}_{\xi-0},A^{\pm}_{\xi+0}\) exist.  Moreover, since
\(A_\xi^{\pm}\ge A_\xi\) always, while for \(\varepsilon>0\)
\[
  A^{\pm}_{\xi-\varepsilon}\le A_\xi,
  \qquad A_{\xi+\varepsilon}\ge A_\xi^{\pm},
\]
we have
\[
  A_{\xi-0}\le A^{\pm}_{\xi-0}\le A_\xi,
  \qquad
  A^{\pm}_{\xi+0}\ge A_{\xi+0}\ge A_\xi^{\pm}.
\]
Finally, the proposition stated at the end of
\hyperref[sec:91]{no.~91} gives
\(A_{\xi-0}=A_\xi\) and \(A^{\pm}_{\xi+0}=A_\xi^{\pm}\).  Consequently,
\[
  A_{\xi-0}=A^{\pm}_{\xi-0}=A_\xi,
  \qquad
  A_{\xi+0}=A^{\pm}_{\xi+0}=A_\xi^{\pm}.
\]
In other words, the forms \(A_\xi\) and \(A_\xi^{\pm}\), considered as functions
% [Source printed page 142]
of \(\xi\), are continuous and equal to one another at every \(\xi\) which
does not belong to the point spectrum.  At points \(\xi\) of the point
spectrum they are distinct and remain continuous only when
\(P_\xi(x,x)=0\); in general they undergo there a sudden change equal to
\(P_\xi(x,x)\).

Let us add, as is almost evident, that in most of the preceding arguments---for
example in formula \eqref{eq:94:14}---one may replace \(A_\xi\) by
\(A_\xi^{\pm}\), or, for symmetry, by
\(\tfrac12(A_\xi+A_\xi^{\pm})\).

\numberedparagraph{98}
The preceding considerations allow us, following Hilbert, to decompose the
form \(A\) into two forms of a more special kind.  Except perhaps at the point
\(\xi=0\), which will play an exceptional role, one of these forms will have
the same point spectrum and the same corresponding forms \(P_\xi\) as \(A\);
the point spectrum of the other, if it exists, will reduce to \(\xi=0\).

First make a general observation.  Suppose that \(A=B+C\), where the forms
\(B\) and \(C\) are orthogonal to one another, that is,
\(BC=CB=0\).  Under this hypothesis,
\begin{equation}
  f(A)=f(B)+f(C)-f(0)E.
  \tag{18}\label{eq:98:18}
\end{equation}
This equality is evident when \(f(\mu)\) is a polynomial; for the other
functions under consideration it follows by passage to the limit.

Now put
\[
  P(x,x)=\sum_\xi P_\xi(x,x),
\]
the sum extending over all \(\xi\), or equivalently over the values \(\xi\)
which belong to the point spectrum.  The forms \(P_\xi\) are nonnegative and,
for every finite set of distinct points \(\xi_1,\ldots,\xi_n\),
\(P_{\xi_1}+\cdots+P_{\xi_n}\le E\); hence the series converges.  The
relations \(P_\xi^2=P_\xi\) and
\(P_{\xi_1}P_{\xi_2}=0\) for \(\xi_1\ne\xi_2\), after summation, give
\(PP_\xi=P_\xi P=P_\xi\), and another summation gives \(P^2=P\).  Likewise,
from \(A^nP_\xi=P_\xi A^n=\xi^nP_\xi\), we obtain
\begin{equation}
  A^nP=PA^n=\sum_\xi \xi^nP_\xi.
  \tag{19}\label{eq:98:19}
\end{equation}
It follows that
\[
  AP(A-PA)=(A-AP)PA=APA-AP^2A=APA-APA=0.
\]
Thus, on putting \(AP=PA=B\) and \(A-B=C\), we have \(BC=CB=0\).
% [Source printed page 143]
The forms \(B\) and \(C\) are orthogonal, and their sum is \(A\).

Let us calculate \(f(B)\) and \(f(C)\).  I assert that
\begin{equation}
  f(B)=\sum_\xi f(\xi)P_\xi+f(0)(E-P).
  \tag{20}\label{eq:98:20}
\end{equation}
Indeed, \(B^n=A^nP^n=A^nP\); hence by \eqref{eq:98:19},
\(B^n=\sum_\xi\xi^nP_\xi\).  Formula \eqref{eq:98:20} is therefore true
when \(f\) is a polynomial.  Without loss of generality it is enough to extend
it to bounded functions \(f\) which are sums of series of nonnegative
polynomials.  Since the forms \(P_\xi\) are also nonnegative, this merely
amounts to interchanging the order of summation in a double series of
nonnegative terms, which is permissible.  Thus \eqref{eq:98:20} remains valid
for all the functions under consideration.

The expression for \(f(C)\) follows from \eqref{eq:98:18} and
\eqref{eq:98:20}:
\begin{equation}
  f(C)=f(A)+f(0)P-\sum_\xi f(\xi)P_\xi.
  \tag{21}\label{eq:98:21}
\end{equation}
Calculate in particular the forms \(P'_\xi\) and \(P''_\xi\), analogous to
\(P_\xi\), which correspond respectively to \(B\) and \(C\).  In
\eqref{eq:98:20} and \eqref{eq:98:21}, put \(f(\mu)=1\) for \(\mu=\xi\)
and \(f(\mu)=0\) for \(\mu\ne\xi\).  When \(\xi\ne0\), this gives
\[
  P'_\xi=P_\xi,\qquad P''_\xi=0.
\]
For \(\xi=0\), it gives
\[
  P'_0=P_0+E-P,\qquad P''_0=P.
\]
Thus \emph{apart from \(\xi=0\), the form \(B\) has the same point spectrum
and the same corresponding forms as \(A\); as for \(C\), its point spectrum
reduces to the point \(0\), and the corresponding form is
\(P=\sum_\xi P_\xi\).}

\numberedparagraph{99}
The importance of the preceding decomposition lies in the fact that, in
solving the homogeneous system \eqref{eq:96:17}, the form \(A\) may be
replaced by the more special form \(B\).  Conversely, when the totality of
solutions of \eqref{eq:96:17} is known, it uniquely determines \(B\).  To see
this, it is enough to express \(P_\xi\) in terms of the solutions corresponding
to \(\lambda=\xi\).  We saw that the totality of these solutions is obtained
by applying the substitution \(P_\xi\) to the whole Hilbert space.  If it is
% [Source printed page 144]
applied only to a countable subset dense everywhere in that space and the
resulting systems are orthogonalized,
\footnote{Cf. \hyperref[sec:50]{no.~50}.  There it was the equations themselves
that were orthogonalized; here the solutions are orthogonalized by an entirely
analogous procedure.  After numbering them, one suppresses those which are
linear combinations of their predecessors and numbers them anew; one then
adds to each of them a linear combination of those preceding it, chosen so
that the modified solution is orthogonal to the preceding ones.  Finally one
multiplies by suitable constants so that
\(\sum_{k=1}^{\infty}(l_k^{(i)})^2=1\).}
one obtains a finite or countable sequence of normalized, pairwise orthogonal
systems
\[
  (l_k^{(1)}),(l_k^{(2)}),\ldots,
\]
which solve \eqref{eq:96:17} and such that every solution can be approximated
strongly by them or by their linear combinations.

Consider the form
\[
  Q(x,x)=\sum_i\left(\sum_{k=1}^{\infty}l_k^{(i)}x_k\right)^2
\]
and its corresponding substitution \(Q\).  The series in parentheses plainly
converge.  When only finitely many systems \((l_k^{(i)})\) are involved, the
right-hand side therefore has a definite meaning.  In that case one sees at
once that \(Q^2=Q\), and consequently \((E-Q)^2=E-Q\).  Hence
\(E(x,x)-Q(x,x)\) is a nonnegative form, so that \(Q(x,x)\le E(x,x)\).  It
follows immediately that the displayed series converges in every case and
remains at most \(E(x,x)\).

I assert that \(P_\xi=Q\).  Each system \((l_k^{(i)})\) is reproduced upon
application of \(P_\xi\), since it is a solution of \eqref{eq:96:17}; and,
because the systems are normalized and pairwise orthogonal, each is likewise
reproduced upon application of \(Q\).  Thus \(P_\xi-Q\), applied to any one of
the systems \((l_k^{(i)})\), gives zero.  The same is therefore true for every
solution of \eqref{eq:96:17}, since such a solution can be approximated
indefinitely closely by linear combinations of the \((l_k^{(i)})\).  Finally,
\(P_\xi^2=P_\xi\) and \(QP_\xi=Q\), so that
\((P_\xi-Q)P_\xi=P_\xi-Q\).  Hence applying \(P_\xi-Q\) to a system amounts
first to applying \(P_\xi\) and then \(P_\xi-Q\).  The first operation gives
% [Source printed page 145]
a solution of \eqref{eq:96:17}, and the second gives zero on that solution.
Thus \(P_\xi-Q\) carries the whole Hilbert space into zero; hence
\(P_\xi=Q\).

Having decomposed each quadratic form \(P_\xi\) in this way into squares of
linear forms, formula \(B=\sum_\xi\xi P_\xi\) likewise gives a similar
decomposition of \(B\): one merely replaces each \(P_\xi\) by its corresponding
sum \(Q\).  Thus the special form \(B\) has a decomposition analogous to
\eqref{eq:84:2}, except that the summation may extend over a countable infinity
of squares.  In the general case, however, after the part \(B\) corresponding
to the solutions of \eqref{eq:96:17} has been separated from \(A\), a form
\(C\) still remains.  This form too has a special character: its point
spectrum, if it exists, reduces to \(\xi=0\); the form \(C_\xi\) is a
continuous function of \(\xi\), except perhaps at \(\xi=0\); and therefore,
after excluding the point \(0\) if it is isolated, the spectrum of \(C\) is a
perfect set.  This perfect set is called the \emph{continuous spectrum} of the
form \(A\).

\numberedparagraph{100}
When \(A\) is completely continuous, these results become still simpler.  In
that case \(B=AP\), and hence \(C=A-B\), are also completely continuous.  I
assert that \(C=0\).  Indeed, \(C\) has no point spectrum except perhaps at
\(\xi=0\); hence, for \(\xi\ne0\), the homogeneous system corresponding to
\(\xi E-C\) has no solution.  On the other hand, if a point \(\xi\ne0\)
belonged to the spectrum of \(C\), this system could be satisfied with
arbitrary accuracy by elements \((x_k)\) for which \(E(x,x)=1\).  There would
thus exist a sequence \((x_k^{(n)})\) such that
\(E(x^{(n)},x^{(n)})=1\) and such that the sequence obtained by applying
\(\xi E-C\) to these elements tends strongly to \((0)\).  From it one could
extract a subsequence \((x_k^{(n_j)})\) tending to an element \((x_k^*)\), and
the homogeneous system \(\xi E-C=0\) would have the solution \((x_k^*)\).
Moreover, since \(C\) is strongly continuous,
\(C(x_k^{(n_j)})\) would tend strongly to \(C(x_k^*)\), and then
\((x_k^{(n_j)})\) would tend strongly to its limit.  We would therefore have
\(E(x^*,x^*)=1\), so that \(\xi E-C=0\) would possess a nonzero solution.
Consequently no point \(\xi\ne0\) belongs to the spectrum, and the form
\(C_\xi\), regarded as a function of \(\xi\),
% [Source printed page 146]
must be constant for \(\xi<0\) and also for \(\xi>0\).  Thus
\[
  C=\int_{-\infty}^{+\infty}\xi\,dC_\xi=0.
\]

Since \(C=0\), we have \(A=B\).  Consequently \(A(x,x)\) has a decomposition
analogous to \eqref{eq:84:2}:
\begin{equation}
  A(x,x)=\sum_{j=1}^{\infty}\mu_j
  \left(\sum_{k=1}^{\infty}l_{jk}x_k\right)^2,
  \qquad
  E(x,x)=\sum_{j=1}^{\infty}
  \left(\sum_{k=1}^{\infty}l_{jk}x_k\right)^2.
  \tag{22}\label{eq:100:22}
\end{equation}
The linear forms \(\sum_k l_{jk}x_k\) are normalized and pairwise orthogonal.
The coefficients \(\mu_j\) run through the spectrum of \(A\), taking each
value \(\xi\) in it one or more times.  For each such value, the corresponding
linear forms are obtained by decomposing \(P_\xi\).  In the general case this
decomposition may yield infinitely many terms; in the present special case it
does not.  Here only \(\xi=0\) can yield infinitely many terms; each other
value yields only finitely many.  Moreover, I assert that \(\mu_j\to0\), which
plainly implies the fact just stated.  Indeed, the normalized orthogonal
elements \((l_{jk})\), \(j=1,2,\ldots\), tend to \((0)\); applying \(A\) gives
the sequence \((\mu_jl_{jk})\).  Since \(A\) is completely continuous, this
sequence tends strongly to \((0)\), that is, \(\mu_j^2\to0\).

In summary, \emph{every completely continuous form \(A(x,x)\) has the
expansion \eqref{eq:100:22}, with \(\mu_j\to0\).}

It is natural to expect that this result can be established more directly,
without using the general theory of the spectrum.  Hilbert gave such a
proof.\footnote{Hilbert, fourth memoir, p.~201.}  Its essential idea rests on
the following fact.  The completely continuous form \(A(x,x)\), varied under
the condition \(E(x,x)=1\), attains its maximum \(M\) at an element
\((l_k^{(1)})\) and its minimum \(m\) at an element \((l_k^{(2)})\); these
elements satisfy respectively the homogeneous systems \(ME-A=0\) and
\(mE-A=0\).  Further, if \(A(x,x)\) is varied while the \((x_k)\) are also
required to be orthogonal to certain
% [Source printed page 147]
elements satisfying systems of the same type \(\mu E-A=0\), the extremal
elements satisfy systems of that type as well.  This fact permits one to
exhaust successively all solutions of the homogeneous systems and thus leads
to expression \eqref{eq:100:22}.

Von Koch showed that, when
\(\sum_i|a_{ii}|\) and \(\sum_{i,k}a_{ik}^2\) converge, expansion
\eqref{eq:100:22} can also be obtained by the method of infinite
determinants.\footnote{H.~von Koch, ``Sur un théorème de M.~Hilbert,''
\emph{Mathematische Annalen}, vol.~LXIX (1910), pp.~266--283.}

\section{The Continuous Spectrum and Differential Solutions}

\numberedparagraph{101}
Let us now study the form \(C\) corresponding to the continuous spectrum.  For
simplicity suppose that \(C=A\), that is, suppose that \(A\) has no point
spectrum.  The modifications needed in the general case are almost evident;
only the point \(\xi=0\) is involved.

Our hypothesis says that, for no value of \(\lambda\), does system
\eqref{eq:96:17} have solutions \((x_k)\) such that \(\sum_kx_k^2\) converges
and is nonzero.  Hilbert gave a very simple example of such a form:
\begin{equation}
  A(x,x)=x_1x_2+x_2x_3+x_3x_4+\cdots.
  \tag{23}\label{eq:101:23}
\end{equation}
Let us discuss the corresponding systems \eqref{eq:96:17}.  They are
recurrence systems: after choosing \(\lambda\) and \(x_1\) arbitrarily, one
can calculate \(x_2,x_3,\ldots\) successively.  Since the Cauchy--Lagrange
inequality gives \(-E\le A\le E\), we need consider only values of
\(\lambda\) between \(-1\) and \(1\).  Following Hellinger, put
\(\lambda=\cos t\) and \(x_1=c\sin t\).  A calculation easily restored gives
\[
  x_k=c\sin kt\qquad(k=1,2,\ldots).
\]
As \(t\) and \(c\) vary, this formula gives all solutions of
\eqref{eq:96:17} for \(-1\le\lambda\le1\).  For these solutions,
\(\sum_kx_k^2\) does not converge.  Thus the \((x_k)\) are not admissible
solutions from our present point of view, and one might think that our methods
do not apply to the study of these singular solutions.  But not all is lost.

% [Source printed page 148]
For example, regard the solution
\(x_k(\lambda)=\sin(k\arccos\lambda)\) as a function of \(\lambda\), and let
\(y_k\) be the integral of \(x_k(\lambda)\) with respect to \(\lambda\) over
an interval \(I\).  Since \(ky_k\) remains bounded as \(k\to\infty\), the
series \(\sum_ky_k^2\) converges.  If \(z_k\) are analogous quantities
corresponding to another interval which does not overlap \(I\), then the
elements \((y_k)\) and \((z_k)\) are orthogonal.  This follows from an easy
calculation; we shall soon see the deeper reasons for it.

When the point spectrum is absent, the analogy between these integrals of
singular solutions and the admissible solutions in the other case suggested
to Hellinger the idea of basing the study of the continuous spectrum on these
integrals, without considering the solutions themselves.  Another important
reason also compelled him to do so.  To explain it, I must first indicate the
standpoint adopted by Hilbert and his school in studying the spectrum.  The
necessary and sufficient condition for two quadratic forms in \(n\) variables
to be transformable into one another by an orthogonal substitution of the
variables is that their singular values and respective multiplicities
coincide.  What are the corresponding conditions for infinitely many
variables?  When the theory of quadratic forms is organized around this
problem, the facts to be considered are naturally delimited: one confines
oneself to facts which persist under every orthogonal substitution of the
variables; in short, to the \emph{orthogonal invariants} of the quadratic
form.\footnote{The following observation connects the study of orthogonal
invariants with our present line of thought.  Let \(O\) be a real orthogonal
substitution; thus, for its coefficients \(o_{ik}\),
\[
  \sum_{k=1}^{\infty}o_{ik}o_{jk}
  =\sum_{k=1}^{\infty}o_{ki}o_{kj}
  =\begin{cases}1,&i=j,\\0,&i\ne j.\end{cases}
\]
Applying \(O\) to Hilbert space replaces the forms \(A\) by
\(OAO^{-1}\).  Plainly
\(f(OAO^{-1})=Of(A)O^{-1}\); thus the correspondence between \(A\) and
\(f(A)\) is unchanged.}
The existence
% [Source printed page 149]
of an admissible solution of \eqref{eq:96:17} for a given \(\lambda\) is such
an invariant; indeed, admissible solutions correspond under the same
substitutions as the forms themselves.  This is not so for solutions \((x_k)\)
for which \(\sum_kx_k^2\) diverges: orthogonal substitutions are not defined
on such systems and, in most cases, their definition cannot be extended to
them.

Regard, then, \(x_k\) as a function of \(\lambda\), and suppose that the
equations \eqref{eq:96:17} may be integrated term by term with respect to
\(\lambda\).  Put
\[
  \rho_k(\xi)=\int_0^\xi x_k(\lambda)\,d\lambda.
\]
Using Stieltjes notation, we obtain
\begin{equation}
  \int_{\lambda_1}^{\lambda_2}\xi\,d\rho_i(\xi)
  -\sum_{k=1}^{\infty}a_{ik}
  \bigl[\rho_k(\lambda_2)-\rho_k(\lambda_1)\bigr]=0,
  \qquad i=1,2,\ldots.
  \tag{24}\label{eq:101:24}
\end{equation}

\begin{definition}
In Hellinger's terminology, a \emph{differential solution} is a sequence of
continuous functions \([\rho_k(\xi)]\) such that
\(\sum_{k=1}^{\infty}[\rho_k(\xi)]^2\) converges to a continuous function and
such that the \(\rho_k\) satisfy \eqref{eq:101:24} for every pair
\(\lambda_1,\lambda_2\).
\end{definition}

To obtain such solutions, apply \(A_\xi\) to an arbitrary system \((x_k)\)
for which \(\sum_kx_k^2\) converges.  For every \(\xi\) there results a
definite system \([\rho_k(\xi)]\).  Under our hypothesis on \(A\), the form
\(A_\xi\) is continuous in \(\xi\), and the functions \(\rho_k(\xi)\) are
continuous as well.  Moreover,
\[
  \sum_{k=1}^{\infty}[\rho_k(\xi)]^2
  =A_\xi^2(x,x)=A_\xi(x,x),
\]
which is a continuous function of \(\xi\).  Finally,
\[
  \int_{\lambda_1}^{\lambda_2}\xi\,dA_\xi
  =A\bigl(A_{\lambda_2}-A_{\lambda_1}\bigr).
\]
Indeed, both sides equal \(f(A)\), where \(f\) is the function equal to
\(\xi\) for \(\lambda_1\le\xi<\lambda_2\) and zero elsewhere.  It follows
that the functions \(\rho_k\) satisfy \eqref{eq:101:24}.

Differential solutions share, \emph{mutatis mutandis}, all the essential
properties of admissible solutions of
% [Source printed page 150]
\eqref{eq:96:17}; all this amounts to saying that their variation over small
intervals behaves approximately like exact solutions of
\eqref{eq:96:17}.  Among other things, they may be used to decompose the form
\(A\) into simpler parts and to study the equivalence of two forms with
respect to orthogonal substitutions.  For all details we refer the reader to
the original memoirs.\footnote{E.~Hellinger, \emph{Die Orthogonalinvarianten
quadratischer Formen von unendlich vielen Variabeln} (doctoral dissertation),
Göttingen, 1907; and ``Neue Begründung der Theorie quadratischer Formen von
unendlich vielen Veränderlichen,'' \emph{Journal für die reine und angewandte
Mathematik}, vol.~CXXXVI (1909), pp.~210--271.  Hellinger makes very successful
use of a highly interesting notion of generalized integral analogous to
Stieltjes's.  At first sight the use of these integrals seems indispensable.
Very recently H.~Hahn discovered their intimate relation to ordinary
Lebesgue integrals, in particular to the summable functions with summable
squares discussed below; by this discovery he was able to solve completely
the problem of equivalence of two quadratic forms.  See H.~Hahn, ``Über die
Integrale des Herrn Hellinger und die Orthogonalinvarianten der quadratischen
Formen von unendlich vielen Veränderlichen,'' \emph{Monatshefte für Mathematik
und Physik}, vol.~XXIII (1912), pp.~161--224.}

\numberedparagraph{102}
Let us mention one more very interesting class of forms for which all these
arguments become much simpler.  It includes, among others, form
\eqref{eq:101:23}.\footnote{Hilbert, fourth memoir, p.~207.  Hilbert treats only
the special case \(u(\mu)=1/\mu\).  The class studied here is contained in a
larger class treated by Hellinger.}

Suppose that the functions
\(\varphi_1(\mu),\varphi_2(\mu),\ldots\), defined on the interval \((a,b)\),
form a normalized orthogonal system, that is,
\[
  \int_a^b\varphi_i(\mu)\varphi_k(\mu)\,d\mu
  =\begin{cases}1,&i=k,\\0,&i\ne k.\end{cases}
\]
Suppose further that, for every function \(f(\mu)\) which is integrable and
whose square is integrable, the integral
\[
  \int_a^b\left[f(\mu)-\sum_{k=1}^{n}c_k\varphi_k(\mu)\right]^2d\mu
\]
can be made arbitrarily small
% [Source printed page 151]
by a suitable choice of \(n\) and the \(c_k\).
This last property is expressed by saying that the system
\([\varphi_k(\mu)]\) is \emph{complete}.\footnote{As an example of a
normalized, complete orthogonal system on \((-\pi,\pi)\), take
\[
  \frac1{\sqrt{2\pi}},\quad
  \frac1{\sqrt\pi}\cos\mu,\quad
  \frac1{\sqrt\pi}\sin\mu,\quad
  \frac1{\sqrt\pi}\cos2\mu,\quad
  \frac1{\sqrt\pi}\sin2\mu,\ldots.
\]
The system is plainly normalized and orthogonal.  Its completeness follows,
for example, from Weierstrass's theorem that every continuous function of
period \(2\pi\) can be approximated uniformly by trigonometric polynomials,
together with the fact that for every integrable \(f(\mu)\) whose square is
integrable, \(\int_{-\pi}^{\pi}[f(\mu)-\varphi(\mu)]^2d\mu\) may be made
arbitrarily small by a suitable continuous function \(\varphi\) of period
\(2\pi\).  Applied to this special case, the formula in the text is Parseval's
theorem, also called the fundamental theorem of Fourier constants.  Cf.
A.~Hurwitz, ``Über die Fourierschen Konstanten integrierbarer Funktionen,''
\emph{Mathematische Annalen}, vol.~LVII (1903), pp.~425--446, and vol.~LIX
(1904), p.~553; P.~Fatou, ``Séries trigonométriques et séries de Taylor,''
\emph{Acta Mathematica}, vol.~XXX (1906), pp.~335--400.  Every system used
below can easily be reduced to the case just considered.}
For a system which is also normalized and orthogonal, the best choice of the
\(c_k\), for each \(n\), is
\[
  c_k=f_k=\int_a^b f(\mu)\varphi_k(\mu)\,d\mu,
\]
as follows immediately from
\[
  \int_a^b\left[f(\mu)-\sum_{k=1}^{n}c_k\varphi_k(\mu)\right]^2d\mu
  =\int_a^b\left[f(\mu)-\sum_{k=1}^{n}f_k\varphi_k(\mu)\right]^2d\mu
   +\sum_{k=1}^{n}(f_k-c_k)^2.
\]
It follows that, as \(n\to\infty\),
\[
  \int_a^b[f(\mu)]^2d\mu-\sum_{k=1}^{n}f_k^2
  =\int_a^b\left[f(\mu)-\sum_{k=1}^{n}f_k\varphi_k(\mu)\right]^2d\mu
  \longrightarrow0;
\]
that is,
% [Source printed page 152]
\[
  \int_a^b[f(\mu)]^2d\mu=\sum_{k=1}^{\infty}f_k^2.
\]
Finally, applying this formula to
\(\tfrac12[f(\mu)+g(\mu)]\) and \(\tfrac12[f(\mu)-g(\mu)]\), and
subtracting, gives
\begin{equation}
  \int_a^b f(\mu)g(\mu)\,d\mu
  =\sum_{k=1}^{\infty}
  \left(\int_a^b f(\mu)\varphi_k(\mu)\,d\mu\right)
  \left(\int_a^b g(\mu)\varphi_k(\mu)\,d\mu\right),
  \tag{24}\label{eq:102:24}
\end{equation}
where \(f\) and \(g\), as well as their squares, are assumed integrable.

All these considerations hold when \emph{integral} is understood in
Lebesgue's sense, and therefore \emph{a fortiori} for the classical notion of
integral.  To fix ideas, suppose that the functions in question are continuous
or consist of finitely many continuous pieces.  Put
\[
  a_{ik}=\int_a^b u(\mu)\varphi_i(\mu)\varphi_k(\mu)\,d\mu,
  \qquad
  A(x,x)=\sum_{i,k=1}^{\infty}a_{ik}x_ix_k.
\]
The sections \([A(x,x)]_n\) and \([E(x,x)]_n\) are plainly represented by
\[
  \int_a^b u(\mu)\left[\sum_{k=1}^{n}x_k\varphi_k(\mu)\right]^2d\mu,
  \qquad
  \int_a^b\left[\sum_{k=1}^{n}x_k\varphi_k(\mu)\right]^2d\mu.
\]
Consequently \(A(x,x)\) is bounded and, when \(E(x,x)=1\), the set of its
values lies between the lower and upper bounds of \(u(\mu)\).  Let
\(A^*(x,x)\) be the form defined analogously by a function \(u^*(\mu)\).
Putting \(f(\mu)=u(\mu)\varphi_i(\mu)\) and
\(g(\mu)=u^*(\mu)\varphi_k(\mu)\) in \eqref{eq:102:24}, we find that the
coefficients of \(AA^*\) are
\[
  \int_a^b u(\mu)u^*(\mu)\varphi_i(\mu)\varphi_k(\mu)\,d\mu.
\]
In particular, the coefficients of \(A^2,A^3,\ldots\) are obtained by
replacing \(u(\mu)\) with \(u^2(\mu),u^3(\mu),\ldots\) in the expression for
\(a_{ik}\).

% [Source printed page 153]
It follows that the coefficients of \(f(A)\) are
\[
  \int_a^b f[u(\mu)]\varphi_i(\mu)\varphi_k(\mu)\,d\mu.
\]
This expression immediately opens the way to the study of the spectrum and
the spectral forms.  Let \(\mathfrak E_\xi,\mathfrak E_\xi^{\pm},\mathfrak E'_\xi\)
denote the sets of values \(\mu\) for which, respectively,
\(u(\mu)<\xi\), \(u(\mu)\le\xi\), and \(u(\mu)=\xi\).  The coefficients of
\(A_\xi,A_\xi^{\pm},P_\xi\) are given by
\[
  \int \varphi_i(\mu)\varphi_k(\mu)\,d\mu,
\]
where the integral extends over \(\mathfrak E_\xi,\mathfrak E_\xi^{\pm}\), and
\(\mathfrak E'_\xi\), respectively.  If, for example, \(u(\mu)\) is continuous
and has \(m\) and \(M\) as its extreme values, the spectrum is the interval
\([m,M]\).  If it takes neither the value \(\xi\) nor neighboring values,
then \(\xi\) and its neighborhood do not belong to the spectrum.  Finally, if
it takes the value \(\xi\) throughout an interval, then \(\xi\) belongs to the
point spectrum.

Form \eqref{eq:101:23} belongs to this class on putting
\[
  a=0,\qquad b=\pi,\qquad
  \varphi_k(\mu)=\sqrt{\frac2\pi}\sin k\mu,\qquad u(\mu)=\cos\mu.
\]

\numberedparagraph{103}
Consider another special case, intimately connected with the theory of
trigonometric series.  Here it is convenient to let the indices range from
\(-\infty\) to \(+\infty\).  This merely renumbers the variables \(x_k\) and
therefore has no effect on the essential results of the theory.

Put
\[
  a=-\pi,\qquad b=\pi,\qquad
  \varphi_k(\mu)=\frac{1}{\sqrt{2\pi}}(\sin k\mu+\cos k\mu)
  =\frac1{\sqrt\pi}\cos\left(k\mu-\frac\pi4\right).
\]
As \(k\) runs through all integers from \(-\infty\) to \(+\infty\), the
\(\varphi_k(\mu)\) form a normalized, complete orthogonal system.  Consider
the form
\[
  A(x,x)=\sum_{i,k=-\infty}^{+\infty}a_{ik}x_ix_k,
\]
% [Source printed page 154]
where
\begin{align*}
  a_{ik}
  &=\int_{-\pi}^{\pi}u(\mu)\varphi_i(\mu)\varphi_k(\mu)\,d\mu\\
  &=\frac1{2\pi}\int_{-\pi}^{\pi}u(\mu)\cos(i-k)\mu\,d\mu
   +\frac1{2\pi}\int_{-\pi}^{\pi}u(\mu)\sin(i+k)\mu\,d\mu.
\end{align*}
Denote the two terms on the right by \(\alpha_{i-k}\) and \(\beta_{i+k}\).
Then \(A(x,x)\) is the sum of the two forms
\[
  A^{(1)}(x,x)=\sum_{i,k=-\infty}^{+\infty}\alpha_{i-k}x_ix_k,
  \qquad
  A^{(2)}(x,x)=\sum_{i,k=-\infty}^{+\infty}\beta_{i+k}x_ix_k.
\]
The form \(A^{(1)}\) corresponds to the even function
\[
  u^{(1)}(\mu)=\frac{u(\mu)+u(-\mu)}2,
\]
and \(A^{(2)}\) to the odd function
\[
  u^{(2)}(\mu)=\frac{u(\mu)-u(-\mu)}2.
\]
These forms belong to two very special types.  The first, corresponding to
even functions, was first considered by Toeplitz; we shall discuss it in the
\hyperref[chap:6]{next chapter}.  The second corresponds to odd functions.  For example, put
\[
  u(\mu)=
  \begin{cases}
    -\pi-\mu,&-\pi<\mu<0,\\
    \pi-\mu,&0<\mu<\pi.
  \end{cases}
\]
The preceding integrals give
\[
  a_{ik}=\frac1{i+k}\quad(i+k\ne0),
  \qquad a_{ik}=0\quad(i+k=0).
\]

The lower and upper bounds of the piecewise-defined function \(u(\mu)\) on
\((-\pi,\pi)\) are \(-\pi\) and \(\pi\)\translatornote{The French text
mistakenly writes \(u(\mu)=\mu\) here; the statement concerns the piecewise
function just defined.}; hence these are also the
lower and upper bounds of \(A(x,x)\), and therefore, by
\hyperref[sec:85]{no.~85}, of the corresponding bilinear form
\(A(x,y)\).  These forms include, in particular, the two forms considered by
Hilbert:
\[
  \sum_{i,k=1}^{\infty}\frac{x_ix_k}{i+k},
  \qquad
  \mathop{\sum_{i,k=1}^{\infty}}_{i\ne k}\frac{x_iy_k}{i-k}.
\]
The first is obtained from \(A(x,x)\) by putting
\(x_0=x_{-1}=x_{-2}=\cdots=0\); the second
% [Source printed page 155]
is obtained from \(A(x,y)\) by putting
\(x_0=x_{-1}=x_{-2}=\cdots=0\),
\(y_0=y_1=y_2=\cdots=0\), and then reversing the signs of the indices of the
\(y\)'s.  Thus, under \(E(x,x)\le1\), \(E(y,y)\le1\), both forms remain
between \(-\pi\) and \(\pi\).

The special interest of these forms lies in the fact that they stand, so to
speak, on the boundary of the class of bounded forms.  Indeed, very similar
expressions fail to represent bounded forms.  Consider, for example, the
following expression.\translatornote{The French proof treats all \(\alpha<1\)
with the first estimate below, which directly applies only when
\(0\le\alpha<1\).  The omitted case \(\alpha<0\) is supplied separately.}
\[
  \sum_{i,k=1}^{\infty}\frac{x_ix_k}{(i+k)^\alpha},
  \qquad \alpha<1.
\]
Put \(x_1=x_2=\cdots=x_n=1/\sqrt n\) and
\(x_{n+1}=x_{n+2}=\cdots=0\), consistently with \(E(x,x)=1\).  If
\(0\le\alpha<1\), all terms for which \(i+k\le n+1\) are at least
\(1/[n(n+1)^\alpha]\), and there are
\(1+2+\cdots+n=n(n+1)/2\) of them; thus the value of the expression is at least
\(\tfrac12(n+1)^{1-\alpha}\), and grows without bound with \(n\).  If
\(\alpha<0\), all \(n^2\) nonzero terms are at least \(2^{-\alpha}/n\), so
their sum is at least \(2^{-\alpha}n\) and is again unbounded.  Likewise,
making the same choice in
\[
  \mathop{\sum_{i,k=1}^{\infty}}_{i\ne k}\frac{x_ix_k}{|i-k|}
\]
gives
\[
  2\left(1+\frac12+\cdots+\frac1{n-1}\right)-\frac{2(n-1)}n,
\]
which also grows without bound.\footnote{The latter fact may be expressed by
saying that the bounded bilinear form
\(\mathop{\sum_{i,k=1}^{\infty}}_{i\ne k}x_iy_k/(i-k)\) is not absolutely
bounded.  It was noted by Hellinger and Toeplitz in their memoir cited in
\hyperref[sec:57]{no.~57}, p.~308.  For analogous forms, cf. O.~Toeplitz,
``Zur Theorie der quadratischen Formen von unendlich vielen Veränderlichen,''
\emph{Nachrichten der Königlichen Gesellschaft der Wissenschaften zu
Göttingen} (1910), pp.~489--506; and I.~Schur, ``Bemerkungen zur Theorie der
beschränkten Bilinearformen mit unendlich vielen Veränderlichen,''
\emph{Journal für die reine und angewandte Mathematik}, vol.~CXL (1911),
pp.~1--28.}

% [Source printed page 156]
\chapter{Applications}\label{chap:6}

\section{Linear Differential Equations}

\numberedparagraph{104}
While studying the method of infinite determinants in \hyperref[chap:2]{Chapter~II}, we stated
that this method was invented and first applied by Hill to the linear
differential equation
\begin{equation}
  \frac{d^2w}{dt^2}
  +\bigl(\theta_0+2\theta_1\cos t+2\theta_2\cos2t+\cdots\bigr)w=0.
  \tag{1}\label{eq:104:1}
\end{equation}
The argument of the learned astronomer, too concise for the pure geometer, was
perfected and fully justified by Poincaré.  We have presented Poincaré's
argument, as well as later investigations into the convergence of infinite
determinants and their application to systems of linear equations in
infinitely many unknowns.  We did not, however, return to equation
\eqref{eq:104:1}.  We do so now, but instead of considering only the special
equation \eqref{eq:104:1},\footnote{Equation \eqref{eq:104:1} was studied in
detail by Poincaré, loc. cit.; see \hyperref[sec:17]{nos.~17} and
\hyperref[sec:19]{19}.} we shall follow von Koch in placing ourselves in a
more general setting.

Consider the homogeneous linear differential equation of order \(n\)
\begin{equation}
  \frac{d^nu}{dz^n}+P_1(z)\frac{d^{n-1}u}{dz^{n-1}}
  +P_2(z)\frac{d^{n-2}u}{dz^{n-2}}+\cdots+P_n(z)u=0,
  \tag{2}\label{eq:104:2}
\end{equation}
where the coefficients \(P_1(z),\ldots,P_n(z)\) have Laurent expansions
\begin{equation}
  P_m(z)=\sum_{k=-\infty}^{+\infty}a_{mk}z^k,
  \tag{3}\label{eq:104:3}
\end{equation}
% [Source printed page 157]
convergent in the circular annulus \(r<|z|<R\).  From the investigations of
Fuchs it is known that there exists at least one integral of the form
\[
  u(z)=z^\rho\sum_{k=-\infty}^{+\infty}\alpha_kz^k,
\]
where the Laurent series converges in the annulus.  Our task is to determine
the constant \(\rho\) and the coefficients \(\alpha_k\) so that \(u(z)\) is an
integral of \eqref{eq:104:2}.

Equation \eqref{eq:104:1} plainly belongs to this class if the function
\(\theta(t)=\theta_0+2\theta_1\cos t+\cdots\) is holomorphic for real values
of \(t\), and if in addition we make the change of variable \(z=e^{it}\).

\numberedparagraph{105}
First observe that, without loss of generality, we may suppose
\(P_1(z)=0\) and \(r<1<R\).  Indeed, these assumptions amount to replacing
\(u\) and \(z\), respectively, by
\[
  u'=u\exp\left(-\frac1n\int P_1(z)\,dz\right),
  \qquad z'=\frac{z}{\sqrt{Rr}};
\]
the modified equation plainly satisfies the assumptions made.

Substitute the expansions of \(P_m(z)\) and \(u(z)\) into
\eqref{eq:104:2}.  Comparing coefficients gives the equations
\begin{equation}
  \sum_{k=-\infty}^{+\infty}\varphi_{i-k}(\rho+k)\alpha_k=0,
  \qquad i=-\infty,\ldots,+\infty,
  \tag{4}\label{eq:105:4}
\end{equation}
where
\begin{align*}
  \varphi_0(\rho)
  ={}&\rho(\rho-1)\cdots(\rho-n+1)\\
  &+\rho(\rho-1)\cdots(\rho-n+3)a_{2,-2}
  +\cdots+\rho a_{n-1,-n+1}+a_{n,-n},
\end{align*}
and, for \(j\ne0\),
\[
  \varphi_j(\rho)
  =\rho(\rho-1)\cdots(\rho-n+3)a_{2,j-2}
  +\cdots+\rho a_{n-1,j-n+1}+a_{n,j-n}.
\]
To put \eqref{eq:105:4} into the form
\begin{equation}
  \sum_{k=-\infty}^{+\infty}\psi_{ik}(\rho)\alpha_k=0,
  \qquad i=-\infty,\ldots,+\infty,
  \tag{5}\label{eq:105:5}
\end{equation}
% [Source printed page 158]
with \(\psi_{ii}=1\) for every \(i\), put
\[
  \psi_{ik}(\rho)=
  \frac{\varphi_{i-k}(\rho+k)}{\varphi_0(\rho+i)}.
\]
Let \(\rho_1,\ldots,\rho_n\) be the roots of
\(\varphi_0(\rho)=0\).  The poles of the rational functions
\(\psi_{ik}(\rho)\) then belong to the set of values
\(\rho_1-i,\ldots,\rho_n-i\), \(i=-\infty,\ldots,+\infty\), which are the
roots of one or another of the equations \(\varphi_0(\rho+i)=0\).

In the \(\rho\)-plane consider a finite or infinite domain \(D\), lying
between two lines parallel to the imaginary axis, such that none of these
roots lies in its interior or on its boundary.  We shall show that
\begin{equation}
  \mathop{\sum_{i,k=-\infty}^{+\infty}}_{i\ne k}
  |\psi_{ik}(\rho)|
  \tag{6}\label{eq:105:6}
\end{equation}
converges uniformly in \(D\).

For this purpose, write the polynomials
\(\varphi_j(\rho+i-j)\), \(j\ne0\), in the form
\[
  H_2(j)(\rho+i)^{n-2}+H_3(j)(\rho+i)^{n-3}
  +\cdots+H_n(j).
\]
Each \(H_r(j)\) is a sum of finitely many terms of the form
\(j^qa_{mj}\), \(q=0,1,2,\ldots\).  The Laurent series
\eqref{eq:104:3} converge in an annulus containing \(|z|=1\); hence the
series
\[
  \sum_{j=-\infty}^{+\infty}|j^qa_{mj}|
  \qquad(q=0,1,2,\ldots)
\]
converge, and therefore so do
\[
  S_r=\sum_{j=-\infty}^{+\infty}|H_r(j)|,
  \qquad r=2,3,\ldots,n.
\]
By the hypothesis on \(D\), there is a constant \(C_1\) such that
\[
  \frac1{|\varphi_0(\rho+i)|}<\frac{C_1}{(1+|i|)^n},
\]
% [Source printed page 159]
and a constant \(C_2\) such that
\(|\rho+i|<C_2(1+|i|)\).  Consequently,
\[
  \mathop{\sum_{i,k=-\infty}^{+\infty}}_{i\ne k}
  |\psi_{ik}(\rho)|
  \le (S_2+S_3+\cdots+S_n)C_3
  \sum_{i=-\infty}^{+\infty}\frac1{(1+|i|)^2},
\]
which proves the uniform convergence of \eqref{eq:105:6}.

It follows that the infinite determinant \(\Omega(\rho)\) corresponding to
\eqref{eq:105:5} is of normal type.  It and its minors converge for every
\(\rho\), except perhaps at roots of the equations
\(\varphi_0(\rho+i)=0\), and the convergence is uniform in every domain \(D\)
of the kind just described.  Thus \(\Omega(\rho)\) is holomorphic except at
those roots, where it has the character of a rational function.  The last
assertion follows because any one value \(\rho\) can satisfy only finitely many
of the equations \(\varphi_0(\rho+i)=0\).  More precisely, \(\rho\) is a pole
of order \(m\) of \(\Omega(\rho)\) if exactly \(m\) roots of
\(\varphi_0(\rho)=0\), counted with multiplicity, differ from \(\rho\) by an
integer or by zero.

The function \(\Omega(\rho)\) has period one, since
\[
  \psi_{ik}(\rho+1)=\psi_{i+1,k+1}(\rho),
\]
and replacing \(\rho\) by \(\rho+1\) therefore does not change the
determinant.  Moreover, as \(\rho\to\infty\) while its real part remains
bounded, the functions \(\psi_{ik}(\rho)\), \(i\ne k\), tend to zero; hence
\(\Omega(\rho)\to1\).  Consequently, \(\Omega(\rho)\) is a rational function
\(R(\omega)\) of
\[
  \omega=e^{2\pi i\rho},
\]
and \(R(0)=R(\infty)=1\).  The poles of \(R\) necessarily have the form
\(\omega_k=e^{2\pi i\rho_k}\); more precisely, \(\omega\) is a pole of order
\(m\) if it occurs \(m\) times among \(\omega_1,\ldots,\omega_n\).  Thus
\[
  R(\omega)=\frac{P(\omega)}
  {(\omega-\omega_1)(\omega-\omega_2)\cdots(\omega-\omega_n)},
\]
where \(P\) is a polynomial.  Let
\(\omega^{(1)},\omega^{(2)},\ldots\)
% [Source printed page 160]
be the roots of \(P\), and write
\[
  P(\omega)=C(\omega-\omega^{(1)})(\omega-\omega^{(2)})\cdots.
\]
The relation \(R(\infty)=1\) shows that \(P\) too has degree \(n\) and that
\(C=1\); the relation \(R(0)=1\) gives
\[
  \omega^{(1)}\omega^{(2)}\cdots\omega^{(n)}
  =\omega_1\omega_2\cdots\omega_n.
\]
On putting
\[
  \rho^{(k)}=\frac1{2\pi i}\log\omega^{(k)},
\]
we may therefore choose the branches of the logarithms so that
\[
  \rho^{(1)}+\cdots+\rho^{(n)}=\rho_1+\cdots+\rho_n.
\]
It follows that
\[
  \Omega(\rho)=\prod_{k=1}^{n}
  \frac{\sin\pi(\rho-\rho^{(k)})}{\sin\pi(\rho-\rho_k)}.
\]

\numberedparagraph{106}
Consider now the system obtained from \eqref{eq:105:4} by multiplying each
equation by the corresponding function
\[
  \begin{aligned}
    h_0(\rho)&=1,\\
    h_i(\rho)&=\frac1{i^n}
    \exp\left(-\frac{\rho-\rho_1}{i}\right)
    \exp\left(-\frac{\rho-\rho_2}{i}\right)\cdots
    \exp\left(-\frac{\rho-\rho_n}{i}\right),
    \qquad i\ne0.
  \end{aligned}
\]
where the \(\rho_k\) are the roots of \(\varphi_0(\rho)=0\).  The
coefficients of this new system are entire functions of \(\rho\), and the same
is true of its determinant \(\Delta(\rho)\).  This determinant is of normal
type and, together with its minors, converges for every \(\rho\).  Moreover,
\[
  \Delta(\rho)=\Pi(\rho)\Omega(\rho),
  \qquad
  \Pi(\rho)=\prod_{k=1}^{n}\frac{\sin\pi(\rho-\rho_k)}{\pi}.
\]
Indeed, the entries of \(\Delta(\rho)\) are
\[
  \chi_{ik}(\rho)=h_i(\rho)\varphi_0(\rho+i)\psi_{ik}(\rho),
\]
and the infinite product formed from the diagonal entries
\[
  \chi_{ii}(\rho)=
  \prod_{k=1}^{n}
  \left(1+\frac{\rho-\rho_k}{i}\right)
  \exp\left(-\frac{\rho-\rho_k}{i}\right)
\]
converges absolutely to \(\Pi(\rho)\), uniformly in every bounded domain.

% [Source printed page 161]
Thus
\[
  \Delta(\rho)=\prod_{k=1}^{n}
  \frac{\sin\pi(\rho-\rho^{(k)})}{\pi}.
\]
In \(\Delta(\rho)\), replace the entries \(\chi_{ik}\) of the row numbered
\(i\) by the corresponding quantities \(z^k\), and denote the modified
determinant by \(\Delta_i(\rho,z)\).  When \(|z|=1\), these quantities are
bounded as a whole, so the modified determinant converges absolutely.  More is
true: it converges absolutely for every \(z_0\) in the annulus
\(r<|z_0|<R\).

To see this, first replace \(\chi_{ik}(\rho)\) in \(\Delta(\rho)\) by
\(\chi_{ik}(\rho)z_0^{i-k}\), or equivalently replace
\(\varphi_j(\rho)\) by \(\varphi_j(\rho)z_0^j\).  This determinant is again of
normal type.  Indeed, on introducing the new variable \(z'=z/z_0\) in the
differential equation, \(\Delta(\rho)\) is replaced by precisely this modified
determinant, and all the preceding arguments apply.  Now replace all entries
of its row numbered \(i\) by \(z_0^i\); the resulting determinant converges
absolutely.  Absolute convergence is preserved when each entry is multiplied
by the corresponding quantity \(z_0^{k-i}\).  This produces
\(\Delta_i(\rho,z_0)\), proving the assertion.

\numberedparagraph{107}
When \(\Delta(\rho)\ne0\), system \eqref{eq:105:4} has, besides the evident
zero solution, no solution \((\alpha_k)\) for which the \(\alpha_k\) are
bounded as a whole.  Consequently, \eqref{eq:104:2} has no solution of the
form \(z^\rho\sum_k\alpha_kz^k\).

Let \(\rho'\) now be a root of \(\Delta(\rho)\).  Suppose there is an index
\(i\) for which at least one first minor \(M_{ik}(\rho')\) of
\(\Delta(\rho')\) does not vanish.  Then, apart from an arbitrary common
factor, equations \eqref{eq:105:4} have the unique solution
\(\alpha_k=M_{ik}(\rho')\).  From the result just proved for
\(\Delta_i(\rho,z)\), the series
\[
  \sum_{k=-\infty}^{+\infty}M_{ik}(\rho')z^k
  =\Delta_i(\rho',z)
\]
converges
% [Source printed page 162]
whenever \(r<|z|<R\).  Therefore
\[
  z^{\rho'}\Delta_i(\rho',z)
\]
is the required solution \(u(z)\) for \(\rho=\rho'\), and it is unique up to
a constant factor.

We have supposed that for \(\rho=\rho'\) there is a minor
\(M_{ik}(\rho')\ne0\).  This is always so when \(\rho'\) is a simple root of
\(\Delta(\rho)\).  Indeed, the formula
\[
  \frac{d\Delta}{d\rho}
  =\sum_{i,k=-\infty}^{+\infty}
  \frac{d\chi_{ik}}{d\rho}M_{ik}
\]
shows at once that if \(M_{ik}(\rho')=0\) for every \(i,k\), then
\(d\Delta/d\rho=0\); hence \(\rho'\) is not a simple root.

If the quantities \(\rho^{(1)},\ldots,\rho^{(n)}\) are all distinct and no
two differ by an integer, all roots of \(\Delta(\rho)\) are simple.  Otherwise
there are multiple roots, and at such a root all first minors may vanish.  Let
\(\rho'\) be a root of order \(m\).  Expanding
\(\left[d^m\Delta/d\rho^m\right]_{\rho=\rho'}\) in terms of minors of order
\(m\), one sees immediately that at least one such minor is nonzero.  It
follows that system \eqref{eq:105:4}, for
\(\rho=\rho'\), has a certain number \(g\le m\) of independent solutions,
given by minors of \(\Delta(\rho')\); each of these solutions gives rise to a
solution of the differential equation.

Von Koch carried the study of \(\Delta(\rho)\) still further by applying to
infinite determinants the principles of the theory of elementary divisors.
In this way he was able to calculate and study anew the fundamental systems of
\(n\) independent integrals of the differential equation.  I shall not enter
into the details; the interested reader will profit greatly from consulting
the original memoir.\footnote{H.~von Koch, ``Sur les intégrales régulières des
équations différentielles linéaires,'' \emph{Acta Mathematica}, vol.~XVIII
(1894), pp.~337--419.}

\section{Integral Equations}

\numberedparagraph{108}
The reader is doubtless aware of the special importance recently attached to
the theory of
% [Source printed page 163]
integral equations.

The fundamental principles of this young theory have already been set forth in
excellent treatises, and I may confine myself to the questions most closely
connected with our subject.

In his fifth memoir on linear integral equations, published in 1906, Hilbert
reduced the solution of these functional equations to the solution of infinite
systems of linear equations in infinitely many unknowns.  The fourth memoir,
published shortly before and already discussed by us, had as its principal
purpose to prepare this application.  Since then, most of the results forming
the subject of the present volume have been translated into corresponding
results on integral equations.\footnote{Besides Hilbert's fifth memoir, cf.
H.~Weyl, ``Singuläre Integralgleichungen,'' \emph{Mathematische Annalen},
vol.~LXVI (1908), pp.~273--324; M.~Plancherel, ``Note sur les équations
intégrales singulières,'' \emph{Rivista di Fisica, Matematica e Scienze
Naturali} (Pavia, 1909), pp.~37--53; A.~J.~Pell, ``Biorthogonal Systems of
Functions,'' \emph{Transactions of the American Mathematical Society},
vol.~XII (1911), pp.~135--164.}\nobreak\hspace{0.12em}\translatornote{The French original prints
``vol.~XIX''; the article appeared in vol.~XII.}

The relations between the two theories are not limited to a formal analogy;
their objects can be put into an actual correspondence.  The germ of this
correspondence is the application of the method of undetermined coefficients
to series in orthogonal functions.  Consider, for example, the integral
equation of the second kind
\begin{equation}
  \varphi(s)-\int_a^b H(s,t)\varphi(t)\,dt=f(s),
  \tag{7}\label{eq:108:7}
\end{equation}
and suppose it can be satisfied by a function of the form
\[
  \varphi(s)=\sum_{k=1}^{\infty}\varphi_k\alpha_k(s),
\]
where the functions \(\alpha_k(s)\) form a normalized orthogonal system on
\((a,b)\).  Substitute this expansion in the equation, multiply successively
by the \(\alpha_i(s)\), and integrate with respect to \(s\).  The integral
equation becomes the system
\begin{equation}
  \varphi_i-\sum_{k=1}^{\infty}h_{ik}\varphi_k=f_i,
  \qquad i=1,2,\ldots,
  \tag{8}\label{eq:108:8}
\end{equation}

\pagebreak[2]
% [Source printed page 164]
in the unknowns \(\varphi_k\), where
\[
  h_{ik}=\int_a^b\!ds\int_a^b H(s,t)\alpha_i(s)\alpha_k(t)\,dt,
  \qquad
  f_i=\int_a^b f(s)\alpha_i(s)\,ds.
\]
To make this argument rigorous, one must first state precisely the hypotheses
on the functions entering the calculation and then prove that, under those
hypotheses, the passage from \eqref{eq:108:7} to \eqref{eq:108:8} is
permissible.  Suppose further that the hypotheses have been chosen so that
\eqref{eq:108:8} can be treated by one of the methods developed above and has
a solution \((\varphi_k)\).  A new difficulty then arises.  It is not enough
to pass from \eqref{eq:108:7} to \eqref{eq:108:8}; one must also be able to
return.  Once the quantities \(\varphi_k\) have been determined, one must
construct the function \(\varphi(s)\).  In general the series
\(\sum_k\varphi_k\alpha_k(s)\) will not converge; it will be only, so to speak,
a formal expansion of the desired function.

All these questions belong to the theory of integration, and different means
of treating them are known in different cases.  In many special cases the
classical theory of integration, going back to Cauchy, already supplies what
is needed.  But if one always confined oneself to that classical theory,
unexpected complications would arise even in apparently simple cases.  We
prefer from the outset to base the discussion on Lebesgue's notion of the
integral.

\numberedparagraph{109}
Consider real functions \(\alpha_1(s),\alpha_2(s),\ldots\), defined on
\((a,b)\), normalized and pairwise orthogonal.  Let \(f(s)\) be a real
function on the same interval, integrable and with integrable square.  The
\emph{coordinates} of \(f(s)\) relative to the system \([\alpha_k(s)]\) are
the quantities
\[
  f_k=\int_a^b f(s)\alpha_k(s)\,ds.
\]

% [Source printed page 165]
Bessel's familiar identity
\[
  \int_a^b\left[f(s)-\sum_{k=1}^{n}f_k\alpha_k(s)\right]^2ds
  =\int_a^b[f(s)]^2ds-\sum_{k=1}^{n}f_k^2
\]
gives the inequality
\[
  \sum_{k=1}^{n}f_k^2\le\int_a^b[f(s)]^2ds.
\]
Thus the series \(\sum_kf_k^2\) converges.  One may ask whether the converse
is also true: given a sequence of numbers \(f_k\) such that
\(\sum_kf_k^2\) converges, does there exist a function \(f(s)\) having these
numbers as its coordinates?

For the answer to be affirmative, it is indispensable to understand the word
\emph{integral} in Lebesgue's sense.  With this understanding, I proved that
there do indeed exist summable functions \(f(s)\) (Lebesgue-integrable
functions) with summable squares and prescribed coordinates
\(f_k\).\footnote{F.~Riesz, ``Sur les systèmes orthogonaux de fonctions,''
\emph{Comptes rendus}, March~18, 1907.  The same fact was discovered
independently by E.~Fischer, ``Sur la convergence en moyenne,'' \emph{Comptes
rendus}, May~13, 1907.  For a detailed proof and further bibliographical
references, see my memoir ``Untersuchungen über Systeme integrierbarer
Funktionen,'' \emph{Mathematische Annalen}, vol.~LXIX (1910), pp.~449--497,
or the recent exposition by W.~H.~Young and G.~C.~Young, ``On the Theorem of
Riesz--Fischer,'' \emph{Quarterly Journal of Mathematics}, vol.~XLIV (1912),
pp.~49--88.}
Such a function is obtained by considering the uniformly convergent
series\translatornote{The printed formula uses \(s\) both as the upper limit
and as the variable of integration; \(t\) has been substituted as the dummy
variable.}
\[
  F(s)=\sum_{k=1}^{\infty}f_k\int_a^s\alpha_k(t)\,dt.
\]
The function \(F(s)\) has a derivative except perhaps at values of \(s\)
forming a set of measure zero.  Define \(f(s)\) to equal this derivative
wherever it exists and give it arbitrary values at all remaining points.  Then
\(f(s)\) is summable, has summable square, and has the \(f_k\) as coordinates.
Moreover,
\begin{equation}
  \int_a^b[f(s)]^2ds=\sum_{k=1}^{\infty}f_k^2.
  \tag{9}\label{eq:109:9}
\end{equation}

% [Source printed page 166]
Recall also that if the system \([\alpha_k(s)]\) is complete, formula
\eqref{eq:109:9} holds for every real summable function \(f(s)\) with
summable square and coordinates \(f_k\); see \hyperref[sec:102]{no.~102}.  In
particular, if \(f_1(s)\) and \(f_2(s)\) have the same coordinates, then
\(f_1(s)-f_2(s)\) has every coordinate zero, and the integral of
\([f_1(s)-f_2(s)]^2\) vanishes.  Thus \(f_1(s)=f_2(s)\), except perhaps on a
set of measure zero.  In view of the role of such sets in Lebesgue's theory,
one may simply say that a function \(f(s)\) is uniquely determined by its
coordinates.

Finally recall the formula
\begin{equation}
  \int_a^b f(s)g(s)\,ds=\sum_{k=1}^{\infty}f_kg_k,
  \tag{10}\label{eq:109:10}
\end{equation}
obtained by applying \eqref{eq:109:9} to
\(\tfrac12[f(s)+g(s)]\) and \(\tfrac12[f(s)-g(s)]\), and subtracting the
second expansion from the first.

These results extend immediately when \(f(s),g(s)\) and the coordinates
\(f_k,g_k\) are no longer assumed real; one merely replaces squares by
squares of absolute values.  It is not even essential that the functions
\(\alpha_k(s)\) be real; the definitions and statements need only slight
modification.

All the results likewise extend to functions of several variables.  In
particular, if \([\alpha_k(s)]\) is a normalized, complete orthogonal system of
real functions on \((a,b)\), then \([\alpha_i(s)\alpha_k(t)]\) is such a
system on the square \(a\le s\le b\), \(a\le t\le b\), and
\begin{align}
  &\int_a^b\!\int_a^b f(s,t)g(s,t)\,ds\,dt \notag\\
  &\quad=\sum_{i,k=1}^{\infty}
  \left(\int_a^b\!\int_a^b
  f(s,t)\alpha_i(s)\alpha_k(t)\,ds\,dt\right)
  \left(\int_a^b\!\int_a^b
  g(s,t)\alpha_i(s)\alpha_k(t)\,ds\,dt\right),
  \tag{11}\label{eq:109:11}
\end{align}
where \(f(s,t)\) and \(g(s,t)\), as well as their squares, are assumed
summable.

\numberedparagraph{110}
Consider the integral equation \eqref{eq:108:7}.  To begin with a simple
% [Source printed page 167]
case, suppose that \(H(s,t)\) and \(f(s)\) are continuous.  Using a normalized,
complete orthogonal system of real functions \(\alpha_k(s)\), the passage from
\eqref{eq:108:7} to \eqref{eq:108:8} is made by replacing \(g(s)\) in
\eqref{eq:109:10} by \(\alpha_i(s)\), and by replacing \(f(s,t)\) and
\(g(s,t)\) in \eqref{eq:109:11} by \(H(s,t)\) and
\(\alpha_i(s)\alpha_k(t)\), respectively.  Since
\begin{equation}
  \sum_{i,k=1}^{\infty}|h_{ik}|^2
  =\int_a^b\!\int_a^b|H(s,t)|^2\,ds\,dt
  \tag{12}\label{eq:110:12}
\end{equation}
converges, the substitution defined by \((h_{ik})\) is completely continuous.
All the results concerning this special class of linear substitutions may
therefore be applied both to \eqref{eq:108:8} and to its associated
homogeneous system.

At the same time one may consider the transposed equation
\begin{equation}
  \psi(t)-\int_a^b H(s,t)\psi(s)\,ds=g(t),
  \tag{13}\label{eq:110:13}
\end{equation}
which, by means of the orthogonal system \([\alpha_k(s)]\), becomes
\begin{equation}
  \psi_k-\sum_{i=1}^{\infty}h_{ik}\psi_i=g_k,
  \qquad k=1,2,\ldots,
  \tag{14}\label{eq:110:14}
\end{equation}
in the unknowns \(\psi_k\).  Every solution \((\varphi_k)\) of
\eqref{eq:108:8} for which \(\sum_k|\varphi_k|^2\) converges corresponds to a
summable function \(\varphi(s)\) with summable \(|\varphi(s)|^2\); the same is
true of every admissible solution \((\psi_k)\) of \eqref{eq:110:14}.\translatornote{The French original refers here to equation~(15); equation~(14) is meant.}  The
function
\[
  \varphi(s)-\int_a^bH(s,t)\varphi(t)\,dt,
\]
whose coordinates are
\(\varphi_i-\sum_{k=1}^{\infty}h_{ik}\varphi_k=f_i\), equals \(f(s)\), except
perhaps on a set of measure zero.  But the coordinates determine
\(\varphi(s)\) only up to the addition of a function that vanishes almost
everywhere.\translatornote{The French original says ``a function whose integral
is zero.''  Completeness of the orthogonal system instead implies that the
difference has squared modulus with integral zero, hence vanishes almost
everywhere.}
This indeterminacy can be used to obtain an exact solution of
\eqref{eq:108:7}: the equation itself makes the necessary correction.  Indeed,
\[
  \varphi^*(s)=f(s)+\int_a^bH(s,t)\varphi(t)\,dt
\]
is continuous, has the same coordinates as \(\varphi(s)\), and hence satisfies
\eqref{eq:108:7} exactly.

% [Source printed page 168]
The same argument applies to \eqref{eq:110:13}.

These brief indications should suffice to show how every result about
completely continuous substitutions and their corresponding systems of
equations can be translated into a result about integral equations.  It is
also easy to see that continuity of \(H(s,t)\) is not essential.  In fact,
most of the preceding argument remains unchanged if one assumes only that
\[
  H(s,t),\quad |H(s,t)|^2,\quad f(s),\quad |f(s)|^2,
  \quad g(s),\quad |g(s)|^2
\]
are summable.  Even if, for example, one wants to ensure that the solution
\(\varphi^*(s)\) is continuous, it is enough, instead of assuming the kernel
\(H(s,t)\) continuous, to choose it so that the first of the
integrals\translatornote{In the first printed integral, \(t\) appears both as
the upper limit and as the variable of integration; \(\tau\) has been
substituted as the dummy variable.}
\[
  \int_a^tH(s,\tau)\,d\tau,
  \qquad
  \int_a^b|H(s,t)|^2\,dt
\]
represents a continuous function of \(s\) for every \(t\), and the second
represents a bounded function.  The kernels \(|s-t|^{-\alpha}\), with
\(\alpha<\tfrac12\), are examples.

Finally, convergence of \eqref{eq:110:12} also permits the method of infinite
determinants to be applied to systems \eqref{eq:108:8} and
\eqref{eq:110:14}.  By an easy calculation, one can even recover in this way
all Fredholm's formulas.\footnote{Cf. E.~Marty, ``Transformation d'un
déterminant infini en un déterminant de Fredholm,'' \emph{Bulletin des Sciences
mathématiques}, vol.~XXXIII (1909), pp.~296--300; J.~Mollerup, ``Sur l'identité
du déterminant de Fredholm et d'un déterminant infini de v.~Koch,'' ibid.,
vol.~XXXVI (1912), pp.~130--136.  Marty assumes that
\(\sum_{i,k}|h_{ik}|\) converges; Mollerup adopts the assumption used in the
text.}

\numberedparagraph{111}
To approach equations whose kernels have higher singularities, we introduce
the notion of a linear functional transformation.  To every function \(f(s)\)
of the kind considered---summable together with \(|f|^2\)---associate a
function \(f'(s)\) of the same kind, in such a way that
\[
  f_1'+f_2'=(f_1+f_2)',
  \qquad cf'=(cf)',
\]
and suppose in addition that there is a constant \(M^2\) such that, for
% [Source printed page 169]
every function \(f(s)\),
\begin{equation}
  \int_a^b|f'(s)|^2ds\le M^2\int_a^b|f(s)|^2ds.
  \tag{15}\label{eq:111:15}
\end{equation}
\begin{definition}
Every correspondence of this kind is called a \emph{linear transformation}.
\end{definition}
It is evident that, by associating the coordinates of \(f'(s)\) with those of
\(f(s)\), one defines a linear substitution of Hilbert space.  Conversely, if
to each function \(f(s)\) of the kind considered we associate the function
\(f'(s)\) whose coordinates are obtained by applying a linear substitution
\(A\) to those of \(f(s)\), then a linear functional transformation has been
defined.

I do not propose to pursue every detail of this relation between
transformations and linear substitutions.  Plainly the whole theory of linear
substitutions developed in \hyperref[chap:4]{Chapter~IV} can be translated into a theory of
linear functional transformations.

In particular, under the conditions previously imposed on \(H(s,t)\), the
integral
\begin{equation}
  f'(s)=\int_a^bH(s,t)f(t)\,dt
  \tag{16}\label{eq:111:16}
\end{equation}
represents a linear transformation.  The conditions imposed on \(H(s,t)\) are
not essential; in any event, it is enough that \eqref{eq:111:16} exist and
that \eqref{eq:111:15} hold.  I hasten to observe, however, that
\eqref{eq:111:16} does not exhaust the possibilities.  The identity
transformation \(f'(s)=f(s)\), and the transformation
\(f'(s)=g(s)f(s)\), where \(g(s)\) is summable and bounded, cannot be put in
this form.\footnote{A detailed study of the analytic representation of linear
functional transformations was given by M.~Plancherel in ``Contribution à
l'étude de la représentation d'une fonction arbitraire par des intégrales
définies,'' \emph{Rendiconti del Circolo Matematico di Palermo}, vol.~XXX
(1910), pp.~289--335.}  This observation gives rise to the question discussed
in the next paragraph.

\numberedparagraph{112}
The study of a linear substitution \(A\) brings in
% [Source printed page 170]
other related substitutions: recall its transpose \(\mathfrak A\), its
iterates \(A^k\), the inverses \(A^{-1}\) and \((E-A)^{-1}\), the substitutions
\(A^{(k)}\) studied in \hyperref[sec:82]{no.~82}, and, when \(A\) is real and
symmetric, all substitutions \(f(A)\).  Suppose now that a linear
transformation of the form \eqref{eq:111:16} is under consideration.  By means
of an orthogonal system \([\alpha_k(s)]\), it corresponds to a linear
substitution \(A\).  This substitution gives rise to those just recalled, and
each of them, through the same orthogonal system, corresponds to a linear
functional transformation.  For each such transformation one may ask whether
it can be put in either of the forms
\[
  \int_a^bG(s,t)f(t)\,dt,
  \qquad
  cf(s)+\int_a^bG(s,t)f(t)\,dt.
\]

For the transformations corresponding to \(\mathfrak A\) and the \(A^k\), the
question reduces at once to interchanging the order of certain successive
integrations, and I shall not dwell on it.  For \(A^{-1}\), one is led to the
delicate problem of integral equations of the first kind, which cannot be
summarized in a few lines.  For \((E-A)^{-1}\), the question involves one of
the fundamental problems in the theory of integral equations of the second
kind.  Since Fredholm's investigations it has been known that, in the simplest
cases, the solution of \eqref{eq:108:7} can be put in the form
\[
  \varphi(s)=f(s)+\int_a^bG(s,t)f(t)\,dt,
\]
where \(G(s,t)\), called the \emph{resolvent kernel}, is independent of the
given function \(f(s)\).

This fact may be attached to the relation
\[
  (E-A)^{-1}=E+(E-A)^{-1}A.
\]
It shows that, to calculate \(G(s,t)\), one need only apply to \(H(s,t)\),
regarded as a function of \(s\), the linear transformation corresponding to
\((E-A)^{-1}\).  The legitimacy of this procedure is evident when, for
example,
% [Source printed page 171]
the kernel is continuous; it also applies when the kernel has high
singularities.  Everything plainly reduces to proving that the order of
certain passages to the limit may be interchanged.

The same reasoning applies to all substitutions \(f(A)\) which can be put in
the form \(cE+f_1(A)A\).  I cannot enter into details without going deeply into
the theory of integration; readers familiar with that theory will readily
supply the necessary arguments.

Finally, it should be observed that not only the results, but almost all the
arguments concerning systems of equations in infinitely many unknowns have
their analogues in the theory of integral equations.  This analogy may be used
when the consideration of summable functions with summable squares is not
suited to the data of the problem.

\section{Trigonometric Series}

\numberedparagraph{113}
Given a sequence extending indefinitely in both directions \((\alpha_k)\), attach
to it a matrix \((a_{ik})\) by putting \(a_{ik}=\alpha_{k-i}\).  This matrix belongs
to a special type which we shall call the \emph{Laurent type}, because of its
relation to the series
\[
  \sum_{k=-\infty}^{+\infty}\alpha_kz^k.
\]
The theory of Laurent substitutions and forms corresponding to these special
matrices is due to Toeplitz.\footnote{O.~Toeplitz, ``Zur Transformation der
Scharen bilinearer Formen von unendlich vielen Veränderlichen,''
\emph{Nachrichten der Königlichen Gesellschaft der Wissenschaften zu
Göttingen} (1907), pp.~110--116; ``Zur Theorie der quadratischen Formen von
unendlich vielen Veränderlichen,'' ibid. (1910), pp.~489--506; ``Zur Theorie
der quadratischen und bilinearen Formen von unendlich vielen
Veränderlichen,'' \emph{Mathematische Annalen}, vol.~LXX (1911),
pp.~351--376.}

The problem which leads most immediately to Laurent matrices is the following.
If one seeks the Laurent expansion of the quotient of two functions given by
their Laurent expansions, the method of undetermined coefficients
% [Source printed page 172]
plainly leads to a system of linear equations in infinitely many unknowns;
the coefficients of these equations form the Laurent matrix corresponding to
the series in the denominator.  We encountered a special problem of this kind
in \hyperref[sec:13]{no.~13}, where we discussed a note of Appell.

We prefer to adopt at once a more general standpoint, connecting Laurent
substitutions and forms with the theory of systems of orthogonal functions.
Recall that we already alluded to this in \hyperref[sec:103]{no.~103}, where
the coefficients of certain quadratic forms were Fourier coefficients.

\numberedparagraph{114}
Consider the functions
\[
  \varphi_k(s)=e^{2\pi iks}=\cos2\pi ks+i\sin2\pi ks,
\]
where \(k\) ranges indefinitely in both directions.  Plainly,
\[
  \int_{-1/2}^{1/2}\varphi_k(s)\overline{\varphi_l(s)}\,ds
  =\int_{-1/2}^{1/2}e^{2\pi i(k-l)s}\,ds
  =\begin{cases}1,&k=l,\\0,&k\ne l.\end{cases}
\]
The system is complete on \((-\tfrac12,\tfrac12)\), as follows immediately
from the analogous property of
\((\cos2\pi ks,\sin2\pi ks)\), \(k=0,1,2,\ldots\).  Thus the
\(\varphi_k\) form a normalized, complete orthogonal system on this
interval.\footnote{Instead of \([\varphi_k(s)]\), we could have used the more
familiar system \((e^{iks})\) on \((-\pi,\pi)\) or \((0,2\pi)\); that choice
would, however, have burdened the formulas with inessential numerical
constants.}
The coordinates of a function \(f(s)\) relative to this system are
\[
  f_k=\int_{-1/2}^{1/2}f(s)\overline{\varphi_k(s)}\,ds.
\]
% [Source printed page 173]
All the results concerning orthogonal systems of real functions in
\hyperref[sec:102]{nos.~102} and \hyperref[sec:109]{109} apply, with evident
modifications, to \(\varphi_k(s)\).  In particular, if \(f_k\) and \(g_k\)
are the coordinates of functions \(f(s)\) and \(g(s)\), summable together
with \(|f(s)|^2\) and \(|g(s)|^2\), then
\(\int f(s)g(s)\,ds=\sum_k f_{-k}g_k\).  Moreover, the coordinate numbered
\(i\) of \(f(-s)\varphi_k(s)\) is plainly \(f_{k-i}\), and hence
\begin{equation}
  \int_{-1/2}^{1/2}f(-s)g(s)\overline{\varphi_i(s)}\,ds
  =\sum_{k=-\infty}^{+\infty}f_{k-i}g_k.
  \tag{17}\label{eq:114:17}
\end{equation}

Suppose now that the matrix \((a_{ik}=\alpha_{k-i})\) defines a linear
(distributive and bounded) substitution of Hilbert space,
\begin{equation}
  x'_i=\sum_{k=-\infty}^{+\infty}\alpha_{k-i}x_k,
  \qquad i=-\infty,\ldots,+\infty.
  \tag{18}\label{eq:114:18}
\end{equation}
The series \(\sum_k|\alpha_k|^2\) then converges, and the \(\alpha_k\) are the
coordinates of a function \(\alpha(s)\), summable together with
\(|\alpha(s)|^2\).  Let \(x(s)\) be the function of the same kind whose
coordinates are the \(x_k\).  Consider
\[
  x'(s)=\alpha(-s)x(s).
\]
By \eqref{eq:114:17}, the coordinates \(x'_k\) of this function are given by
\eqref{eq:114:18}.  Thus the orthogonal system \([\varphi_k(s)]\) associates
with the substitution \(A=(\alpha_{k-i})\) the functional transformation
\(x'(s)=\alpha(-s)x(s)\).  Since \(\sum_k|x_k|^2\) and
\(\sum_k|x'_k|^2\) represent the integrals of \(|x(s)|^2\) and
\(|x'(s)|^2\), respectively,
\begin{equation}
  \int_{-1/2}^{1/2}|\alpha(-s)x(s)|^2ds
  \le M_A^2\int_{-1/2}^{1/2}|x(s)|^2ds.
  \tag{19}\label{eq:114:19}
\end{equation}
Let \(\mathfrak E\) be the set of values \(s\) for which
\(|\alpha(-s)|>M_A\), and put \(x(s)=1\) on this set and zero elsewhere.
Inequality \eqref{eq:114:19} gives
\[
  \int_{\mathfrak E}|\alpha(-s)|^2ds\le M_A^2\int_{\mathfrak E}ds.
\]
Since \(|\alpha(-s)|>M_A\) on \(\mathfrak E\), the set \(\mathfrak E\) must
have measure zero.  The function \(\alpha(s)\) may be modified arbitrarily
there
% [Source printed page 174]
without changing its coordinates; consequently, we may suppose that
\(|\alpha(s)|\le M_A\) everywhere.

An analogous argument shows that if the Hermitian form
\begin{equation}
  \sum_{i=-\infty}^{+\infty}
  \left|\sum_{k=-\infty}^{+\infty}\alpha_{k-i}x_k\right|^2
  =\int_{-1/2}^{1/2}|\alpha(-s)x(s)|^2ds
  \tag{20}\label{eq:114:20}
\end{equation}
remains at least \(J^2\) whenever \(\sum_k|x_k|^2=1\), then the set on which
\(|\alpha(-s)|<J\) has measure zero.  Thus we may also suppose that
\(|\alpha(s)|\ge J\) everywhere.

Conversely, given a bounded summable function \(\alpha(s)\), its coordinates
\(\alpha_k\) define a linear substitution \(A=(\alpha_{k-i})\).  Indeed, if
\(|\alpha(s)|\le M\), and if \(x(s)\) and \(|x(s)|^2\) are summable, then so
are \(x'(s)=\alpha(-s)x(s)\) and \(|x'(s)|^2\), and
\[
  \int_{-1/2}^{1/2}|x'(s)|^2ds
  \le M^2\int_{-1/2}^{1/2}|x(s)|^2ds.
\]
Hence the substitution corresponding to \(x'(s)=\alpha(-s)x(s)\) is bounded,
and \(M_A\le M\).

Suppose now that the Laurent substitution \(A\) has an inverse \(A^{-1}\).
Since the value of the Hermitian form \eqref{eq:114:20} must then be at least
\(1/M_{A^{-1}}^2\) for every system \((x_k)\) with
\(\sum_k|x_k|^2=1\), we may suppose
\[
  |\alpha(s)|\ge\frac1{M_{A^{-1}}}.
\]
Thus \(1/\alpha(s)\) is bounded and summable.  Denote its coordinates by
\(\alpha_k^{(-1)}\).  The matrices
\((\alpha_{k-i})\) and \((\alpha_{k-i}^{(-1)})\) plainly satisfy the general
relations between the coefficients of inverse substitutions.  Hence
\emph{the inverse \(A^{-1}\) is the Laurent substitution corresponding to
the function \(1/\alpha(s)\).}

More generally, \emph{the inverse of \(\lambda E-A\), when it exists, is
given by the coordinates of the function \(1/[\lambda-\alpha(s)]\); and,
conversely, if this function is
% [Source printed page 175]
bounded, or can be made bounded after neglecting values of \(s\) forming a
set of measure zero, the inverse \((\lambda E-A)^{-1}\) exists.}

In particular, \emph{if \(\alpha(s)\) is continuous, the set of values
\(\lambda\) for which the inverse of \(\lambda E-A\) does not exist coincides
with the set of values taken by \(\alpha(s)\).}  If the doubly infinite series
\(\sum_k\alpha_kz^k\) converges in a circular annulus containing \(z=1\), the
singular values of \(\lambda\) are the values taken by the series on
\(|z|=1\); this set is therefore an analytic curve.\footnote{Toeplitz gave a
detailed and very interesting study of the special matrices corresponding to
actual Laurent series in the last of the works cited above.}

\numberedparagraph{115}
Suppose now that \(\alpha(s)\) is real, which is plainly equivalent to
\(\alpha_{-k}=\overline{\alpha_k}\) for every \(k\).  Consider the Hermitian
form
\begin{equation}
  A(x,x)=\sum_{i,k=-\infty}^{+\infty}
  \alpha_{k-i}x_i\overline{x_k}
  =\int_{-1/2}^{1/2}\alpha(s)|x(s)|^2ds.
  \tag{21}\label{eq:115:21}
\end{equation}
Let \(M\) be the upper bound of \(\alpha(s)\) in Lebesgue's sense; that is,
let \(M\) be the least number for which the set of \(s\) satisfying
\(\alpha(s)>M\) has measure zero.  Define the lower bound \(m\) analogously.
When the Hermitian form \eqref{eq:115:21} is varied subject to
\begin{equation}
  E(x,x)=\sum_{k=-\infty}^{+\infty}|x_k|^2
  =\int_{-1/2}^{1/2}|x(s)|^2ds=1,
  \tag{22}\label{eq:115:22}
\end{equation}
its values plainly remain between \(m\) and \(M\).  Conversely, let
\(\varepsilon>0\) be arbitrarily small.  The set on which
\(\alpha(s)>M-\varepsilon\) has a nonzero measure \(\mu\); put
\(x(s)=1/\sqrt\mu\) on this set and zero elsewhere.  The coordinates of this
function satisfy
% [Source printed page 176]
\eqref{eq:115:22}, and \eqref{eq:115:21} gives
\[
  A(x,x)\ge M-\varepsilon.
\]
An analogous argument applies to the lower bound \(m\).  Hence \emph{the
upper and lower bounds of \(\alpha(s)\), in Lebesgue's sense, are also the
exact bounds of the Hermitian form \(A(x,x)\).}

In particular, \emph{a necessary and sufficient condition that
\(\alpha(s)\) be nonnegative, or can be made nonnegative after neglecting a
set of measure zero, is that the Hermitian form \(A(x,x)\) be nonnegative.}

\numberedparagraph{116}
The preceding results apply immediately to the problem of determining the two
bounds of a real function \(\alpha(2\pi s)\) given by its Fourier expansion
\begin{equation}
  \frac12a_0+\sum_{k=1}^{\infty}
  \bigl(a_k\cos2\pi ks+b_k\sin2\pi ks\bigr),
  \tag{23}\label{eq:116:23}
\end{equation}
whose coefficients are given by the familiar
formulas\translatornote{In the printed display, the sine term in
\eqref{eq:116:23} has \(b\) instead of \(b_k\), and the coefficient formulas
use \(\alpha(2\pi s)\) while integrating from \(-\pi\) to \(\pi\).  The
evident typographical errors have been corrected by writing \(b_k\) and using
the dummy variable \(t\) in the standard coefficient formulas.  The printed
identifications of \(\alpha_k\) and \(\alpha_{-k}\) also reverse the signs of
\(ib_k\); since \(\varphi_k(s)=e^{2\pi iks}\), those signs, and the
corresponding signs in the determinant below, are corrected here.}
\[
  a_k=\frac1\pi\int_{-\pi}^{\pi}\alpha(t)\cos kt\,dt,
  \qquad
  b_k=\frac1\pi\int_{-\pi}^{\pi}\alpha(t)\sin kt\,dt.
\]
Without concerning ourselves with whether the Fourier series converges, we
agree to regard it as a symbol for the function which gives rise to it.  The
coefficients \(a_k,b_k\) play the role of coordinates.  The coordinates
relative to \([\varphi_k(s)]\) are obtained by putting
\[
  \alpha_0=\frac12a_0,
  \qquad
  \alpha_k=\frac12(a_k-ib_k),
  \qquad
  \alpha_{-k}=\overline{\alpha_k}=\frac12(a_k+ib_k).
\]
To determine the bounds of \(\alpha(2\pi s)\), we therefore consider the
Hermitian form \eqref{eq:115:21} formed with these quantities \(\alpha_k\).

Let \(m_n\) and \(M_n\) denote the bounds of the section
\begin{equation}
  [A(x,x)]_n=\sum_{i,k=0}^{n}\alpha_{k-i}x_i\overline{x_k},
  \tag{24}\label{eq:116:24}
\end{equation}
% [Source printed page 177]
where the \(x_k\) vary subject to
\[
  \sum_{k=0}^{n}|x_k|^2=1.
\]
More generally, consider the sections
\[
  [A(x,x)]_{l,n}=\sum_{i,k=l}^{n}\alpha_{k-i}x_i\overline{x_k},
  \qquad l<n.
\]
Since the coefficient matrix depends only on \(n-l\), their bounds are
\(m_{n-l}\) and \(M_{n-l}\).  Plainly,
\[
  m_n\ge m,\qquad m_n\ge m_{n+1},\qquad
  M_n\le M,\qquad M_n\le M_{n+1}.
\]

Suppose that for some system \((x_k)\), the Hermitian form
\eqref{eq:115:21} comes within \(\varepsilon\) of its upper bound \(M\), and
put
\[
  x_k^{(n)}=
  \frac{x_k}{\left(\sum_{i=-n}^{n}|x_i|^2\right)^{1/2}}
  \quad(k=-n,-n+1,\ldots,n),
\]
with \(x_k^{(n)}=0\) for all other indices.  The relation
\translatornote{The French original uses the subscript \(n\) here and in the
next inequality, although \(x^{(n)}\) is supported on the indices
\(-n,\ldots,n\).  The section is therefore \([A]_{-n,n}\), whose upper bound
is \(M_{2n}\).}
\[
  \lim_{n\to\infty}[A(x^{(n)},x^{(n)})]_{-n,n}=A(x,x)
\]
shows that, for sufficiently large \(n\),
\([A(x^{(n)},x^{(n)})]_{-n,n}>M-2\varepsilon\).  Hence
\(M_{2n}>M-2\varepsilon\); since \(M_n\) is nondecreasing and never exceeds
\(M\), it follows that \(M_n\to M\).  The analogous argument gives
\(m_n\to m\), or equivalently
\[
  m_n\longrightarrow m,
  \qquad M_n\longrightarrow M.
\]

By a familiar theorem, the bounds \(m_n\) and \(M_n\) of the Hermitian form
\eqref{eq:116:24} are respectively the smallest and largest roots
\translatornote{The printed French reverses ``largest'' and
``smallest'' at this point; the order has been corrected to agree with the
definitions of \(m_n\) and \(M_n\).}
of the determinantal equation
\[
  \det\!\begin{pmatrix}
    \alpha_0-x&\alpha_1&\alpha_2&\cdots&\alpha_n\\
    \alpha_{-1}&\alpha_0-x&\alpha_1&\cdots&\alpha_{n-1}\\
    \vdots&\vdots&\vdots&&\vdots\\
    \alpha_{-n}&\alpha_{-n+1}&\alpha_{-n+2}&\cdots&\alpha_0-x
  \end{pmatrix}=0.
\]
In summary, \emph{the upper and lower bounds, in Lebesgue's sense,
% [Source printed page 178]
of the function \(\alpha(2\pi s)\) having expansion \eqref{eq:116:23} are,
as \(n\to\infty\), the limits of the largest and smallest roots of}
\[
  \det\!\begin{pmatrix}
    a_0-2x&a_1-ib_1&\cdots&a_n-ib_n\\
    a_1+ib_1&a_0-2x&\cdots&a_{n-1}-ib_{n-1}\\
    \vdots&\vdots&&\vdots\\
    a_n+ib_n&a_{n-1}+ib_{n-1}&\cdots&a_0-2x
  \end{pmatrix}=0.
\]

\numberedparagraph{117}
Let us finally observe that these beautiful results of Toeplitz form part of a
group of theorems, due to several authors, more or less closely connected with
Laurent matrices and all belonging to the line of thought of the important and
well-known work of Landau and Carathéodory on Picard's celebrated
theorem.\footnote{Cf. the communications of Toeplitz, Carathéodory, Fejér, and
Fischer, \emph{Rendiconti del Circolo Matematico di Palermo}, vol.~XXXII
(1911), pp.~191--256; G.~Herglotz, \emph{Berichte der Königlich Sächsischen
Gesellschaft der Wissenschaften zu Leipzig}, vol.~LXIII (1911), pp.~501--511;
and I.~Schur and G.~Frobenius, \emph{Sitzungsberichte der Königlich Preussischen
Akademie der Wissenschaften} (1912), pp.~4--31.}

We shall indicate the principal facts briefly.  By Picard's theorem, a
nonconstant entire function takes every prescribed finite value, with at most
one exceptional value.  This theorem, discovered in 1879, was deepened in an
unexpected way by Landau in 1904.  Landau proved that, to determine a circle
\(|z|<R\) in which a convergent power series
\[
  a_0+a_1z+\cdots,\qquad a_1\ne0,
\]
takes at least once one of the values \(0\) or \(1\), it is enough to know
only the first two coefficients \(a_0,a_1\).  In the following year,
Carathéodory gave the exact value of \(R\) as a function of \(a_0,a_1\).  He
was led to it by connecting this problem and a large class of analogous
problems with the following one.  The first \(n+1\) terms of the power series
\begin{equation}
  \frac12+c_1z+c_2z^2+\cdots+c_nz^n+\cdots
  \tag{25}\label{eq:117:25}
\end{equation}
are given;
% [Source printed page 179]
can the remaining terms be chosen so that the series converges throughout the
disk \(|z|<1\) and the real part of the function it represents is everywhere
positive?

Carathéodory's answer is the following.  \emph{The series
\eqref{eq:117:25} can be completed in the required manner if and only if the
point}
\[
  y'_1=c'_1,\quad y''_1=-c''_1,\quad\ldots,\quad
  y'_n=c'_n,\quad y''_n=-c''_n,
  \qquad(c_k=c'_k+ic''_k),
\]
\emph{of real \(2n\)-dimensional space belongs to the smallest convex domain
containing the closed curve defined parametrically by}
\[
  y'_1=\cos\theta,\quad y''_1=\sin\theta,\quad\ldots,\quad
  y'_n=\cos n\theta,\quad y''_n=\sin n\theta.
\]

More recently, the investigations of Toeplitz just discussed supplied a purely
algebraic answer to Carathéodory's question.  Put
\(c_0=1\), \(c_{-k}=\overline{c_k}\).  Then \emph{the series
\eqref{eq:117:25} can be continued in the required manner if and only if the
Hermitian form}
\begin{equation}
  \sum_{i,k=0}^{n}c_{k-i}x_i\overline{x_k}
  \tag{26}\label{eq:117:26}
\end{equation}
\emph{is nonnegative.}

The equivalence of the two conditions was also established by Carathéodory,
E.~Fischer, I.~Schur, and Frobenius in the works cited above.

Suppose now that all coefficients of \eqref{eq:117:25} are given.  One may
then ask whether the given series represents in \(|z|<1\) a function whose
real part remains positive there.  This question was posed and solved
independently by Carathéodory and by the author.\footnote{C.~Carathéodory,
``Über den Variabilitätsbereich der Fourierschen Konstanten von positiven
harmonischen Funktionen,'' \emph{Rendiconti del Circolo Matematico di Palermo},
vol.~XXXII (1911), pp.~193--217; F.~Riesz, ``Sur certains systèmes singuliers
d'équations intégrales,'' \emph{Annales de l'École Normale}, 3rd series,
vol.~XXVIII (1911), pp.~33--62.}  The answer is as follows: \emph{the series
% [Source printed page 180]
\eqref{eq:117:25} converges in \(|z|<1\), and the real part of the function it
represents is everywhere positive, if and only if either of the two equivalent
conditions just stated is satisfied for every integer \(n\).}

Indeed, put \(z=re^{2\pi is}\).  The real part of \eqref{eq:117:25} becomes
\begin{equation}
  \frac12+\sum_{k=1}^{\infty}
  \bigl(c'_kr^k\cos2k\pi s-c''_kr^k\sin2k\pi s\bigr),
  \tag{27}\label{eq:117:27}
\end{equation}
and we require that, for \(r<1\), this series converge and represent a
positive function.  By \hyperref[sec:115]{no.~115}, this requires the
Hermitian form
\[
  \sum_{i,k=-\infty}^{+\infty}
  c_{k-i}r^{|k-i|}x_i\overline{x_k}
\]
to exist and be nonnegative for every \(r<1\).  In particular, its sections
must be nonnegative for \(r<1\); on passing to the limit \(r=1\), the
Hermitian forms \eqref{eq:117:26} must therefore be nonnegative.

Conversely, suppose that \eqref{eq:117:26} is nonnegative for every integer
\(n\).  Put \(x_0=1\), \(x_1=\cdots=x_{n-1}=0\), and \(x_n=-c_n\).
\translatornote{The printed text has \(x_1=1\) here; since the form
\eqref{eq:117:26} is indexed from \(0\) through \(n\), the required choice is
\(x_0=1\).} Then
\eqref{eq:117:26} becomes \(1-|c_n|^2\); hence \(|c_n|\le1\) for every
\(n\).  It follows that \eqref{eq:117:25}, \eqref{eq:117:27}, and the series
\begin{equation}
  \sum_{k=-\infty}^{+\infty}|c_kr^{|k|}|^2
  \le\sum_{k=-\infty}^{+\infty}r^{2|k|}
  =\frac{1+r^2}{1-r^2}
  \tag{28}\label{eq:117:28}
\end{equation}
converge for every \(r<1\).  Since \eqref{eq:117:28} converges, the results of
\hyperref[sec:115]{no.~115} apply; the harmonic function defined by
\eqref{eq:117:27} is therefore nonnegative for every \(r<1\).  Finally, a
harmonic function cannot attain its minimum in the interior of a domain in
which it exists; hence the function \eqref{eq:117:27} is positive for every
\(r<1\).

Many other questions of this kind are treated in the works cited at the
beginning of this paragraph, especially in the memoir of Carathéodory and
Fejér.

% End inlined .parts/pages_121_180.tex

\clearpage
\pdfbookmark[0]{Table of Contents}{toc}
\tableofcontents

\end{document}
